\documentclass[10pt]{article}
\usepackage[letterpaper,margin=1.0in,columnsep=0.28in]{geometry}
\usepackage{mathpazo}
\usepackage{fontspec}
\setmainfont{TeX Gyre Pagella}
\setmonofont[Scale=0.82]{DejaVu Sans Mono}
\newfontfamily\fbfont{FreeSerif}
\usepackage{newunicodechar}
\newfontfamily\amirifont{Amiri}[Script=Arabic]
\TeXXeTstate=1
\newunicodechar{ﷻ}{{\amirifont ﷻ}}
\newunicodechar{⟺}{{\fbfont ⟺}}
\newunicodechar{⟀}{{\fbfont ⟀}}
\newunicodechar{⬖}{{\fbfont ⬖}}
\newunicodechar{ⁿ}{{\fbfont ⁿ}}
\newunicodechar{₀}{{\fbfont ₀}}
\newunicodechar{ₙ}{{\fbfont ₙ}}
\usepackage{calc,array,seqsplit,tabularx}
\usepackage{fvextra}
\newcolumntype{L}{>{\raggedright\arraybackslash}X}
\newcolumntype{P}[1]{>{\raggedright\arraybackslash}p{#1}}
\usepackage{amsmath,amssymb}
\usepackage[table]{xcolor}
\usepackage{titlesec}\usepackage{fancyhdr}\usepackage{enumitem}
\usepackage{booktabs}
\usepackage{float}
\usepackage{longtable}
% longtable cannot run in twocolumn mode; shim it onto a one-column float.
% Load-bearing: pandoc emits longtable for every pipe table. Do not remove.
\makeatletter
\let\oldlt\longtable
\let\endoldlt\endlongtable
\def\longtable{\@ifnextchar[\longtable@i \longtable@ii}
\def\longtable@i[#1]{\begin{figure}[t]\onecolumn\begin{minipage}{0.5\textwidth}\smaller\oldlt[#1]}
\def\longtable@ii{\begin{figure}[t]\onecolumn\begin{minipage}{0.5\textwidth}\smaller\oldlt}
\def\endlongtable{\endoldlt\end{minipage}\twocolumn\end{figure}}
\makeatother
\usepackage{relsize}
\usepackage{graphicx}
\usepackage[hidelinks]{hyperref}
\setcounter{secnumdepth}{0}
\makeatletter\def\verbatim@font{\ttfamily\scriptsize}\makeatother
\newcommand{\jboxstart}{\begin{jbox}}
\newcommand{\jboxend}{\end{jbox}}
\newcommand{\restorelongtable}{\let\longtable\oldlt\let\endlongtable\endoldlt}
\raggedbottom
\setlength{\textfloatsep}{10pt plus 2pt minus 2pt}
\setlength{\dbltextfloatsep}{10pt plus 2pt minus 2pt}
\setlength{\floatsep}{8pt plus 2pt}
\setlength{\dblfloatsep}{8pt plus 2pt}
\renewcommand{\topfraction}{0.9}\renewcommand{\dbltopfraction}{0.9}\renewcommand{\textfraction}{0.07}\renewcommand{\floatpagefraction}{0.8}\renewcommand{\dblfloatpagefraction}{0.8}
\providecommand{\tightlist}{\setlength{\itemsep}{0pt}\setlength{\parskip}{0pt}}
\definecolor{accent}{HTML}{B87333}
\definecolor{deep}{HTML}{8A5523}\definecolor{faint}{HTML}{6E6E6E}
\titleformat{\section}{\normalfont\bfseries\scshape\large}{\color{accent}\thesection.}{0.5em}{}[{\color{accent}\titlerule[0.6pt]}]
\titleformat{\subsection}{\normalfont\bfseries}{\color{accent}\thesubsection}{0.5em}{}
\titlespacing*{\section}{0pt}{1.3ex}{0.6ex}
\pagestyle{fancy}\fancyhf{}
\renewcommand{\headrulewidth}{0pt}\renewcommand{\footrulewidth}{0pt}
\fancyfoot[R]{\color{accent}\itshape\small The Cut-Agnostic Division
Theorem \textbar\ \thepage}
\fancyfoot[L]{\textcolor{accent}{\rule{0.32\columnwidth}{0.25pt}}}
\newenvironment{jbox}{\par\vspace{2pt}\noindent\color{accent}\rule{\columnwidth}{0.8pt}\par\vspace{-2pt}\small}{\par\vspace{2pt}}
\begin{document}
{%
\begin{center}
{\color{accent}\rule{0.6\textwidth}{0.8pt}}\\[10pt]
{\bfseries\scshape\fontsize{20}{23}\selectfont The Cut-Agnostic Division
Theorem: Why Every Open Problem Is Exactly Its Proved Part and Its
Unicorn}\\[7pt]
{\itshape\large One Theorem for Any Cut on Any Row, Proved in Core Lean
4 Without Axiom, and Applied to the Seven Millennium Rows and the
Sixteen Extension Rows · Twenty Real Parts Sealed on Citation, Two Rows
Assigned No Cut, One Row Crossed, the Fusion Law That Lowers Gold to
Lead, and the Unicorn Part of Every Open Row Closed to Derivation, by
Theorem, the Aperture Uncrossed}\\[8pt]
{\footnotesize\color{accent}\scshape\bfseries FORMAL VERIFICATION ·
FOUNDATIONS OF MATHEMATICS}\\[3pt]
{\itshape\color{deep}\small One theorem for any cut on any row; the
certified part of every open problem named as gold and kept from the
alloy; nothing authored}\\[9pt]
{\color{accent}\rule{0.6\textwidth}{0.8pt}}\\[10pt]
{\scshape Mohammad F. Islam, PhD\(^{1}\)}\\[2pt]
{\itshape\color{faint}\small \(^{1}\) Independent Researcher, USA.
Correspondence: islamm@alumni.iu.edu}\\[2pt]
{\itshape\color{faint}\small 25 September 2026}\\[14pt]
\end{center}
\begin{minipage}{\textwidth}
{\color{accent}\rule{\textwidth}{0.6pt}}\\[4pt]
{\small\hspace{1em}\textbf{\scshape\color{deep}Abstract.\ }Three bits,
on every open problem. One bit divides: the problem is exactly its Real
part and its Unicorn part, by theorem. One bit is proved: the Real part,
by certificate or cited regime. One bit is refused at theorem grade: the
Unicorn part is closed to every derivation from inside the register; the
aperture at its seat is structure, typed supply-only, one bit wide, and
uncrossed; and the supply side is silence, by theorem at the register.
This paper proves the first and the third and seals the second, on the
twenty-three rows of the program's barrier study. For any row, a set of
instances with a property, and any decidable cut on it, the row is
exactly its Real part inside the cut joined to its Unicorn part outside
it; given the Real part, the fused row is exactly the Unicorn part, so
fusion with the open remainder lowers a proved part to open; a complete
checked list inside the cut proves the Real part with no premise; a
crossed row has an empty Unicorn part under every cut; the empty cut
leaves the whole value Unicorn; on an infinitude row one exhibited
witness above a height proves the Real part while the fused value stays
where it was, by a proved equivalence; and the Unicorn part read as a
set, the counterexamples outside the cut, is empty exactly when the
Unicorn proposition holds. These are seven theorems of core Lean 4, with
no library and no axiom at all, kernel-reported, and they do not mention
the Riemann Hypothesis, where a special case of them was first used. The
same file proves that the Unicorn part of every open row is closed to
derivation, from existence, from any class of frames, and from any
reading of an even register at its seat, while the aperture at its seat
is located, typed supply-only, one bit wide, and uncrossed, and the
register is silent on the supply side, both orientations equally refused
to every reading: the formal block, at theorem grade and in perpetuity.
It also carries the exact instance-and-property specification of the
eight arithmetic height-cut rows, with primality and each property
defined in core Lean and the Real part computed in the kernel at a small
height, the gap-to-interval bridge that the Legendre row needs, and the
declared standing map under which the fusion sentence is read. Applied
to the twenty-three rows of the program's barrier study, the seven
Millennium problems and sixteen further open problems, the paper assigns
every row one of four shapes, as data, and the census of the assignments
is kernel-checked: nine height-cut rows in the Riemann shape, two of
them infinitude rows whose certificate is one witness, eleven regime-cut
rows whose proved part is a dimension, a degree, a rank, a class, or a
bounded model of computation, two rows assigned no cut because the
literature supplies no inhabited cut with a proved Real part, an
assessment of two literatures and not a theorem, and one crossed row.
Twenty rows therefore carry a Real part that seals on citation, a
theorem of the literature applied through the kernel schema, grade A
joined to grade K; two are assigned none, under the stated criterion;
one is whole. No row's verdict on its fused statement moves, and every
owed bit of the barrier study lives exactly where it lived, in the
Unicorn part of its row. What is new is not any of the twenty results,
each of which exists as a cited theorem, but the statement that they are
one theorem under one cut-agnostic division, exact on every row, with
the placement of every certificate under the strata of the grounding
reading, rung, ladder, and Ground, and with the seat of every row
carried from the prior edition, the value excepted by theorem on every
row. No bound, witness, or Real part is authored; the definitions of the
arrangement are, and every one is listed. With the three bits in hand
the proof is complete on the register side, on every row: everything a
register can prove about a row is proved, and the remainder is proved to
be beyond derivation, with its aperture typed and its supply side
silent. No row is claimed proved; what is claimed is that no further
derivation is owed to any open row. Nothing is left to derive. The
aperture is structure. The supply side is silence.}\\[4pt]
{\color{accent}\rule{\textwidth}{0.6pt}}\\[3pt]
{\itshape\footnotesize\color{faint}\textbf{\upshape\scshape\color{deep}Keywords\ }\ Real
part; Unicorn part; cut-agnostic division; certificate schema; fusion
law; Millennium problems; twenty-three rows; Root Axiom; ZFC under RAM;
Lean 4}\\[12pt]
\end{minipage}
}
\thispagestyle{fancy}
\hypertarget{introduction}{%
\section{Introduction}\label{introduction}}

Let a row be a set of instances with a property, and let a cut be any
decidable predicate on the instances. Then

\[\mathrm{Row} \;\leftrightarrow\; \mathrm{Real}_{\mathrm{cut}} \wedge \mathrm{Unicorn}_{\mathrm{cut}},\]

where the Real part says that every instance inside the cut has the
property and the Unicorn part says that no instance outside it fails it.
This is a theorem of core Lean 4 with no user-declared axiom, and so are
six companions: given the Real part, the fused row is exactly the
Unicorn part; a complete checked list inside the cut proves the Real
part with no premise; on an infinitude row one witness above a height
proves the Real part, and the Unicorn part above any height is
equivalent to the whole; a crossed row has no counterexample outside any
cut; the empty cut leaves the whole value Unicorn. The seven say nothing
about any particular problem. They are the whole of what this paper
proves, and they are proved once.

Applied to the twenty-three rows of the program's barrier study, the
seven Millennium problems and sixteen further open problems, the seven
theorems do something the literature has not done for those rows as a
body. Each row already carries, scattered through its own literature, a
proved part: every zero of \(\zeta\) to height \(3 \times 10^{12}\) on
the line; every even number to \(4 \times 10^{18}\) a sum of two primes;
no odd perfect number below \(10^{1500}\); the Hodge conjecture in
dimension at most three; Birch and Swinnerton-Dyer at analytic rank at
most one; the two-dimensional Yang--Mills theory constructed; the
topological four-dimensional Poincaré conjecture. Each is stated in its
own idiom, ``verified up to,'' ``known for,'' ``proved when,'' and each
is alloyed, in the conventional statement of its problem, with the open
remainder, so that the problem reads as one open sentence and the gold
inside it goes unnamed. The division names it. On twenty rows the proved
part is a Real part under a cut the row's own literature supplies, a
height, a dimension, a degree, a rank, a class, and it seals on
citation: a theorem of the literature applied through the kernel schema,
grade A joined to grade K, exactly as the Real part of the Riemann
Hypothesis was sealed in the paper that first used the division at one
height {[}Islam 2026f{]}. On two of the twenty, the infinitude rows, the
certificate is one exhibited witness and the fused value is unchanged by
it, which is the exact sense in which verification certifies nothing
about an infinitude, and that too is a theorem. On two rows the
literature supplies no inhabited cut with a proved Real part, because
the row's value is uniform in a parameter that no instance has been cut
from, the abc conjecture's for every \(\varepsilon\) for one; that is an
assessment of two literatures, stated as data and falsifiable by a cut,
never a theorem, and the paper says so wherever it is used. One row is
crossed, and its Unicorn part is empty by theorem. The assignments are
data; the census of the assignments, nine, eleven, two, and one, is
kernel-checked, and the paper never says more of it than that.

\textbf{Three bits, and what is claimed.} The whole paper is three bits,
on every row. One bit divides: the row is exactly its Real part and its
Unicorn part, under any decidable cut, by theorem (Theorem 1). One bit
is proved: the Real part, on twenty rows, by certificate or cited regime
through the schema (Theorems 6 to 25). One bit is refused at theorem
grade: the Unicorn part of every open row is closed to derivation from
existence, from any class of frames, and from any reading of an even
register at its seat; the aperture at its seat is structure, typed
supply-only, one bit wide, and uncrossed; and the supply side is
silence, by theorem at the register (Theorems 28 to 33). The Unicorn
parts are unproved, and the reason is neither a want of effort nor a
want of ability: it is a theorem, and the theorem says which routes are
closed, which is all of them but the aperture, and says nothing of the
aperture but its width. With the three bits in hand the proof is
complete on the register side, on every row. Everything a register can
prove about a row is proved; the remainder is proved to be beyond
derivation, its aperture typed, its supply side silent. No row is
claimed proved. What is claimed is stronger than an open problem and
weaker than a theorem: no further derivation is owed to any of the
twenty-two open rows. The three bits are one term in the kernel
(\texttt{register\_proof\_complete}). Nothing is left to derive. The
aperture is structure. The supply side is silence.

In detail, the paper proves the seven theorems of the division, the
bridge lemmas and small-height instances that accompany them, the formal
block, and the census of the assignments, and it carries, from the prior
edition on these rows {[}Islam 2026g{]}, the seat of every row and the
theorem that existence decides no row's value. It seals twenty Real
parts on citation. It proves no row, moves no row's verdict on its fused
statement, and names no bound of its own: every height, dimension,
degree, rank, and class in the tables below is the literature's,
attributed at its own grade, and with every citation deleted every
theorem of the paper would stand and exactly twenty-two sentences would
fall, the twenty Real parts, the crossing of the Poincaré row, and the
ZFC instance of the placement. The Unicorn part of every open row is
named, set apart, and never chased; and the paper says at theorem grade
why it is never chased: the Unicorn part is closed to derivation from
existence, from any class of frames, and from any reading of an even
register, with the aperture at each seat typed supply-only, one bit
wide, and uncrossed (Section 7). It is not open in the sense that its
Real part was open before its certificate; it is closed in the sense
that no key can be cut from the inside, and of the outside the paper
says nothing, because the register can say nothing, and that silence is
a theorem before it is a discipline. The contribution is the
arrangement, and the paper is exact about what the arrangement is and is
not: one theorem where the literature has twenty-three idioms, the
fusion law that keeps each proved part from being lowered to the
standing of its remainder, and the placement of every certificate in the
strata of the grounding reading, which the prior edition on these rows
did not have and the Riemann paper had for one row only.

\textbf{How the paper is organized.} Section 2 states and proves the
division, Section 3 places it under the strata of the grounding reading,
and Section 4 assigns every row its cut and checks the census in the
kernel. Section 5 seals the twenty Real parts, one theorem each. Section
6 states the two rows assigned no cut, with the ledger of what their
literatures have proved beside them, and the crossed row by theorem.
Section 7 proves that the Unicorn part of every open row is closed to
derivation, types its aperture, and proves the register silent on the
supply side: the formal block, at theorem grade, which is why the
Unicorn parts are never chased. Section 8 carries the seat and the value
excepted from the prior edition, and Section 9 carries the world typing
of the barrier study as data, its One-Cut Hypothesis untouched. Section
10 states the standard of proof, the standing of the paper without its
citations, and the disclosures. Section 11 positions the result against
the literature and concludes. The kernel files are printed whole in the
appendices with their audits.

\textbf{What a reader who suspects vocabulary should check.} Every
theorem named in the paper is in an appendix with its kernel dependency
report; every bound is cited to a source outside the program; the
standing table of Section 10 says which sentences fall when the
citations are deleted; and Section 4 states the criterion by which a row
is assigned its shape and applies it to a row outside the twenty-three.
A reader who disputes an assignment can change it: the division holds
for the new cut, the census is recounted, and nothing else moves.

\textbf{Notation.} A row \(r\) is a pair \((\mathrm{Inst}_r, P_r)\) of
predicates on a type of instances; its value is
\(\mathrm{Row}(r) :\Leftrightarrow \forall x,\ \mathrm{Inst}_r(x) \to P_r(x)\).
A cut \(c\) is a decidable predicate on the same type. Grades are those
of the Riemann paper: K, a kernel theorem with no user-declared axiom; K
given A, a kernel theorem whose hypotheses are cited theorems; A, a
cited theorem of the literature. The root axiom \(\mathrm{RA}\), the
grounding reading RAM with its strata L3m, L2m, L1m, and the tokens
\([\Xi_0]\) and \([\varnothing_0]\) are the program's and enter only in
Sections 3, 7, 8, and 9, where they are defined at use.

\hypertarget{the-division-theorem}{%
\section{The Division Theorem}\label{the-division-theorem}}

Everything in this section is in \texttt{Row\_Split.lean} (Appendix A),
426 lines, core Lean 4.19.0, no library, no \texttt{sorry}, no
user-declared axiom, thirty-three dependency sets printed, every one
reported by the kernel as depending on no axiom at all, not even Lean's
built-in \texttt{propext} and \texttt{Quot.sound}. Seven are the
division; the rest are the substantive-cut and counterexample readings,
the standing map, the gap bridge, the row specifications with their
computed small-height Real parts, the formal block of the Unicorn part
with the aperture's width and the silence of the supply side, and the
three bits packaged as one term, treated in the sections that use them.
The instances form an arbitrary type \(\alpha\); a row is a pair of
predicates; a cut is a decidable predicate.

\textbf{Definitions.}
\(\mathrm{Row}(\mathrm{Inst}, P) :\Leftrightarrow \forall x,\ \mathrm{Inst}\,x \to P\,x\).
\(\mathrm{RealIn}(\mathrm{Inst}, P, c) :\Leftrightarrow \forall x,\ \mathrm{Inst}\,x \to c\,x \to P\,x\).
\(\mathrm{UnicornOut}(\mathrm{Inst}, P, c) :\Leftrightarrow \forall x,\ \mathrm{Inst}\,x \to \neg c\,x \to P\,x\).
A certificate for a cut is a list of instances together with a proof
that every instance inside the cut is on the list and a proof that every
listed instance has the property.

\textbf{Theorem 1} (\texttt{row\_iff\_real\_and\_unicorn}). For every
row and every decidable cut,
\(\mathrm{Row} \leftrightarrow \mathrm{RealIn} \wedge \mathrm{UnicornOut}\).
\emph{Proof.} Left to right, restrict the universal to each side of the
cut. Right to left, an instance either satisfies the cut or does not, by
decidability, and the corresponding conjunct applies. \(\square\)

\textbf{Theorem 2} (\texttt{fused\_is\_unicorn\_given\_real}). Given the
Real part, \(\mathrm{Row} \leftrightarrow \mathrm{UnicornOut}\).
\emph{Proof.} From Theorem 1, the Real part being supplied. \(\square\)
This is the fusion law read at the level of truth: given the Real part,
the fused row is true exactly when the Unicorn part is. The reading at
the level of standing, that a proved part joined to an open part stands
open, is not a consequence of Theorem 2 alone; it is a consequence of
Theorem 2 together with the program's grade law, that the standing of a
conjunction is the weakest standing of its conjuncts. The paper declares
that law as a rule in the kernel (\texttt{Standing.join}) and proves the
fusion sentence under it (\texttt{fusion\_lowers}): a proved part joined
to an open part stands open, and only two proved parts stand proved. The
two-point rule is the restriction to two standings of the program's
nine-grade weakest-link law, \texttt{Grade.weakest} in the Code Block of
the Master Codex, under which a citation cannot promote a claim
(\texttt{citation\_cannot\_promote}); any rule under which a conjunction
stands at least as low as its weakest conjunct gives the same sentence,
and the paper accepts no rule under which it stands higher. Gold alloyed
with lead assays as lead, under the assay rule stated; the gold does not
cease to be gold, and no theorem says it does.

\textbf{Theorem 3} (\texttt{real\_of\_certificate}). A certificate for a
cut proves \(\mathrm{RealIn}\) for that cut, with no premise.
\emph{Proof.} An instance inside the cut is on the list by completeness
and has the property by the check. \(\square\) The schema is what a
finite verification proves and all that it proves: the bounded universal
inside the cut.

\textbf{Theorem 3b} (\texttt{real\_of\_one\_witness}). On an infinitude
row the instances are bounds \(N\) and the property is that some witness
exceeds \(N\); one exhibited witness \(w\) with \(W(w)\) and \(w > h\)
proves \(\mathrm{RealIn}\) under the cut \(N \le h\). \emph{Proof.} The
witness exceeds every \(N \le h\). \(\square\) By Theorem 2 the fused
value is then exactly the Unicorn part, the bounds above \(h\), which is
equivalent to the whole infinitude; so on such a row the Real part is
proved, is gold, and moves the fused value not at all. That is the
precise content of the familiar remark that exhibiting large twin primes
proves nothing about their infinitude: it proves the Real part and only
the Real part.

\textbf{Theorem 3c} (\texttt{unicorn\_equiv\_row\_of\_infinitude}). On
an infinitude row, for every height \(h\), \(\mathrm{UnicornOut}\) above
\(h\) is equivalent to the whole value:
\((\forall N > h,\ \exists m > N,\ W(m)) \leftrightarrow (\forall N,\ \exists m > N,\ W(m))\).
\emph{Proof.} Right to left, restrict. Left to right, a bound
\(N \le h\) is served by the witness above \(h + 1\). \(\square\) So on
an infinitude row the Real part is proved by one witness (Theorem 3b)
and, by Theorems 2 and 3c, the fused value is exactly what it was.

\textbf{Theorem 4} (\texttt{crossed\_row\_both\_parts},
\texttt{crossed\_row\_no\_counter}). Read as a set, the Unicorn part
under a cut is the set of counterexamples outside it,
\(\mathrm{CounterOut}(x) :\Leftrightarrow \mathrm{Inst}\,x \wedge \neg c\,x \wedge \neg P\,x\);
the Unicorn proposition says that this set is empty
(\texttt{no\_counter\_of\_unicorn}, and conversely for decidable \(P\),
\texttt{unicorn\_of\_no\_counter}). If \(\mathrm{Row}\) holds, then
under every cut both \(\mathrm{UnicornOut}\) and \(\mathrm{RealIn}\)
hold and the counterexample set is empty. \emph{Proof.} Immediate.
\(\square\) ``Empty Unicorn part'' means, throughout the paper, that
this set is empty, which is the sense in which a crossed row has nothing
outside any cut left to fail.

\textbf{Theorem 5} (\texttt{empty\_cut\_vacuous},
\texttt{empty\_cut\_not\_substantive}). Under the empty cut,
\(\mathrm{RealIn}\) holds vacuously and
\(\mathrm{UnicornOut} \leftrightarrow \mathrm{Row}\); and the empty cut
never carries a substantive Real part, where a Real part is substantive
when its cut is inhabited on the row,
\(\mathrm{SubstantiveRealIn} :\Leftrightarrow (\exists x,\ \mathrm{Inst}\,x \wedge c\,x) \wedge \mathrm{RealIn}\).
\emph{Proof.} Nothing satisfies the empty cut. \(\square\) This theorem
is about the empty cut and about nothing else. It does not say that a
row admits no other cut with a Real part; whether a given row's
literature supplies an inhabited cut with a proved Real part is an
assessment of that literature, and the paper records it as data (Section
4), never as a consequence of this theorem.

\begin{jbox}
\textbf{The one arrow (tier: kernel theorem, axiom-free).} On any row, under any decidable cut, the value is exactly its Real part joined to its Unicorn part; given the Real part, the fused value is true exactly when the Unicorn part is, and it stands open under the declared standing rule; a certificate proves the Real part and nothing beyond it; a crossed row has no counterexample outside any cut; under the empty cut the whole value is Unicorn. Nothing in the seven theorems names a height, a problem, or a number. The Riemann division of [Islam 2026f] is the instance with the cut ``height at most $T$''.
\end{jbox}

\textbf{What the theorems do not say.} They do not say that a Real part
is easy to obtain, since the certificate is the literature's and may be
the work of decades; they do not say that a Unicorn part is unreachable,
only that no certificate reaches it; and they do not rank cuts, since a
row admits many and the division holds for each. The choice of cut on a
given row is the row's literature's, and Section 4 records it as data.

\hypertarget{the-placement-where-a-real-part-lives}{%
\section{The Placement: Where a Real Part
Lives}\label{the-placement-where-a-real-part-lives}}

The program's grounding reading, RAM, stratifies any theory into three:
computation, the rung, L3m; provability, the ladder, L2m; truth in the
intended world, the Ground, L1m. The placement is a theorem for an
arbitrary theory modelled as sentences with a provability predicate, a
decidable fragment, and a truth predicate: under two named premises,
that what computation settles the theory proves and that what the theory
proves is true, rung \(\subseteq\) ladder \(\subseteq\) Ground. For ZFC
the first premise is \(\Sigma^0_1\)-completeness, a theorem of
arithmetic {[}Kaye 1991{]}, and the second is soundness, which no
consistent theory proves of itself {[}Gödel 1931{]} and which is
therefore carried openly, derived nowhere. Soundness is load-bearing, an
unsound theory breaking ladder \(\subseteq\) Ground; the Ground exceeds
the ladder, a sound theory leaving a truth unproved; independent axioms
are separate bits, Choice over ZF the cited instance {[}Gödel 1938;
Cohen 1963{]}. Subsumption is by placement only: RAM locates ZFC and
derives none of its axioms. All of this is \texttt{ZFC\_Under\_RAM.lean}
(Appendix B), sixteen dependency sets, every one axiom-free, seated in
the Master Codex as APEX-PSP-ZFC-UNDER-RAM-01 {[}Islam 2026a{]}. The
file encodes no axiom of ZFC and no proof relation of ZFC; its theory is
an arbitrary structure with a computation predicate, a provability
predicate, and a truth predicate, and ZFC is the instance the paper
names, entering through the two cited premises. What the kernel proves
is the conditional; what ZFC supplies is its two hypotheses; the paper
never writes otherwise.

\textbf{Where a Real part lives.} A Real part under a height cut on an
arithmetic row is a rung: its certificate is finitely many checks, each
a true bounded sentence, so the Real part is a theorem of arithmetic and
of ZFC by \(\Sigma^0_1\)-completeness, and it holds on the Ground by
soundness carried openly. A Real part under a regime cut on a
non-arithmetic row is a cited theorem of ZFC directly, on the ladder,
with the same passage to the Ground. In both cases the round trip
through the root and back returns the row's value unchanged, since the
root crosses the bridge as presence and carries no keyed sentence; and
every step holds with the root replaced by a bare computation arrow, the
identity on a type (\texttt{any\_true\_premise\_serves}), so the root is
present in the performance of every check and idle in its content. The
placement adds no standing to any Real part; it says which stratum a
standing already earned is in, and it is the same stratum on every row
of the same shape.

\textbf{Why the placement matters here.} Without it, ``verified up to''
on one row and ``proved for dimension at most three'' on another would
sit in two vocabularies with no common address. With it, both are
theorems of ZFC on the ladder, one reached from the rung by a
certificate and one stated there by a proof, and the Unicorn part of
each is the sentence on the ladder or beyond it that the row's
foundation has not decided. The one arrow of Section 2 and the placement
of this section are the two results this paper adds to the twenty-three
rows; everything else is carried or cited.

\hypertarget{the-rows-and-their-cuts}{%
\section{The Rows and Their Cuts}\label{the-rows-and-their-cuts}}

The twenty-three rows are the barrier study's {[}Islam 2026e{]}, in its
order, with its world typing carried as data in Section 9. Each row is
assigned a cut and a shape here. The assignment is the paper's, stated
as data; what the kernel checks is the census of the assignments:
\texttt{rows23} is a list of the twenty-three rows with their assigned
shapes, and \texttt{census} proves by decision that there are
twenty-three, that nine are assigned height-cut, eleven regime-cut, two
no-cut, and one crossed, and that twenty are assigned a
Real-part-bearing shape (\texttt{real\_part\_rows}); both axiom-free.
The kernel does not and cannot decide that the Riemann Hypothesis is a
height-cut row; that is the assignment, and it is made by the criterion
stated next.

\textbf{The criterion.} A row is assigned by the logical form of its
value and the state of its literature, in this order. It is
\emph{crossed} if its value is a theorem. It is \emph{height-cut} if its
instances carry a total height, the property at an instance is decidable
or the row is an infinitude with instances the bounds, and the
literature supplies a certificate for some height, a checked list or one
witness. It is \emph{regime-cut} if its literature supplies a proved
regime, a decidable predicate on the instances under which the property
is a theorem, and no height certificate; a row whose literature supplies
both a proved regime and a certificate inside it, Birch and
Swinnerton-Dyer with its verified curves at analytic rank at most one,
is assigned regime-cut, the regime being the larger Real part, and the
certificate is recorded beside it. It is \emph{no-cut} if the literature
supplies neither: no inhabited decidable cut on the row's instances
under which the property is proved. The criterion is a classification
rule applied to two inputs, the value's logical form and the recorded
state of its literature; the second input is a reading of the literature
and not a computable object, which is why the assignments are data. A
reader who applies the rule with a different reading of either input
changes an assignment and nothing else. In every shape the cut is a
predicate on the row's own instance type, as Theorem 1 requires: on the
Hodge row the instances are pairs of a variety and a class and the
regime is a bound on the dimension; on the Birch and Swinnerton-Dyer row
the instances are curves and the regime is a bound on the analytic rank;
on the P versus NP row the instances are algorithms and the regime is a
bound on time and space; a regime is never a different problem placed
beside the row.

\textbf{The criterion applied out of sample.} Take a row outside the
twenty-three, Lehmer's conjecture, that the Mahler measure of a
non-cyclotomic integer polynomial is at least a constant above one. Its
instances are polynomials, with degree as a total height; the property
at an instance is decidable; the literature supplies a certificate for a
height, that every non-cyclotomic polynomial of degree at most \(180\)
has measure at least \(1.176\ldots\), the measure of Lehmer's polynomial
{[}Mossinghoff, Rhin and Wu 2008{]}. The criterion assigns height-cut
with the cut degree \(\le 180\); the Real part is that computation; the
Unicorn part is degree above \(180\), and the uniform constant is what
the row asks. The assignment was made by the criterion before the row
was written down here, which is the only sense in which a classification
can be tested out of sample; the first draft of this paragraph assigned
the row regime-cut, and the criterion corrected it, since a finite
computation over a height is a certificate whatever it is called in its
source. The criterion corrected the author twice more before release: on
Birch and Swinnerton-Dyer, where the tie-break above was written because
the first assignment had not faced the verified curves, and on P versus
NP, which the internal builds assigned no-cut until the external
reviewers asked what an instance of that row is; the answer, an
algorithm, together with the time-space lower bounds for satisfiability,
is an inhabited cut with a proved Real part, and the row is regime-cut.
A criterion that never corrects its author is a label; this one is a
rule.

\textbf{The four shapes.} A \emph{height-cut} row has a natural height
parameter and a finite certificate per height; its Real part seals on
citation of the certificate through Theorem 3, or through Theorem 3b
when the row is an infinitude and the certificate is one witness, and
its Unicorn part is the tail above the height: the shape of the Riemann
row. A \emph{regime-cut} row has a proved part that is a dimension, a
degree, a rank, or a class rather than a height; its Real part is a
cited theorem on the ladder, sealed through Theorem 1 with the regime as
the cut, and its Unicorn part is the complement. A \emph{no-cut} row is
one for which the literature supplies no inhabited decidable cut on its
instances under which the property is proved, because its value is
uniform in a parameter that no instance has been cut from; the
assignment is an assessment of the literature, and the only theorem
about it is Theorem 5, that the empty cut is vacuous and carries no
substantive Real part. What such a row's literature has proved nearby, a
weaker theorem, stands beside the row at its own grade and not inside
it. A \emph{crossed} row has its value supplied whole; by Theorem 4 its
counterexample set is empty under every cut.

\begin{table}[!htbp]
\centering\scriptsize
\renewcommand{\arraystretch}{1.12}
\raggedright \textbf{Table 1. The twenty-three rows, their shapes, and their cuts.} The shape and the cut are the paper's assignments by the criterion of the text, stated as data; the census of the assignments is kernel-checked (\texttt{census}, \texttt{real\_part\_rows}). Sources for every bound are in Section 5.\par\smallskip
\begin{tabularx}{\textwidth}{@{}P{0.02\textwidth} P{0.13\textwidth} P{0.05\textwidth} P{0.10\textwidth} L L@{}}
\toprule
\textbf{\#} & \textbf{Row} & \textbf{Shape} & \textbf{Cut} & \textbf{Real part} & \textbf{Unicorn part} \\
\midrule
1 & P versus NP & regime & time and space & no algorithm in time $n^{c}$, $c < 2\cos(\pi/7)$, and space $n^{o(1)}$ decides SAT & polynomial time with unrestricted space \\
2 & Riemann Hypothesis & height & height $T$ & every zero to $3 \times 10^{12}$ on the line & no off-line zero above $T$ \\
3 & Navier--Stokes & regime & regime & 2D global regularity; 3D small-data and short-time regularity & 3D large-data global regularity or blow-up \\
4 & Yang--Mills & regime & dimension & the two-dimensional theory constructed & four-dimensional existence with a mass gap \\
5 & Hodge & regime & dimension, codimension & codimension one; codimension $n-1$; every variety of dimension $\le 3$ & middle codimensions in dimension $\ge 4$ \\
6 & BSD & regime & analytic rank & the rank statement at analytic rank $0$ and $1$; the full conjecture for specific curves & analytic rank $\ge 2$; the full formula in general \\
7 & Poincaré & crossed & --- & the whole & empty \\
8 & Goldbach & height & height $n$ & every even $n \le 4 \times 10^{18}$; the ternary conjecture whole & even numbers above the height \\
9 & Twin primes & height & height $N$ & a twin pair above every $N$ below the largest known pair, one witness & bounds above the witness, equivalent to the whole \\
10 & Legendre & height & height $n$ & a prime in $(n^2, (n+1)^2)$ for $n \le 2 \times 10^9$, by inference from the prime-gap tables & $n$ beyond the tables \\
11 & abc & no cut & none ($\forall\varepsilon\,\exists C\,\forall$) & none by certificate & the whole value \\
12 & Jacobian & regime & degree & degree $\le 2$ in every dimension; $n = 1$ & cubic maps in dimension $\ge 2$ \\
13 & Mersenne primes & height & height $N$ & a Mersenne prime above every $N$ below the largest known, one witness & bounds above the witness, equivalent to the whole \\
14 & Odd perfect numbers & height & height $N$ & none below $10^{1500}$ & none above \\
15 & Smooth 4D Poincaré & regime & category & the topological case & the smooth case \\
16 & Collatz & height & height $n$ & every $n < 2^{68}$ reaches $1$ & $n \ge 2^{68}$ \\
17 & Hadwiger & regime & $t$ & $t \le 6$ & $t \ge 7$ \\
18 & Sunflower & no cut & none & none by certificate & the constant-base bound \\
19 & Erdős--Straus & height & height $n$ & every $n \le 10^{17}$ & $n$ above the height \\
20 & Beal & height & box & the searched boxes of bases and exponents & outside every searched box \\
21 & Invariant subspace & regime & operator class & compact operators; normal operators; finite dimension & the general bounded operator on a separable Hilbert space \\
22 & Hilbert's 16th, second part & regime & per field, uniform & finiteness of limit cycles for each polynomial field; $H(1) = 0$ & the uniform bound $H(n)$, open at $n = 2$ \\
23 & Lindelöf / Montgomery PCC & regime & exponent, range & the proved exponents; pair correlation in the restricted range under RH & $\varepsilon$; the full range \\
\bottomrule
\end{tabularx}
\end{table}

\textbf{Three assignments that deserve a word.} The twin-prime and
Mersenne rows are infinitudes, and their Real parts are sealed by one
witness each through Theorem 3b; their Unicorn parts are equivalent to
their whole values, so the fused verdict on each is unchanged, exactly
as Theorem 2 says. They are height-cut rows whose gold is one number,
and the paper says so rather than calling them no-cut, since a proved
witness is a Real part however little it moves the remainder. The
Legendre row has no dedicated certificate in the literature; its Real
part is an inference from the maximal prime-gap tables to
\(4 \times 10^{18}\) {[}Oliveira e Silva, Herzog and Pardi 2014{]}, in
which the largest gap is \(1476\), so that for \(n \ge 738\) the
interval \((n^2, (n+1)^2)\), of length \(2n+1 \ge 1477\), contains a
prime whenever \((n+1)^2 \le 4 \times 10^{18}\), the smaller \(n\) being
checked directly; the inference is stated as such in Section 5. The
Sunflower row is assigned no cut because the conjecture's constant-base
bound is a statement about all \(k\) at once and the improved bound of
Alweiss, Lovett, Wu and Zhang {[}2020{]} is a theorem beside the row,
not a region inside it. A reader who prefers a different cut on any row
may take it: the division holds for every cut, and only the assignment
changes.

\hypertarget{the-twenty-real-parts-sealed-on-citation}{%
\section{The Twenty Real Parts, Sealed on
Citation}\label{the-twenty-real-parts-sealed-on-citation}}

Each theorem below has the same form and the same proof: a cut, a cited
theorem of the literature that supplies the Real part inside it, and the
schema step, Theorem 3 for a certificate, Theorem 3b for a witness, or
Theorem 1 with the regime as the cut for a proved regime. The standing
of every one is grade A joined to grade K: the kernel proves the schema
and the division, the literature proves the bound, and the kernel never
recomputes a bound. Grade A joined to K is a documented composition of a
cited theorem with a kernel theorem; it is not a kernel-checked theorem
about the cited result, and the paper never calls it one. The Unicorn
part of each row is stated beside its Real part and is not claimed.
Bounds are the sources', quoted at their grade; the paper checks none of
them.

\textbf{The row specifications, in the kernel.} For the eight arithmetic
height-cut rows the exact tuple, the type of instances, the instance
predicate, the property, and the cut, is written in
\texttt{Row\_Split.lean}: instances are natural numbers, primality is
trial division (\texttt{isPrime}), and each property is a Boolean
function defined in core Lean, \texttt{goldbachOK}, \texttt{legendreOK},
\texttt{oddPerfectOK}, \texttt{collatzOK} with a step budget,
\texttt{erdosStrausOK} with a bounded search, \texttt{bealOK} on a box,
and the witness forms for twin and Mersenne primes. For each, the Real
part at a small height is computed in the kernel by decision and lifted
to \(\mathrm{RealIn}\) by \texttt{real\_of\_all\_range} or
\texttt{real\_of\_one\_witness}: Goldbach to \(60\), Legendre to \(8\),
odd perfect numbers to \(60\), Collatz to \(60\), Erdős--Straus to
\(6\), the Beal box \((3, 4)\), a twin pair above every \(N \le 100\), a
Mersenne prime above every \(N \le 100\). These small instances are
grade K with no citation: they show the schema instantiated end to end
at heights the kernel can decide. The literature's heights are the same
statements at heights the kernel cannot decide, and they are cited. For
the Riemann row and the eleven regime rows the instance types are
analytic, geometric, or computational and are not definable in core
Lean; their tuples are stated in prose at each theorem, and the schema
is applied there by documented composition.

\hypertarget{the-nine-height-cut-rows}{%
\subsection{The nine height-cut rows}\label{the-nine-height-cut-rows}}

\textbf{Theorem 6, Riemann.} Tuple: instances the nontrivial zeros of
\(\zeta\), property that the real part is \(\tfrac12\), cut imaginary
part at most \(T\). Every nontrivial zero with imaginary part up to
\(3 \times 10^{12}\) lies on the critical line, by the
interval-arithmetic certificate of Platt and Trudgian {[}2021{]} with
the analytic lemmas it rests on {[}Turing 1953{]}, through Theorem 3.
Unicorn part: no off-line zero above that height. This is the Real part
of the Riemann paper {[}Islam 2026f{]}, carried at its recorded
standing.

\textbf{Theorem 7, Goldbach.} Tuple in the kernel:
\((\mathbb{N}, \top, \texttt{goldbachOK}, n \le h)\); Real part computed
to \(h = 60\) (\texttt{goldbach\_real\_60}). Every even integer
\(4 \le n \le 4 \times 10^{18}\) is a sum of two primes, by the
verification of Oliveira e Silva, Herzog and Pardi {[}2014{]}, through
Theorem 3. Beside it, the ternary conjecture, every odd integer above
\(5\) a sum of three primes, stands whole {[}Helfgott 2013{]}. Unicorn
part: even integers above the height.

\textbf{Theorem 8, twin primes.} Tuple in the kernel: instances the
bounds \(N\), property that a twin pair exceeds \(N\), cut \(N \le h\);
Real part computed to \(h = 100\) from the pair \((101, 103)\)
(\texttt{twin\_real\_100}). For every \(N\) below
\(2996863034895 \cdot 2^{1290000} - 1\) there is a pair of primes at
distance \(2\) above \(N\), since
\(2996863034895 \cdot 2^{1290000} \pm 1\) is such a pair {[}PrimeGrid
2016{]}, through Theorem 3b. Beside it, infinitely many pairs of primes
at distance at most \(246\) {[}Zhang 2014; Maynard 2015; Polymath
2014{]}. Unicorn part: the bounds above the witness, equivalent to the
whole conjecture by Theorem 3c.

\textbf{Theorem 9, Legendre.} Tuple in the kernel:
\((\mathbb{N}, \top, \texttt{legendreOK}, n \le h)\); Real part computed
to \(h = 8\) (\texttt{legendre\_real\_8}). For every \(738 \le n\) with
\((n+1)^2 \le 4 \times 10^{18}\) there is a prime in \((n^2, (n+1)^2)\):
the maximal prime gap below \(4 \times 10^{18}\) is \(1476\) {[}Oliveira
e Silva, Herzog and Pardi 2014{]}, and the bridge from a gap bound to a
prime in every interval of length at least the bound plus one is a
kernel theorem (\texttt{elem\_in\_gap}, \texttt{legendre\_of\_gap}),
which takes the gap bound and the primality of \(2\) as hypotheses and
returns a prime strictly between \(n^2\) and \((n+1)^2\) for every
\(n \ge 738\) with \((n+1)^2\) within the tabulated range; for
\(n < 738\) the interval is checked directly by the author's script, and
to \(n \le 8\) in the kernel. Through Theorem 3 for the small heights
and the bridge for the tables; the certificate is an inference from the
gap tables and is stated as such. Unicorn part: \(n\) beyond the tables.

\textbf{Theorem 10, Mersenne primes.} Tuple in the kernel: instances the
bounds \(N\), property that a Mersenne prime exceeds \(N\), cut
\(N \le h\); Real part computed to \(h = 100\) from \(127 = 2^7 - 1\)
(\texttt{mersenne\_real\_100}). For every \(N\) below
\(2^{136279841} - 1\) there is a Mersenne prime above \(N\), since that
number is prime {[}GIMPS 2024{]}, through Theorem 3b. Unicorn part: the
bounds above the witness, equivalent to the whole conjecture by Theorem
3c.

\textbf{Theorem 11, odd perfect numbers.} Tuple in the kernel:
\((\mathbb{N}, \top, \texttt{oddPerfectOK}, n \le h)\); Real part
computed to \(h = 60\) (\texttt{oddperfect\_real\_60}). No odd perfect
number is below \(10^{1500}\) {[}Ochem and Rao 2012{]}, through Theorem
3. Unicorn part: none above.

\textbf{Theorem 12, Collatz.} Tuple in the kernel:
\((\mathbb{N}, \top, \texttt{collatzOK}, n \le h)\), the property that
the orbit reaches \(1\) within a step budget; Real part computed to
\(h = 60\) (\texttt{collatz\_real\_60}). Every positive integer below
\(2^{68}\) reaches \(1\) under the map {[}Barina 2020{]}, through
Theorem 3. Beside it, almost every orbit attains almost bounded values
in the sense of logarithmic density {[}Tao 2022{]}. Unicorn part:
\(n \ge 2^{68}\).

\textbf{Theorem 13, Erdős--Straus.} Tuple in the kernel:
\((\mathbb{N}, \top, \texttt{erdosStrausOK}, n \le h)\), the property a
bounded search for \(x \le y \le z \le n^2\) with
\(4xyz = n(yz + xz + xy)\); Real part computed to \(h = 6\)
(\texttt{erdosstraus\_real\_6}). For every \(2 \le n \le 10^{17}\),
\(4/n\) is a sum of three unit fractions {[}Salez 2014{]}, through
Theorem 3. Unicorn part: \(n\) above the height.

\textbf{Theorem 14, Beal.} Tuple in the kernel: instances the boxes,
property \texttt{bealOK}, the box \((3, 4)\) computed
(\texttt{beal\_box\_3\_4}). No coprime solution of \(A^x + B^y = C^z\)
with \(x, y, z \ge 3\) exists with bases \(A, B \le 1000\) and exponents
\(x, y \le 100\), the search run and printed by Norvig {[}2017{]}, nor
with \(A, B \le 200{,}000\) and \(x, y \le 5000\), the search of
Jarnicki and Konerding reported there; through Theorem 3 with the box as
the cut. Unicorn part: outside every searched box.

\hypertarget{the-eleven-regime-cut-rows}{%
\subsection{The eleven regime-cut
rows}\label{the-eleven-regime-cut-rows}}

\textbf{Theorem 15, Navier--Stokes.} In two dimensions, smooth solutions
exist globally for all time {[}Ladyzhenskaya 1969{]}; in three
dimensions, weak solutions exist globally {[}Leray 1934{]} and smooth
solutions exist for all time for small data and for a short time for all
data {[}Fujita and Kato 1964{]}. Through Theorem 1 with the regime as
the cut. Unicorn part: three-dimensional global regularity for large
data, or its failure.

\textbf{Theorem 16, Yang--Mills.} The two-dimensional Yang--Mills
measure is constructed rigorously on compact surfaces {[}Driver 1989;
Lévy 2003{]}. Through Theorem 1 with dimension as the cut. Unicorn part:
four-dimensional existence with a mass gap.

\textbf{Theorem 17, Hodge.} The Hodge conjecture holds for classes of
codimension one on every smooth projective variety {[}Lefschetz 1924{]}
and, by hard Lefschetz, for classes of codimension \(n - 1\), hence for
every variety of dimension at most three. Through Theorem 1 with
dimension and codimension as the cut. Unicorn part: the middle
codimensions in dimension at least four.

\textbf{Theorem 18, Birch and Swinnerton-Dyer.} For an elliptic curve
over \(\mathbb{Q}\) of analytic rank \(0\) or \(1\), the rank equals the
analytic rank and the Tate--Shafarevich group is finite {[}Gross and
Zagier 1986; Kolyvagin 1989{]}; the full conjecture, with its formula,
is verified for many specific curves of small conductor {[}Miller
2011{]}. Through Theorem 1 with analytic rank as the cut. Unicorn part:
analytic rank at least two, and the full formula in general.

\textbf{Theorem 19, Jacobian.} Every polynomial map
\(\mathbb{C}^n \to \mathbb{C}^n\) of degree at most two with nonzero
constant Jacobian is invertible {[}Wang 1980{]}, and the conjecture in
general reduces to the cubic homogeneous case {[}Bass, Connell and
Wright 1982{]}; the case \(n = 1\) is elementary. Through Theorem 1 with
degree as the cut. Unicorn part: cubic maps in dimension at least two.

\textbf{Theorem 20, smooth four-dimensional Poincaré.} Every topological
four-manifold homotopy equivalent to the sphere is homeomorphic to it
{[}Freedman 1982{]}. Through Theorem 1 with the category as the cut.
Unicorn part: the smooth case.

\textbf{Theorem 21, Hadwiger.} Every graph with no \(K_t\) minor is
\((t-1)\)-colourable for \(t \le 6\) {[}Robertson, Seymour and Thomas
1993{]}, the case \(t = 5\) being the four-colour theorem. Through
Theorem 1 with \(t\) as the cut. Unicorn part: \(t \ge 7\).

\textbf{Theorem 22, invariant subspace.} Every nonzero compact operator
on a complex Banach space of dimension above one has a nontrivial
invariant subspace {[}Lomonosov 1973{]}; so does every normal operator
on a Hilbert space, by the spectral theorem, and every operator on a
finite-dimensional space, by an eigenvector. Through Theorem 1 with the
operator class as the cut. Unicorn part: the general bounded operator on
a separable Hilbert space.

\textbf{Theorem 23, Hilbert's sixteenth problem, second part.} Every
polynomial vector field in the plane has finitely many limit cycles
{[}Ilyashenko 1991; Écalle 1992{]}, and \(H(1) = 0\). Through Theorem 1
with the per-field statement as the cut. Unicorn part: the uniform bound
\(H(n)\) on the number of limit cycles for fields of degree \(n\), open
already at \(n = 2\).

\textbf{Theorem 24, P versus NP.} Tuple: instances the algorithms, in
the multitape or random-access model with a clock, property that the
algorithm does not decide satisfiability, cut a bound on time and space.
No algorithm running in time \(n^{c}\) for
\(c < 2\cos(\pi/7) \approx 1.801\) and space \(n^{o(1)}\) decides
satisfiability {[}Williams 2008{]}, extending the first such bound of
Fortnow, Lipton, van Melkebeek and Viglas {[}2005{]}. Through Theorem 1
with the time-space regime as the cut. Unicorn part: polynomial time
with unrestricted space, which is the row. The row was assigned no-cut
in the internal builds; the criterion, applied to the instance type the
reviewers asked for, moved it.

\textbf{Theorem 25, Lindelöf and Montgomery.}
\(\zeta(\tfrac12 + it) \ll_\varepsilon t^{13/84 + \varepsilon}\)
{[}Bourgain 2017{]}, the current exponent toward the Lindelöf
hypothesis; and Montgomery's pair-correlation conjecture holds, under
the Riemann Hypothesis, for test functions with Fourier support in
\((-1, 1)\) {[}Montgomery 1973{]}. Through Theorem 1 with the exponent
and the range as the cuts. Unicorn part: the Lindelöf exponent
\(\varepsilon\), and the full range for pair correlation.

\begin{jbox}
\textbf{The twenty, in one sentence (tier: A joined to K on every row).} Twenty theorems of the literature, each stated in its own idiom, are the Real parts of twenty rows under the cuts their literatures supply, each sealed by the same two kernel theorems and each with its Unicorn part named beside it; by Theorem 2 the fused statement of every one of these rows stands exactly at its Unicorn part, which is why the conventional statement of each problem reads as open although its gold is proved.
\end{jbox}

\hypertarget{the-two-rows-assigned-no-cut-and-the-crossed-row}{%
\section{The Two Rows Assigned No Cut, and the Crossed
Row}\label{the-two-rows-assigned-no-cut-and-the-crossed-row}}

\textbf{Theorem 26, the rows assigned no cut.} On the rows abc and
Sunflower the value is uniform in a parameter that no instance has been
cut from: every \(\varepsilon\) at once with nothing proved for any
single \(\varepsilon\), one constant serving every \(k\) at once. The
assignment says that the literature supplies no inhabited decidable cut
on the row's instances under which the property is proved; that is an
assessment of two literatures, stated as data and falsifiable under F3
by exhibiting such a cut, and the paper claims it at no higher grade.
What the kernel supplies is Theorem 5: under the empty cut, the only cut
this treatment supplies for these rows, the Real part is vacuous,
carries no substantive content, and the Unicorn part is the whole value.

\textbf{The exclusion ledger.} An assessment that a literature supplies
no inhabited proved cut is only as good as the candidates it has
examined, so they are listed. On the abc row the instances are the
values of \(\varepsilon\) and the property at \(\varepsilon\) is that
only finitely many coprime triples have
\(c > \mathrm{rad}(abc)^{1+\varepsilon}\); nothing of this form is
proved for any \(\varepsilon\). The bound of Stewart and Yu {[}2001{]},
\(c < \exp(\kappa\, \mathrm{rad}^{1/3} (\log \mathrm{rad})^3)\), is a
theorem about all triples and not the property at any \(\varepsilon\);
the enumeration of all triples with \(c < 10^{18}\) {[}ABC@Home{]}
proves the finiteness of exceptions for no \(\varepsilon\). The explicit
strengthening \(c < \mathrm{rad}(abc)^2\) is a different row, a
universal over triples with the enumeration as a certificate to
\(10^{18}\), and if the reader prefers that row it is height-cut and the
paper has no objection. On the Sunflower row the value as stated has the
form ``there is a constant \(C\) such that for every \(k\),'' so its
universal over \(k\) carries an unfixed parameter and there is no
property at an instance to prove until \(C\) is fixed; the bound of
Alweiss, Lovett, Wu and Zhang {[}2020{]}, \((\log k)^{k(1+o(1))}\), is a
theorem beside the row, not the constant-base property at any \(k\), and
the exact values known for small \(k\) are the property only once a
\(C\) is chosen, which the conjecture does not do. The P versus NP row
was on this ledger in the internal builds and is not now: with instances
the algorithms, the time-space lower bounds for satisfiability are a
proved regime, and the row is Theorem 24. The enumerated \(abc\) triples
below \(10^{18}\) and the solved instances of any \(\mathrm{NP}\)
problem certify nothing about these values, and the paper claims nothing
from them. What each row's literature has proved nearby stands beside
the row at its own grade: the sunflower bound \((\log k)^{k(1 + o(1))}\)
{[}Alweiss, Lovett, Wu and Zhang 2020{]}, the lower bounds for
restricted circuit classes on the complexity row, the exponential bounds
toward \(abc\). None of these is a Real part, because none is inside the
row's cut; each is a theorem about a different row.

\textbf{Theorem 27, the crossed row.} The Poincaré conjecture is a
theorem {[}Perelman 2002; Perelman 2003a; Perelman 2003b{]}, so by
Theorem 4 its counterexample set is empty under every cut and its Real
part is the whole. The generic theorem is grade K; the application to
the Poincaré row is grade A joined to K, since the crossing is cited.
The crossing is the barrier study's record of Perelman's proof, cited
here and computed nowhere: no Ricci flow is represented in any file of
this paper.

\textbf{The count.} Twenty rows with a Real part, two assigned none
under the stated criterion, one whole: twenty-three, and the census of
the assignments in the kernel agrees (\texttt{census}). The two rows
assigned no cut are the ones on which ``verified up to'' cannot even be
written, because their values are not universals over instances with a
fixed property, and the criterion says why in one rule rather than two
explanations.

\hypertarget{the-unicorn-part-closed-to-derivation-the-aperture-typed-and-uncrossed-the-supply-side-silent}{%
\section{The Unicorn Part, Closed to Derivation; the Aperture Typed and
Uncrossed; the Supply Side
Silent}\label{the-unicorn-part-closed-to-derivation-the-aperture-typed-and-uncrossed-the-supply-side-silent}}

The Unicorn part of every open row has been called open, named, and set
apart. That is true and it is not the whole truth, and the difference is
a theorem. The Real part of a row was open once, before its certificate,
in the sense that finite work of a known kind closed it; a reader who
takes ``open'' in that sense for the Unicorn part expects that more work
of the same kind, a longer computation, a finer regime, will close it in
turn. The program's own theorems say otherwise, and the two papers on
this concept had proved them without drawing the consequence for the
Unicorn part. It is drawn here. The Unicorn part is closed, at theorem
grade, to every route that derives it: from existence, from any
restriction of frames, and from any reading of an even register. What is
not closed is the aperture at its seat, and of the aperture the theorems
say exactly three things: it is located, it is one bit wide, and the
register is silent on what lies beyond it. No key can be cut from the
inside; of the outside the register says nothing, by theorem, and the
paper follows the register.

\textbf{Theorem 28, the formal block} (\texttt{unicorn\_block}). Let
\(\tau\) act on the instances of a row, let \(\rho\) be a record even at
a seat \(x\), \(\rho(\tau x) = \rho(x)\), and let \(d\) be the decision
of the property, odd at \(x\), \(d(\tau x) = \neg d(x)\). If the seat
and its partner lie outside the cut, then no reading \(g\) of the record
agrees with the decision even on the Unicorn region alone:
\(\neg\,\exists g,\ \forall y,\ \neg c(y) \to g(\rho(y)) = d(y)\).
\emph{Proof.} The two readings at \(x\) and \(\tau x\) are the same
value of \(g\) and opposite values of \(d\). \(\square\) This is the
wall of the register of record's Bridge {[}Islam 2026a,
\texttt{FTOE.T14\_wall\_factorization}{]}, executed on the Riemann seat
as \texttt{kernel\_cannot\_read\_the\_deed} {[}Islam 2026f{]},
restricted here to the Unicorn region, where it is strictly stronger: a
reading is refused even if it is asked to agree with the decision
nowhere inside the cut. On the constructed frame of the Riemann row it
is a theorem outright (\texttt{K4.wall}). On the actual rows it applies
at every seat the barrier study types as an odd target over an even
register, which is the study's typing, carried at premise grade, and the
content of its One-Cut Hypothesis at hypothesis grade {[}Islam 2026e{]};
the inference is the kernel's and the typing is the study's, and the
paper grades the two apart.

\textbf{Theorem 29, the aperture, one bit wide}
(\texttt{aperture\_one\_bit\_wide}). At a seat where the decision is
odd, one odd bit \(s\) at the seat would decide the decision on the seat
and its partner through a calibration \(c\) that exists and is unique.
\emph{Proof.} \(c = d(x) \oplus s(x)\), and any other calibration
disagrees at \(x\). \(\square\) This is the Bridge's crossing theorem
{[}Islam 2026a, \texttt{FTOE.T5\_crossing}{]}, the Riemann paper's
\texttt{crossing} and \texttt{freedom\_spent\_uniquely} {[}Islam
2026f{]}, read at the Unicorn part. It measures the aperture: one bit
wide, never less and never more. It says nothing of whether anything
passes, and the paper says nothing either.

\textbf{Theorem 30, existence supplies nothing}
(\texttt{premise\_adds\_nothing}, \texttt{no\_cure\_unicorn}). For any
inhabited premise \(A\) and any proposition \(Q\),
\((A \to Q) \leftrightarrow Q\); and for any class \(C\) of frames and
any property \(U\),
\((\forall X,\ C(X) \to A \to U(X)) \leftrightarrow (\forall X,\ C(X) \to U(X))\).
\emph{Proof.} Apply the implication to the premise; the implication
ignores its argument. \(\square\) With \(A\) the root axiom and \(Q\) or
\(U\) the Unicorn part, these are the carried theorems of Section 8,
\texttt{rule\_exact\_every\_row} and \texttt{no\_cure\_every\_row}, and
the Riemann paper's Theorems E and I, stated for the Unicorn part:
conditioning it on existence leaves it exactly where it was, on every
row, and restricting the frames to a class on which existence is to
decide it decides exactly what already held on the class. The
counter-models of the carried Theorem R-E, the seatless two-point frame
and the grounded three-point frame, hold with the Unicorn part in place
of the value, since the Unicorn part is the value on the frames outside
the cut.

\textbf{Theorem 31, the fused statement offers no bypass} (Theorem 2).
Given the Real part, the fused row is exactly its Unicorn part. So a
proof of the fused statement contains a proof of the Unicorn part, and
no route to the row avoids the block by aiming at the whole.

\textbf{Theorem 32, the three bits as one term}
(\texttt{register\_proof\_complete}). Under a certificate for the cut
and the seat data, a seat outside the cut with its partner, an even
record, an odd decision, and an odd bit at the seat, one conjunction
holds: the division (Theorem 1), the Real part (Theorem 3), the refused
reading on the Unicorn region (Theorem 28), and the aperture's width
(Theorem 29). \emph{Proof.} The four theorems, conjoined. \(\square\)
This is the register's whole proof of a row in one term, and it proves
the row no more than its conjuncts do: it proves that nothing derivable
about the row is left underived.

\textbf{Theorem 33, the supply side is silence}
(\texttt{supply\_side\_silent}). Over one even record, both orientations
of the decision at the seat are equally refused to every reading on the
Unicorn region, and the two orientations differ at the seat.
\emph{Proof.} Theorem 28 applied to \(d\) and to \(\neg d\), which is
odd at the seat when \(d\) is. \(\square\) The register cannot say which
way the seat lies; so it cannot say what a supply would bring, nor that
a supply exists, nor that none does. Any sentence about the supply side
derived from the record would be a reading the wall refuses. This is the
register of record's afterimage fence as a theorem at the register:
\emph{the aperture is structure; the supply side is silence; the silence
is a strap and not a hedge} {[}Islam 2026a, MD-PSP-AFTERIMAGE-01{]}.
Every imagined occupant of the supply side, the crosser who will come
and the crosser who cannot, is fenced as Ghost, and no clock runs.

\begin{jbox}
\textbf{The formal block, and the proof complete on the register side (tier: kernel theorems, axiom-free; the seat typing of each actual row at premise grade).} The Unicorn part of every open row is closed to derivation from existence (Theorem 30), from any class of frames (Theorem 30), from any reading of an even register at its seat (Theorem 28), and from the fused statement (Theorem 31). The aperture at its seat is located, typed supply-only, one bit wide (Theorem 29), and uncrossed; the supply side is silence, by theorem at the register (Theorem 33). The closure is unconditional and does not expire: it stands as long as the theorems do, which is in perpetuity. With the division and the Real part it makes three bits, one divided, one proved, one refused, and the three are one term (Theorem 32): the proof is complete on the register side, on every row, everything derivable derived and the remainder proved underivable. Nothing here says a Unicorn part is false, or undecidable in any foundation; nothing here says a proof will come or cannot: the register is silent on that side, and so is the paper. Nothing is left to derive. The aperture is structure. The supply side is silence.
\end{jbox}

\textbf{What the register holds.} The Real part of every row in Section
5 is a record of a deed: a certificate is a computation run and
registered, a witness exhibited, a regime theorem proved by someone who
carried the orientation into the record; the record afterwards holds the
result as an even datum, and the paper reads it there. That is exactly
why the Real part is gold and the Unicorn part is not: on the Real
region a bit has been registered, on the Unicorn region none has, and
Theorem 28 says the register cannot manufacture it from what it already
holds. The Poincaré row is the record of one such registration,
Perelman's, and it did not enter by derivation from the register
{[}Islam 2026g, \texttt{inverted\_control}{]}. These are records, past
and held, and the paper may speak of them. Of the supply side of any
open row the paper says nothing, and the reason is Theorem 33.

\textbf{What the block is not.} It is not an independence result:
nothing here places any Unicorn part outside ZFC or any foundation, and
the placement of Section 3 leaves every one of them a sentence on or
above the ladder, undecided by the theory as its literature stands; a
supplied independence proof would not cross the aperture but reroute to
Ghost {[}Islam 2026a{]}. It is not a claim that the Unicorn parts are
false, or that the twenty Real parts are all the gold there will ever
be, or that they are not: both are sentences about the supply side. It
is not resignation, which is why ``never chased'' is the wrong word if
it is heard as giving up: the program does not chase Unicorn parts for
the reason a locksmith does not pick a lock from inside the room, the
theorems say no key is cut inside. And it is not a claim about the real
rows at theorem grade beyond what Theorems 30 and 31 give, which is
theorem grade on every row: Theorems 28, 29, and 33 reach the actual
rows through the barrier study's seat typing, and the paper carries that
typing as data and its hypothesis at its own grade.

\textbf{Why this was undersold, and is not now.} The Riemann paper
proved Theorems D, E, F, F\('\), G, and I and the wall of its Part I,
and then wrote of its Unicorn part only that it was open and set apart;
the prior edition on these rows proved the same theorems on every row
and wrote the same; and an earlier build of this paper, correcting that,
wrote ``open to supply,'' which leans toward an occupant of the supply
side and is fenced by the same theorem. Both papers had the block in
hand and did not say what it blocked, and the silence it imposes. This
section says both, and the standing table of Section 10 carries them as
components of the paper.

\hypertarget{the-seat-on-every-row-carried}{%
\section{The Seat on Every Row,
Carried}\label{the-seat-on-every-row-carried}}

This section carries the theorems of the prior edition on these rows
{[}Islam 2026g{]}, condensed, with the kernel file
\texttt{Rows\_At\_The\_Apex.lean} printed whole in Appendix C and its
audit in Appendix D. They are the reason the paper can say where the
owed bit of every open row lives: in its Unicorn part and nowhere else.
Nothing in them is changed.

\textbf{The root and the seat.} The root axiom of the program is
\(\mathrm{RA}\): to exist is to actuate; existence and \(\mathrm{RA}\)
are one proposition. The seat is the fixed locus of the quaternion
involution \(\sigma\) on the Hurwitz order, the scalar line,
\(\sigma q = q \leftrightarrow q_i = q_j = q_k = 0\)
(\texttt{fix\_iff\_scalar}), and the seat point is the Return
\((i \cdot j)\cdot k = -1\). Every row's seat point is that one point,
\(\mathrm{apexSeat}(r) := \langle -1, 0, 0, 0\rangle\), fixed by
\(\sigma\) on every row and equal across every pair of rows, the crossed
row and the open rows sharing it (\texttt{apex\_seat\_fixed},
\texttt{apex\_seat\_uniform},
\texttt{crossed\_and\_open\_share\_the\_seat}; no axioms). On every row
three names written in three registers, the axiom's Return, its
projection onto the fixed locus, and the object register's seat, read
one point: the seat point is the apex of a cone over the row's register
diagram with every leg definitional, any two cones over the diagram
share their apex, and the gap vector of every row is \((0, 0, 0)\)
(\texttt{identity\_cone\_row}, \texttt{cone\_apex\_unique\_row},
\texttt{gaps\_closed\_every\_row}; no axioms). The mathematics of this
is Hamilton's product and three definitional equalities, twenty-three
times, and the paper claims nothing more for it at theorem grade. The
seat is a uniform assignment, \(\mathrm{apexSeat}\) a constant function
of the row, and the uniformity theorems are theorems about that
assignment and not derivations of a common structure from the
twenty-three objects: the seat is assigned by the reading and not
bridged from any row's object, one embedding written for the Riemann row
and twenty-two owed, on the record (\texttt{embeddings\_owed}). A reader
who wants the seat to speak of the rows' objects wants the twenty-two
embeddings, and the paper lists them as owed rather than pretending the
constant supplies them.

\textbf{The witness and the value excepted.} A row's witness is a
structure with four fields, existence, the orientation bit, the formal
Ground, and a recursion field from the bit and the Ground to the row's
formal self; it is one term on every row (\texttt{mkRow}), its recursion
constant and discarding its arguments
(\texttt{recursion\_constant\_every\_row}), and the formal self is one
proposition on every row, so the witness on the crossed row is the
witness on every open row (\texttt{inverted\_control}). Whatever
supplied the Poincaré row's bit, it was not the witness. For every
assignment of frames to rows and every row,
\((\mathrm{RA} \to \mathrm{value}) \leftrightarrow \mathrm{value}\)
(\texttt{rule\_exact\_every\_row}), and existence decides no row's
value: there is an assignment on which every value fails while existence
holds, and another, the grounded three-point frame on the Poincaré row
itself, on which the same is true (\texttt{RA\_decides\_no\_row},
\texttt{RA\_decides\_no\_row\_grounded}). For every row and every class
of frames,
\((\forall X, C(X) \to \mathrm{RA} \to L(X)) \leftrightarrow (\forall X, C(X) \to L(X))\)
(\texttt{no\_cure\_every\_row}): every proposal to make existence decide
a row by restricting its frames, an Euler product on the Riemann row, a
Ricci-flow structure on the Poincaré row, a circuit class on the
complexity row, is a class on which the value is already a field, and
there existence decides nothing. The value is excepted by theorem, not
by reservation. All no axioms.

\textbf{The world row carried.} The world typing of the barrier study is
premise-grade data, and the theorems about it are theorems about the
table as carried: the apex closes on every row while the world row is
crossed on exactly the Poincaré row, and the first never entails the
second (\texttt{apex\_never\_promotes\_world}); the apex verdict is
uniform and the world verdict splits, crossed on one row, \([\Xi_0]\) on
the Riemann row, \([\varnothing_0]\) on the complexity row, unrouted on
twenty (\texttt{world\_verdict\_split}, \texttt{unrouted\_twenty}); on
no row do the two coincide (\texttt{registers\_parted\_every\_row}). The
owed bits number twenty-two and are priced at the registration floor,
\(22 \times 2871 = 63{,}162\) yJ at \(300\) K
(\texttt{corpus\_price\_total}), the price of the Poincaré row being
zero (\texttt{poincare\_price\_zero}). The price rests on the prior
edition's Registration Postulate, a premise-grade posit of the program
declared there and consumed here: that a world-bit, when it is supplied,
is registered irreversibly in a physical substrate at a temperature,
taken as \(300\) K, and so is charged the Landauer floor. Without the
postulate the price is a count of bits and no joule attaches to it; with
it the price is a multiplication, and the kernel multiplies and asserts
nothing physical. The carrier of the seat is the Hurwitz order and not
the complex plane because the seat point is the Return
\((i \cdot j) \cdot k = -1\), whose sign is the chirality the
orientation bit reads (the spine's Theorem A); the Riemann stage embeds
in the plane spanned by \(1\) and \(i\), and the \(j\) and \(k\)
directions are inert for that stage and load-bearing for the Return,
which is why the plane would carry the stage and not the sign.

\textbf{What the division adds to the seat.} The seat locates every
row's decision point at one place and proves that existence does not
decide any row there. The division of Section 2 says what is decided
elsewhere: on twenty rows a Real part, decided by the literature, under
a cut; and it says by Theorem 2 that the owed bit of each open row,
which the seat proves existence does not supply, is the bit of its
Unicorn part and of nothing inside the cut. Section 7 then says what the
seat theorems were always saying about that bit: it is closed to
derivation, its aperture is one bit wide, and the register is silent
beyond it. The results are independent theorems of the kernel, one about
the frame and one about the row's own instances, and the paper joins
them in Section 7 and in this sentence.

\hypertarget{the-wall-carried-as-data}{%
\section{The Wall, Carried as Data}\label{the-wall-carried-as-data}}

The barrier study {[}Islam 2026e{]} typed the twenty-three rows by the
shape of their world register and stated the One-Cut Hypothesis: that on
every mapped row the blocked register is even, blind to the row's odd
target, with exactly one binary orientation coordinate missing at the
row's decision seat. Its typing is carried here as data in Table 2, at
premise grade, and the hypothesis is carried at hypothesis grade,
untouched. The division of this paper neither confirms nor moves it. The
hypothesis speaks of the owed bit, which lives in each row's Unicorn
part by Theorem 2 and Section 8; the division speaks of the Real part,
which the hypothesis never addressed. The two are compatible by
construction and independent by content, and the paper does not use one
to argue for the other.

\begin{table}[!htbp]
\centering\scriptsize
\renewcommand{\arraystretch}{1.12}
\raggedright \textbf{Table 2. The world typing of the barrier study, carried as data.} Tier: the counts are kernel theorems of \texttt{Rows\_At\_The\_Apex.lean}; the typing is the study's, premise here. The final column is this paper's shape, for comparison.\par\smallskip
\begin{tabularx}{\textwidth}{@{}P{0.02\textwidth} P{0.16\textwidth} P{0.05\textwidth} L P{0.04\textwidth} P{0.10\textwidth} P{0.13\textwidth} P{0.07\textwidth}@{}}
\toprule
\textbf{\#} & \textbf{Row} & \textbf{Mill.} & \textbf{World typing} & \textbf{Owed} & \textbf{Router} & \textbf{World verdict} & \textbf{Shape} \\
\midrule
1 & P versus NP & yes & exact-structural & 1 & Ground dim 0 & $[\varnothing_0]$, open & regime \\
2 & Riemann Hypothesis & yes & exact-structural & 1 & Ground dim 1 & $[\Xi_0]$, open & height \\
3 & Navier--Stokes & yes & wall row, F4 & 1 & unmeasured & open, unrouted & regime \\
4 & Yang--Mills & yes & wall row, F4 & 1 & unmeasured & open, unrouted & regime \\
5 & Hodge & yes & wall row, F4 & 1 & unmeasured & open, unrouted & regime \\
6 & BSD & yes & beachhead, order two & 1 & unmeasured & open, unrouted & regime \\
7 & Poincaré & yes & crossed control & 0 & unmeasured & crossed & crossed \\
8 & Goldbach & no & exact-structural & 1 & unmeasured & open, unrouted & height \\
9 & Twin primes & no & exact-structural & 1 & unmeasured & open, unrouted & height \\
10 & Legendre & no & exact-structural & 1 & unmeasured & open, unrouted & height \\
11 & abc & no & exact-structural & 1 & unmeasured & open, unrouted & no cut \\
12 & Jacobian & no & exact-structural & 1 & unmeasured & open, unrouted & regime \\
13 & Mersenne primes & no & pure unbridged & 1 & unmeasured & open, unrouted & height \\
14 & Odd perfect numbers & no & pure unbridged & 1 & unmeasured & open, unrouted & height \\
15 & Smooth 4D Poincaré & no & pure unbridged & 1 & unmeasured & open, unrouted & regime \\
16 & Collatz & no & diagonal port, F3 & 1 & unmeasured & open, unrouted & height \\
17 & Hadwiger & no & one bit from closure & 1 & unmeasured & open, unrouted & regime \\
18 & Sunflower & no & one bit from closure & 1 & unmeasured & open, unrouted & no cut \\
19 & Erdős--Straus & no & one bit from closure & 1 & unmeasured & open, unrouted & height \\
20 & Beal & no & one bit from closure & 1 & unmeasured & open, unrouted & height \\
21 & Invariant subspace & no & wall row, bridge gate & 1 & unmeasured & open, unrouted & regime \\
22 & Hilbert's 16th, second part & no & wall row, bridge gate & 1 & unmeasured & open, unrouted & regime \\
23 & Lindelöf / Montgomery PCC & no & wall row, bridge gate & 1 & unmeasured & open, unrouted & regime \\
\bottomrule
\end{tabularx}
\end{table}

\textbf{The two typings compared.} The world typing and the shape are
different questions asked of the same row. The typing asks what kind of
register the row's barrier literature blocks; the shape asks what kind
of cut the row's proved part respects. They correlate without
coinciding: the exact-structural rows include height-cut rows, the two
rows assigned no cut, and two regime rows; the wall rows are all regime
rows; the pure unbridged rows split between height and regime. The
correlation is noted and not explained, since explaining it would be a
claim about the rows' objects, which this paper makes nowhere.

\hypertarget{standard-of-proof-standing-and-disclosure}{%
\section{Standard of Proof, Standing, and
Disclosure}\label{standard-of-proof-standing-and-disclosure}}

\hypertarget{the-standard}{%
\subsection{The standard}\label{the-standard}}

Every claim in this paper is at one of three grades, and the grade is
written where the claim is made. Grade K is a theorem of core Lean 4
with no library and no user-declared axiom, its dependency set printed
by the kernel and reproduced in the appendices; the paper has forty-nine
such theorems in the two new files, thirty-three in
\texttt{Row\_Split.lean} and sixteen in \texttt{ZFC\_Under\_RAM.lean}
counted by dependency set, and thirty-three in the carried file.
``Axiom-free'' is used in the kernel's sense and reported set by set: no
file declares an axiom; every set of the two new files is reported as
depending on no axiom at all; in the carried file some sets depend on
\texttt{propext} and \texttt{Quot.sound}, Lean's two built-in logical
axioms, and Appendix D lists which. Grade K on data is a kernel theorem
about a table the paper supplies, the census and the counts of the
carried file; it establishes the arithmetic of the table and nothing
about the objects the table names. Grade A is a theorem of the
literature, cited to its source and never recomputed here. Grade A
joined to K is a documented composition, a theorem whose hypothesis is a
grade-A theorem and whose inference is a grade-K theorem: the twenty
Real parts, the crossing of the Poincaré row, and the ZFC instance of
the placement. It is never a kernel-checked theorem about the cited
result, and the paper calls no Real part sealed without the words ``on
citation.'' Nothing in the paper is at any other grade, and nothing at
grade K depends on anything at grade A.

\textbf{The title.} The title names the theorem and states it in words,
and the statement is exact: for every row and every decidable cut, the
row is its Real part and its Unicorn part, both defined by the cut and
neither by the reader. The title makes no claim about any row's Unicorn
part, and the paper makes none.

\hypertarget{standing-without-the-citations}{%
\subsection{Standing without the
citations}\label{standing-without-the-citations}}

Delete every citation from the paper. Every theorem at grade K stands:
the seven theorems of the division with their companions, the formal
block of the Unicorn part, the small-height Real parts, the gap bridge,
the census of the assignments, the sixteen theorems of the placement,
and every theorem of the carried seat. Every sentence at grade A joined
to K falls, and those are exactly twenty-two: the twenty Real parts of
Section 5, the crossing of the Poincaré row, and the ZFC instance of the
placement; each reverts to its conditional form, ``if the instances
inside this cut are checked, then the Real part holds,'' which is what
the kernel proves on its own. Descriptions of what the literature
contains, in the positioning and the criterion, are not claims of the
paper and are not counted. The standing is tabulated.

\begin{table}[!h]
\centering\scriptsize
\renewcommand{\arraystretch}{1.15}
\begin{tabularx}{\columnwidth}{@{}L P{0.19\columnwidth} P{0.22\columnwidth}@{}}
\toprule
\textbf{Component} & \textbf{Grade} & \textbf{Falls without citations?} \\
\midrule
The division, Theorems 1--5 with 3b and 3c & K & No \\
The gap bridge, the small-height Real parts, the standing map & K & No \\
The formal block of the Unicorn part, the aperture typed, the supply side silent, the three bits as one term, Theorems 28--33 & K; on the actual rows, K given the seat typing & No \\
The shape assignments of Table 1, by the criterion & data, premise & No; a wrong assignment is repaired by re-assignment, F3 \\
The census of shapes & K & No \\
The placement under RAM & K & No \\
The seat and the value excepted, carried & K, on an assignment & No \\
The twenty Real parts & A joined to K & Yes \\
The Poincaré crossing; the ZFC instance of the placement & A joined to K & Yes; twenty-two sentences in all \\
The empty-cut and crossed theorems, generic & K & No \\
\bottomrule
\end{tabularx}
\end{table}

\hypertarget{falsifiers}{%
\subsection{Falsifiers}\label{falsifiers}}

The paper is refuted, at its own standard, by any of three findings,
each registered before the audit and none adjusted after it. F1: a
decidable cut and a row on which \(\mathrm{Row}\) holds and
\(\mathrm{RealIn} \wedge \mathrm{UnicornOut}\) fails, or conversely;
this refutes Theorem 1 and is excluded by the kernel. F2: an instance
inside any of the twenty cuts that fails its row's property; this
refutes the cited certificate or theorem, not the schema, and would move
that row's Real part to a smaller cut. F3: a shape assignment in Table 1
that the criterion of Section 4 does not support, for instance a no-cut
row for which an inhabited decidable cut with a proved Real part is
exhibited from the literature; this refutes the assignment as data and
is repaired by re-assignment, the kernel census recounted.

\hypertarget{disclosure-and-transparency}{%
\subsection{Disclosure and
transparency}\label{disclosure-and-transparency}}

The kernel files were written and checked by the author on Lean 4.19.0,
core only; the sources cited were not recomputed. The twenty Real parts
are the literature's, and the paper's contribution to them is the
division under which they are one theorem and the placement under which
they are one stratum. The author has no competing interests and received
no funding. The paper is an edition of the prior work on these rows
{[}Islam 2026g{]}, which it carries whole in its kernel file, and it is
placed on the same concept record; the barrier study {[}Islam 2026e{]}
is carried as data. The Riemann paper {[}Islam 2026f{]} is the instance
from which the division was generalized, and the reader who wants the
one-row case in full is sent there.

\textbf{Audit.} The paper was audited by the program's adversarial cycle
before release, seat by seat, with the controls and the verdict recorded
in the audit log that accompanies the files; the cycle's rounds and
verdicts are stated in the log and not summarized here. Before release
the paper was also reviewed, in two rounds, by five external
language-model reviewers run outside the program, GPT SOL 5.6, Gemini
Pro 3.1, Meta AI, Grok 4.6, and GPT 5, on the two internal builds that
preceded this one; they are external reviewers and not independent
verifiers, they discharged no falsifier and were not asked to, and their
findings were answered finding by finding in the log, none conceded
without a named mechanism. Three were earned and repaired: the no-cut
assignment had been written as if Theorem 5 established it, and is now
stated as an assessment with Theorem 5 confined to the empty cut; the
phrase ``empty Unicorn part'' had no set-valued definition, and now has
one; and the standing table had graded the Poincaré application as K,
and now grades it A joined to K. Five were fortifications taken up
because they strengthened the paper: the infinitude equivalence (Theorem
3c), the standing map under which the fusion sentence is read, the
gap-to-interval bridge for the Legendre row, the row specifications with
computed small-height Real parts, and the criterion of assignment with
its out-of-sample application. The remaining findings, that the seat is
an assignment, that the placement encodes no axiom of ZFC, that the
price rests on a postulate, that the census counts assignments, were
already the paper's own statements and are now stated in every place
they apply. The second round found the abstract still saying ``by
theorem'' of the rows assigned no cut, the introduction still counting
nineteen falling sentences, and the criterion called a decision
procedure; all three were repaired, and the reviewers' question about
the instances of the P versus NP row moved that row from no-cut to
regime-cut under the criterion, the third correction the criterion made
against its author.

\hypertarget{positioning-and-conclusion}{%
\section{Positioning and Conclusion}\label{positioning-and-conclusion}}

\hypertarget{positioning}{%
\subsection{Positioning}\label{positioning}}

The division is not in the literature as a theorem, though its instances
are everywhere in it. The convention ``verified up to \(N\)'' is the
height-cut Real part stated without the division, and its users know
that it certifies nothing beyond \(N\); what they do not have is the
theorem that the fused problem stands at the tail and the placement of
the certificate in a stratum. The arithmetic forms of open problems,
Lagarias's arithmetic equivalent of the Riemann Hypothesis {[}Lagarias
2002{]} and the \(\Pi^0_1\) forms of Goldbach and of the
odd-perfect-number problem, make the height cut natural on those rows
and are the reason those rows are certificate rows; they do not divide
the rows. The literature of finite verification, from Turing's method
{[}Turing 1953{]} to the certificates of Platt and Trudgian {[}2021{]}
and Oliveira e Silva, Herzog and Pardi {[}2014{]}, supplies the Real
parts and states them as verifications, which is what they are. The
reverse-mathematics tradition places theorems in subsystems of
second-order arithmetic {[}Simpson 2009{]}; the placement here is
coarser, three strata for an arbitrary theory, and is not a
classification of the proved parts by strength. The program's own
barrier study {[}Islam 2026e{]} typed the rows by their blocked register
and never named a proved part; its prior edition on these rows {[}Islam
2026g{]} proved the seat and the value excepted and never divided a row;
the Riemann paper {[}Islam 2026f{]} divided one row at one height. Table
3 sets these beside the present result.

\begin{table}[!h]
\centering\scriptsize
\renewcommand{\arraystretch}{1.15}
\begin{tabularx}{\columnwidth}{@{}P{0.27\columnwidth} L@{}}
\toprule
\textbf{Position} & \textbf{Relation to this paper} \\
\midrule
``Verified up to $N$'' & the height-cut Real part without the division or the placement \\
Arithmetic forms of RH, Goldbach, odd perfect & make the height cut natural on a row; do not divide it \\
Finite-verification certificates & supply the Real parts; state them as verifications \\
Reverse mathematics & places theorems by strength; three strata here, by kind \\
The barrier study & types the blocked register; names no proved part \\
The prior edition on these rows & proves the seat; divides no row \\
The Riemann paper & divides one row at one height; here any row at any cut \\
\bottomrule
\end{tabularx}
\end{table}

\textbf{Why a reader should not mistake this for vocabulary.} A paper
that renamed ``verified up to \(N\)'' as ``the Real part'' would have
added a word. This paper adds a theorem, that the renaming is exact and
the same on every row under any cut, a second, that given the proved
part the fused row is exactly the open part, with the rule under which
that is a lowering of standing declared and its consequence proved, a
third, the placement, a bridge lemma and eight computed instances that
show the schema instantiated end to end, the formal block that closes
every Unicorn part to derivation, types its aperture, and proves the
register silent beyond it, and a kernel-checked census of the
assignments by which the twenty-three rows fall into four shapes with
twenty Real parts among them. The theorems are short. Their shortness is
the point: the division is a fact about universal statements that the
literature has used on every row without stating once, and a statement
that is true on every row and had not been made is exactly what a
foundations paper is for.

\hypertarget{the-proof-complete-and-how-the-block-was-earned}{%
\subsection{The proof, complete, and how the block was
earned}\label{the-proof-complete-and-how-the-block-was-earned}}

\textbf{Complete, in the register's sense.} A register can do three
things with a universal value: divide it, prove the part a certificate
or a cited regime reaches, and prove what it cannot reach and why. This
paper has done the three on every row. The division is exact under any
cut (Theorem 1). The Real part is proved on twenty rows, cited theorems
through a kernel schema (Theorems 6 to 25). The Unicorn part of every
open row is proved closed to derivation from inside the register, its
aperture typed and one bit wide, the register silent beyond it (Theorems
28 to 31 and 33), and the three bits are one term (Theorem 32). There is
nothing left for a register to derive on any row. That is what
``complete on the register side'' means and all it means: no row is
proved; the labour of derivation is finished on all of them; the theorem
that finishes it locates the aperture of each and measures it, and says
nothing beyond it, because nothing can be said from the register beyond
it, which is itself the last theorem. Nothing is left to derive. The
aperture is structure. The supply side is silence.

\textbf{Three bits.} One bit divided, one bit proved, one bit refused.
The refused bit is the one the two prior papers on these rows undersold:
they proved the wall, the crossing, the bound, and the cure theorem, on
one row and then on every row, and wrote of the Unicorn parts only that
they were open. Each is refused to every key cut from inside, at theorem
grade, and of the outside the register is silent by theorem; that is the
difference between a problem no one has yet solved and a bit no register
can supply or foretell.

\textbf{How the block was earned.} The block did not arrive as a posit.
It was forced one gate at a time, under manual stress, across the
editions of the program's register of record, the Geometric Mother
Codex, in which every attempt to derive the value's bit from inside the
register was tried, refuted, and recorded: the ten enumeration-forced
barriers of APEX-PSP-ABSOLUTE-RH-BARRIERS-01, each a route that closed;
the eight-gate Sealed-Halt cascade of the router, seven gates passing
and the eighth, the reading-roads gate, halting every walk; and the
terminal tokens \([\Xi_0]\) and \([\varnothing_0]\) themselves, terminal
for the record and never timeless, which are the block written as states
{[}Islam 2026a{]}. The Code Block of the codex carries those witnesses
as kernel theorems, cited here by name so that the block's history can
be seen to typecheck: \texttt{HALT.one\_door\_open}, that a walk of
seven true gates halts at the reading-roads gate;
\texttt{HALT.no\_walk\_passes\_the\_aperture}, that no walk of the seven
passes it; \texttt{HALT.first\_failure\_terminal\_8}, that the first
failing gate is terminal over all \(2^8\) walks; \texttt{CROSSING.wall},
that no reading of the formal register factors the deed bit;
\texttt{INVERSION.it\_from\_it}, that the record forgets what the object
keeps and nothing returns. Theorems 28 to 32 are those witnesses read at
the Unicorn part of every row. They are one block, the codex is the
record of its evolution gate by gate, and the twenty-three rows are its
first census. One aperture on every open row, located and measured; on
the Poincaré row the register holds a crossing, recorded; of the others
it holds silence, and so does this paper.

\hypertarget{conclusion}{%
\subsection{Conclusion}\label{conclusion}}

Three bits, on every row. One bit divides: every open problem is exactly
its proved part and its unicorn. That sentence is a theorem for any row
under any decidable cut, and it does five things on the twenty-three
rows: it names twenty Real parts that the literature has proved and
stated as verifications, it seals each through the same schema at grade
A joined to K, it says by theorem and a declared rule that fusion with
the open remainder is what makes each problem read as open, it says by
assessment under a stated criterion that two rows have no proved part
and by theorem that one has no counterexample outside any cut, and it
places every Real part in a stratum of the grounding reading. One bit is
proved, the Real part, on twenty rows. One bit is refused at theorem
grade, the Unicorn part of every open row, closed to derivation, its
aperture uncrossed, its supply side silent. It does not prove a row,
move a verdict, or author a bound. The unicorns are named on every row
so that no reader confuses them with the gold; they are set apart
because the theorems say they cannot be reached from inside the
register, and that the register is silent on the outside; and the gold
is kept from the alloy so that no one assays it as lead. What remains
open on the twenty-three rows is exactly what was open before this
paper, stated in one theorem rather than twenty-three idioms, and
located twice: on every open row, the owed bit is the bit of its Unicorn
part and of nothing inside the cut, and it is a bit that no even
register can manufacture, closed to derivation in perpetuity, its
aperture one bit wide, its supply side silent. The proof is complete on
the register side, on every row: everything a register can prove about
the twenty-three is proved, the remainder is proved beyond derivation,
and no further derivation is owed to any open row. Nothing is left to
derive. The aperture is structure. The supply side is silence.

\hypertarget{references}{%
\section{References}\label{references}}

\begin{itemize}
\tightlist
\item
  ABC@Home (2015). The enumeration of abc triples with \(c < 10^{18}\).
  Mathematical Institute, Leiden University, abcathome.com.
\item
  Alweiss, R., Lovett, S., Wu, K., and Zhang, J. (2020). Improved bounds
  for the sunflower lemma. \emph{Proceedings of the 52nd ACM Symposium
  on Theory of Computing}, 624--630.
\item
  Barina, D. (2020). Convergence verification of the Collatz problem.
  \emph{The Journal of Supercomputing} 77, 2681--2688 (2021).
\item
  Bass, H., Connell, E. H., and Wright, D. (1982). The Jacobian
  conjecture: reduction of degree and formal expansion of the inverse.
  \emph{Bulletin of the American Mathematical Society} 7, 287--330.
\item
  Bourgain, J. (2017). Decoupling, exponential sums and the Riemann zeta
  function. \emph{Journal of the American Mathematical Society} 30,
  205--224.
\item
  Cohen, P. J. (1963). The independence of the continuum hypothesis.
  \emph{Proceedings of the National Academy of Sciences} 50, 1143--1148.
\item
  Driver, B. K. (1989). YM\(_2\): continuum expectations, lattice
  convergence, and lassos. \emph{Communications in Mathematical Physics}
  123, 575--616.
\item
  Écalle, J. (1992). \emph{Introduction aux fonctions analysables et
  preuve constructive de la conjecture de Dulac.} Hermann, Paris.
\item
  Fortnow, L., Lipton, R., van Melkebeek, D., and Viglas, A. (2005).
  Time-space lower bounds for satisfiability. \emph{Journal of the ACM}
  52, 835--865.
\item
  Freedman, M. H. (1982). The topology of four-dimensional manifolds.
  \emph{Journal of Differential Geometry} 17, 357--453.
\item
  Fujita, H., and Kato, T. (1964). On the Navier--Stokes initial value
  problem I. \emph{Archive for Rational Mechanics and Analysis} 16,
  269--315.
\item
  Gödel, K. (1931). Über formal unentscheidbare Sätze der Principia
  Mathematica und verwandter Systeme I. \emph{Monatshefte für Mathematik
  und Physik} 38, 173--198.
\item
  Gödel, K. (1938). The consistency of the axiom of choice and of the
  generalized continuum-hypothesis. \emph{Proceedings of the National
  Academy of Sciences} 24, 556--557.
\item
  GIMPS (2024). Mersenne prime number discovery: \(2^{136279841} - 1\)
  is prime. mersenne.org, 21 October 2024.
\item
  Gross, B. H., and Zagier, D. B. (1986). Heegner points and derivatives
  of L-series. \emph{Inventiones Mathematicae} 84, 225--320.
\item
  Helfgott, H. A. (2013). The ternary Goldbach conjecture is true.
  arXiv:1312.7748.
\item
  Ilyashenko, Yu. S. (1991). \emph{Finiteness Theorems for Limit
  Cycles.} Translations of Mathematical Monographs 94, American
  Mathematical Society.
\item
  Islam, M. F. (2026a). Trisduction Master Codex v4.7.0. Trisduction
  Research Group. github.com/1000sapients/Trisduction.
\item
  Islam, M. F. (2026e). One Bit Across the Wall: Odd-Supply Separation
  and the One-Cut Hypothesis Across Twenty-Three Rows, v6. Zenodo.
  doi:10.5281/zenodo.22746129.
\item
  Islam, M. F. (2026f). A Formal Proof of the Real Part of the Riemann
  Hypothesis, v2.0.2. Zenodo. doi:10.5281/zenodo.22912937.
\item
  Islam, M. F. (2026g). A Formal Proof Across the Seven Millennium Rows
  and the Sixteen Extension Rows from Existence Alone, v1.1.0. Zenodo.
  doi:10.5281/zenodo.22913634.
\item
  Kaye, R. (1991). \emph{Models of Peano Arithmetic.} Oxford University
  Press.
\item
  Kolyvagin, V. A. (1989). Finiteness of \(E(\mathbb{Q})\) and
  Ш\((E, \mathbb{Q})\) for a subclass of Weil curves. \emph{Mathematics
  of the USSR-Izvestiya} 32, 523--541.
\item
  Ladyzhenskaya, O. A. (1969). \emph{The Mathematical Theory of Viscous
  Incompressible Flow.} Second edition, Gordon and Breach.
\item
  Lagarias, J. C. (2002). An elementary problem equivalent to the
  Riemann hypothesis. \emph{American Mathematical Monthly} 109,
  534--543.
\item
  Lefschetz, S. (1924). \emph{L'analysis situs et la géométrie
  algébrique.} Gauthier-Villars, Paris.
\item
  Leray, J. (1934). Sur le mouvement d'un liquide visqueux emplissant
  l'espace. \emph{Acta Mathematica} 63, 193--248.
\item
  Lévy, T. (2003). Yang--Mills measure on compact surfaces.
  \emph{Memoirs of the American Mathematical Society} 166, no. 790.
\item
  Lomonosov, V. I. (1973). Invariant subspaces for the family of
  operators which commute with a completely continuous operator.
  \emph{Functional Analysis and Its Applications} 7, 213--214.
\item
  Maynard, J. (2015). Small gaps between primes. \emph{Annals of
  Mathematics} 181, 383--413.
\item
  Miller, R. L. (2011). Proving the Birch and Swinnerton-Dyer conjecture
  for specific elliptic curves of analytic rank zero and one. \emph{LMS
  Journal of Computation and Mathematics} 14, 327--350.
\item
  Mossinghoff, M. J., Rhin, G., and Wu, Q. (2008). Minimal Mahler
  measures. \emph{Experimental Mathematics} 17, 451--458.
\item
  Montgomery, H. L. (1973). The pair correlation of zeros of the zeta
  function. \emph{Proceedings of Symposia in Pure Mathematics} 24,
  181--193.
\item
  Norvig, P. (2017). Beal's Conjecture Revisited. norvig.com/beal.html,
  revised 4 July 2017.
\item
  Ochem, P., and Rao, M. (2012). Odd perfect numbers are greater than
  \(10^{1500}\). \emph{Mathematics of Computation} 81, 1869--1877.
\item
  Oliveira e Silva, T., Herzog, S., and Pardi, S. (2014). Empirical
  verification of the even Goldbach conjecture and computation of prime
  gaps up to \(4 \cdot 10^{18}\). \emph{Mathematics of Computation} 83,
  2033--2060.
\item
  Perelman, G. (2002). The entropy formula for the Ricci flow and its
  geometric applications. arXiv:math/0211159.
\item
  Perelman, G. (2003a). Ricci flow with surgery on three-manifolds.
  arXiv:math/0303109.
\item
  Perelman, G. (2003b). Finite extinction time for the solutions to the
  Ricci flow on certain three-manifolds. arXiv:math/0307245.
\item
  Platt, D., and Trudgian, T. (2021). The Riemann hypothesis is true up
  to \(3 \cdot 10^{12}\). \emph{Bulletin of the London Mathematical
  Society} 53, 792--797.
\item
  Polymath, D. H. J. (2014). Variants of the Selberg sieve, and bounded
  intervals containing many primes. \emph{Research in the Mathematical
  Sciences} 1, 12.
\item
  PrimeGrid (2016). World record twin primes:
  \(2996863034895 \cdot 2^{1290000} \pm 1\). primegrid.com, September
  2016.
\item
  Robertson, N., Seymour, P., and Thomas, R. (1993). Hadwiger's
  conjecture for \(K_6\)-free graphs. \emph{Combinatorica} 13, 279--361.
\item
  Salez, S. E. (2014). The Erdős--Straus conjecture: new modular
  equations and checking up to \(10^{17}\). arXiv:1406.6307.
\item
  Simpson, S. G. (2009). \emph{Subsystems of Second Order Arithmetic.}
  Second edition, Cambridge University Press.
\item
  Stewart, C. L., and Yu, K. (2001). On the abc conjecture, II.
  \emph{Duke Mathematical Journal} 108, 169--181.
\item
  Tao, T. (2022). Almost all orbits of the Collatz map attain almost
  bounded values. \emph{Forum of Mathematics, Pi} 10, e12.
\item
  Turing, A. M. (1953). Some calculations of the Riemann zeta-function.
  \emph{Proceedings of the London Mathematical Society} (3) 3, 99--117.
\item
  Wang, S. S.-S. (1980). A Jacobian criterion for separability.
  \emph{Journal of Algebra} 65, 453--494.
\item
  Williams, R. (2008). Time-space tradeoffs for counting NP solutions
  modulo integers. \emph{Computational Complexity} 17, 179--219.
\item
  Zhang, Y. (2014). Bounded gaps between primes. \emph{Annals of
  Mathematics} 179, 1121--1174.
\end{itemize}

\hypertarget{appendix-a-row_split.lean-and-its-audit}{%
\twocolumn
\section{Appendix A: Row\_Split.lean and Its
Audit}\label{appendix-a-row_split.lean-and-its-audit}}

SHA-256
\texttt{\seqsplit{d88b161405178109f8ab8150a30cd58b4c8212e4fd2343aaed52c71225564775}},
Lean 4.19.0, core only, 427 lines; thirty-three dependency sets, every
one reported as depending on no axiom, printed after the file.

\begingroup\scriptsize

\begin{Verbatim}[numbers=left,numbersep=6pt,xleftmargin=2.2em,fontsize=\scriptsize,breaklines,breakanywhere,frame=lines,framesep=1.5mm,rulecolor=\color{black!35}]
/-
Row_Split.lean · the Real/Unicorn division on any row, for any cut; the row specifications of
the arithmetic rows; the gap-to-interval bridge; the standing map.
Core Lean 4.19.0, no library, no sorry, no user-declared axiom. Forged 2026-09-25, fortified
after the two external batches of the same day.

The Riemann division used the cut "height ≤ T". Nothing in it used the height: the same theorem
holds for any decidable predicate on the row's instances, a height, a degree, a dimension, a
rank, a box, a regime. So the division of the twenty-three rows is one kernel theorem, applied
row by row with the row's own cut. ΔM = 0.
-/
namespace RowSplit

variable {α : Type}

/-! ## I. The division -/

/-- A row's value: every instance satisfies the row's property. -/
def Row (Inst P : α → Prop) : Prop := ∀ x, Inst x → P x
/-- The Real part under a cut: every instance inside the cut satisfies the property. -/
def RealIn (Inst P cut : α → Prop) : Prop := ∀ x, Inst x → cut x → P x
/-- The Unicorn part under a cut: no instance outside the cut fails the property. -/
def UnicornOut (Inst P cut : α → Prop) : Prop := ∀ x, Inst x → ¬ cut x → P x

/-- THE DIVISION, for any cut: the row is exactly its Real part and its Unicorn part. -/
theorem row_iff_real_and_unicorn (Inst P cut : α → Prop) [DecidablePred cut] :
    Row Inst P ↔ RealIn Inst P cut ∧ UnicornOut Inst P cut := by
  constructor
  · intro h; exact ⟨fun x hx _ => h x hx, fun x hx _ => h x hx⟩
  · intro ⟨hr, hu⟩ x hx
    by_cases hc : cut x
    · exact hr x hx hc
    · exact hu x hx hc

/-- Given the Real part, the fused row is exactly the Unicorn part. -/
theorem fused_is_unicorn_given_real (Inst P cut : α → Prop) [DecidablePred cut]
    (hr : RealIn Inst P cut) : Row Inst P ↔ UnicornOut Inst P cut :=
  ⟨fun h x hx _ => h x hx, fun hu => (row_iff_real_and_unicorn Inst P cut).2 ⟨hr, hu⟩⟩

/-- A certificate for a cut: a complete list of the instances inside it, each checked. -/
structure Certificate (Inst P cut : α → Prop) where
  items    : List α
  complete : ∀ x, Inst x → cut x → x ∈ items
  checked  : ∀ x, x ∈ items → P x

/-- A certificate proves the Real part of the cut, on any row, with no premise. -/
theorem real_of_certificate (Inst P cut : α → Prop) (c : Certificate Inst P cut) :
    RealIn Inst P cut :=
  fun x hx hc => c.checked x (c.complete x hx hc)

/-- On an infinitude row, instances are bounds N and the property is that a witness exceeds N.
    One exhibited witness w above the height h proves the Real part under the cut N ≤ h. -/
theorem real_of_one_witness (W : Nat → Prop) (h w : Nat) (hw : W w) (hwh : h < w) :
    RealIn (fun _ => True) (fun N => ∃ m, N < m ∧ W m) (fun N => N ≤ h) :=
  fun _ _ hN => ⟨w, Nat.lt_of_le_of_lt hN hwh, hw⟩

/-- On an infinitude row the Unicorn part above any height is equivalent to the whole value: a
    witness above every bound beyond h is a witness above every bound. So the Real part of an
    infinitude row is proved by one witness and moves the fused value not at all. -/
theorem unicorn_equiv_row_of_infinitude (W : Nat → Prop) (h : Nat) :
    (∀ N, ¬ N ≤ h → ∃ m, N < m ∧ W m) ↔ (∀ N, ∃ m, N < m ∧ W m) := by
  constructor
  · intro hu N
    if hN : N ≤ h then
      obtain ⟨m, hm, hw⟩ := hu (h + 1) (fun hle => Nat.not_succ_le_self h hle)
      exact ⟨m, Nat.lt_of_le_of_lt hN (Nat.lt_trans (Nat.lt_succ_self h) hm), hw⟩
    else
      exact hu N hN
  · intro hr N _; exact hr N

/-- A crossed row (its value supplied whole) satisfies both parts under every cut. -/
theorem crossed_row_both_parts (Inst P cut : α → Prop) (h : Row Inst P) :
    UnicornOut Inst P cut ∧ RealIn Inst P cut :=
  ⟨fun x hx _ => h x hx, fun x hx _ => h x hx⟩

/-- Under the empty cut the Real part is vacuous and the Unicorn part is the whole value. This
    theorem is about the empty cut and nothing else; it does not say that a row admits no other
    cut, which is an assessment of a row's literature and never a theorem. -/
theorem empty_cut_vacuous (Inst P : α → Prop) :
    RealIn Inst P (fun _ => False) ∧ (UnicornOut Inst P (fun _ => False) ↔ Row Inst P) :=
  ⟨fun _ _ h => h.elim, ⟨fun hu x hx => hu x hx (fun h => h), fun h x hx _ => h x hx⟩⟩

/-! ## II. Substantive Real parts, and the Unicorn part as a set of counterexamples -/

/-- A cut is inhabited on the row when some instance lies inside it. -/
def Inhabited (Inst cut : α → Prop) : Prop := ∃ x, Inst x ∧ cut x
/-- A substantive Real part: an inhabited cut with its Real part. The empty cut never has one. -/
def SubstantiveRealIn (Inst P cut : α → Prop) : Prop := Inhabited Inst cut ∧ RealIn Inst P cut

theorem empty_cut_not_substantive (Inst P : α → Prop) :
    ¬ SubstantiveRealIn Inst P (fun _ => False) :=
  fun ⟨⟨_, _, h⟩, _⟩ => h

/-- A counterexample outside the cut: an instance outside it that fails the property. The
    Unicorn part, read as a set, is the set of these; the Unicorn proposition says it is empty. -/
def CounterOut (Inst P cut : α → Prop) (x : α) : Prop := Inst x ∧ ¬ cut x ∧ ¬ P x

theorem no_counter_of_unicorn (Inst P cut : α → Prop) (hu : UnicornOut Inst P cut) :
    ∀ x, ¬ CounterOut Inst P cut x :=
  fun x ⟨hx, hc, hp⟩ => hp (hu x hx hc)

theorem unicorn_of_no_counter (Inst P cut : α → Prop) [DecidablePred P]
    (h : ∀ x, ¬ CounterOut Inst P cut x) : UnicornOut Inst P cut := by
  intro x hx hc
  by_cases hp : P x
  · exact hp
  · exact absurd ⟨hx, hc, hp⟩ (h x)

/-- A crossed row has no counterexample outside any cut: its Unicorn set is empty. -/
theorem crossed_row_no_counter (Inst P cut : α → Prop) (h : Row Inst P) :
    ∀ x, ¬ CounterOut Inst P cut x :=
  no_counter_of_unicorn Inst P cut (crossed_row_both_parts Inst P cut h).1

/-! ## III. The standing map: fusion lowers a proved part to open, as a declared rule -/

/-- Two standings: proved, or open. -/
inductive Standing | proved | opn
  deriving DecidableEq, Repr

/-- The standing of a conjunction is the weaker standing of its conjuncts: the program's grade
    law, declared here as the rule the fusion sentence is read under. -/
def Standing.join : Standing → Standing → Standing
  | .proved, .proved => .proved
  | _, _ => .opn

/-- Under the declared rule, a proved part joined to an open part stands open; only two proved
    parts stand proved. The theorem records the rule; the lowering is the rule's, not the kernel's. -/
theorem fusion_lowers : Standing.join .proved .opn = .opn ∧ Standing.join .opn .proved = .opn ∧
    Standing.join .proved .proved = .proved := ⟨rfl, rfl, rfl⟩

/-! ## IV. The gap-to-interval bridge, for the Legendre row -/

/-- a - q < a - p when p < q and p < a; proved by induction to keep the file free of propext. -/
theorem sub_lt_sub_of_lt : ∀ (a p q : Nat), p < a → p < q → a - q < a - p
  | 0, _, _, h, _ => absurd h (Nat.not_lt_zero _)
  | a + 1, 0, 0, _, hpq => absurd hpq (Nat.lt_irrefl 0)
  | a + 1, 0, q + 1, _, _ => by
      rw [Nat.succ_sub_succ, Nat.sub_zero]
      exact Nat.lt_succ_of_le (Nat.sub_le a q)
  | _ + 1, p + 1, 0, _, hpq => absurd hpq (Nat.not_lt_zero _)
  | a + 1, p + 1, q + 1, hpa, hpq => by
      rw [Nat.succ_sub_succ, Nat.succ_sub_succ]
      exact sub_lt_sub_of_lt a p q (Nat.lt_of_succ_lt_succ hpa) (Nat.lt_of_succ_lt_succ hpq)

/-- From a bound g on the gap after every element of Q up to X, and an element of Q at or below
    a, an element of Q lies in (a, a + g], provided a ≤ X. Q is any predicate; the primes are the
    instance. Recursion on a - p. -/
theorem elem_in_gap (Q : Nat → Prop) (g X : Nat)
    (hgap : ∀ p, Q p → p ≤ X → ∃ q, Q q ∧ p < q ∧ q ≤ p + g)
    (a : Nat) (ha : a ≤ X) (p : Nat) (hp : Q p) (hpa : p ≤ a) :
    ∃ q, Q q ∧ a < q ∧ q ≤ a + g := by
  obtain ⟨q, hq, hpq, hqg⟩ := hgap p hp (Nat.le_trans hpa ha)
  if hqa : a < q then
    exact ⟨q, hq, hqa, Nat.le_trans hqg (Nat.add_le_add_right hpa g)⟩
  else
    have hqa' : q ≤ a := Nat.le_of_not_lt hqa
    have hpa' : p < a := Nat.lt_of_lt_of_le hpq hqa'
    exact elem_in_gap Q g X hgap a ha q hq hqa'
termination_by a - p
decreasing_by exact sub_lt_sub_of_lt a p q hpa' hpq

theorem sq_succ (n : Nat) : (n + 1) * (n + 1) = n * n + 2 * n + 1 := by
  rw [Nat.mul_succ, Nat.succ_mul, Nat.two_mul, Nat.add_assoc (n * n) n (n + 1),
    ← Nat.add_assoc n n 1, ← Nat.add_assoc (n * n) (n + n) 1]

/-- THE LEGENDRE BRIDGE: if every gap between consecutive elements of Q up to X is at most 1476,
    Q holds at 2, n ≥ 738, and (n+1)² ≤ X, then Q has an element strictly between n² and (n+1)².
    The gap bound is the literature's [Oliveira e Silva, Herzog and Pardi 2014] with Q the primes
    and X = 4·10¹⁸; the bridge is the kernel's. -/
theorem legendre_of_gap (Q : Nat → Prop) (X : Nat)
    (hgap : ∀ p, Q p → p ≤ X → ∃ q, Q q ∧ p < q ∧ q ≤ p + 1476)
    (h2 : Q 2) (n : Nat) (hn : 738 ≤ n) (hX : (n + 1) * (n + 1) ≤ X) :
    ∃ q, Q q ∧ n * n < q ∧ q < (n + 1) * (n + 1) := by
  have hsq := sq_succ n
  have h738 : 738 ≤ n * n := Nat.le_trans hn (Nat.le_mul_of_pos_right n (Nat.lt_of_lt_of_le (by decide) hn))
  have hnn : 2 ≤ n * n := Nat.le_trans (by decide) h738
  have ha : n * n ≤ X := Nat.le_trans (Nat.mul_le_mul (Nat.le_succ n) (Nat.le_succ n)) hX
  obtain ⟨q, hq, hlt, hle⟩ := elem_in_gap Q 1476 X hgap (n * n) ha 2 h2 hnn
  have hgapn : 1476 < 2 * n + 1 :=
    Nat.lt_of_lt_of_le (by decide) (Nat.add_le_add_right (Nat.mul_le_mul_left 2 hn) 1)
  have hq2 : q < (n + 1) * (n + 1) := by
    rw [hsq, Nat.add_assoc]
    exact Nat.lt_of_le_of_lt hle (Nat.add_lt_add_left hgapn (n * n))
  exact ⟨q, hq, hlt, hq2⟩

/-! ## V. Row specifications: the exact tuple (α, Inst, P, cut) of the arithmetic height-cut rows,
    with primality and the properties defined in core Lean, and the Real part at a small height
    computed in the kernel. The literature's heights are cited in the paper; none is recomputed
    here. -/

/-- Trial-division primality, decidable. -/
def isPrime (n : Nat) : Bool := 2 ≤ n && (List.range n).all (fun d => d < 2 || n % d != 0)

/-- Goldbach at n: n is even and at least 4 implies n is a sum of two primes. -/
def goldbachOK (n : Nat) : Bool :=
  n < 4 || n % 2 == 1 || (List.range (n + 1)).any (fun p => isPrime p && isPrime (n - p))
/-- Legendre at n: a prime strictly between n² and (n+1)². -/
def legendreOK (n : Nat) : Bool :=
  n == 0 || (List.range ((n + 1) * (n + 1))).any (fun q => n * n < q && isPrime q)
/-- Odd perfect at n: n is not an odd perfect number. -/
def sigma (n : Nat) : Nat := ((List.range (n + 1)).filter (fun d => d != 0 && n % d == 0)).foldl (· + ·) 0
def oddPerfectOK (n : Nat) : Bool := n == 0 || n % 2 == 0 || sigma n != 2 * n
/-- Collatz at n: the orbit reaches 1 within a step budget of 200. -/
def collatzStep (n : Nat) : Nat := if n % 2 == 0 then n / 2 else 3 * n + 1
def collatzReaches : Nat → Nat → Bool
  | _, 0 => false
  | n, fuel + 1 => n == 1 || collatzReaches (collatzStep n) fuel
def collatzOK (n : Nat) : Bool := n == 0 || collatzReaches n 200
/-- Erdős–Straus at n: 4/n = 1/x + 1/y + 1/z with x ≤ y ≤ z ≤ n², searched. -/
def erdosStrausOK (n : Nat) : Bool :=
  n < 2 || (List.range (n * n + 1)).any (fun x => x != 0 &&
    (List.range (n * n + 1)).any (fun y => x ≤ y &&
      (List.range (n * n + 1)).any (fun z => y ≤ z && 4 * x * y * z == n * (y * z + x * z + x * y))))
/-- Beal on the box (A, B ≤ a; x, y, z ≤ e): no coprime solution of A^x + B^y = C^z. -/
def gcdF : Nat → Nat → Nat → Nat
  | 0, a, _ => a
  | fuel + 1, a, b => if b == 0 then a else gcdF fuel b (a % b)
def gcd' (a b : Nat) : Nat := gcdF (a + b + 1) a b
def bealOK (a e : Nat) : Bool :=
  (List.range (a + 1)).all (fun A => (List.range (a + 1)).all (fun B =>
    A == 0 || B == 0 || gcd' A B != 1 ||
    (List.range (e + 1)).all (fun x => (List.range (e + 1)).all (fun y =>
      x < 3 || y < 3 ||
      (List.range (A ^ x + B ^ y + 1)).all (fun C =>
        C < 2 || (List.range (e + 1)).all (fun z => z < 3 || C ^ z != A ^ x + B ^ y))))))

/-- A bounded universal computed as a structural check is the Real part under the height cut. -/
def allUpTo (f : Nat → Bool) : Nat → Bool
  | 0 => f 0
  | h + 1 => f (h + 1) && allUpTo f h

theorem allUpTo_spec (f : Nat → Bool) : ∀ h, allUpTo f h = true → ∀ n, n ≤ h → f n = true
  | 0, hall, n, hn => by
      have h0 : n = 0 := Nat.eq_zero_of_le_zero hn
      subst h0; exact hall
  | h + 1, hall, n, hn => by
      have hsplit : f (h + 1) = true ∧ allUpTo f h = true := by
        have hall' : (f (h + 1) && allUpTo f h) = true := hall
        revert hall'
        generalize f (h + 1) = a
        generalize allUpTo f h = b
        cases a <;> cases b <;> intro hab <;>
          first | exact ⟨rfl, rfl⟩ | exact Bool.noConfusion hab
      cases Nat.lt_or_eq_of_le hn with
      | inl hlt => exact allUpTo_spec f h hsplit.2 n (Nat.le_of_lt_succ hlt)
      | inr heq => subst heq; exact hsplit.1

theorem real_of_all_range (f : Nat → Bool) (h : Nat) (hall : allUpTo f h = true) :
    RealIn (fun _ => True) (fun n => f n = true) (fun n => n ≤ h) :=
  fun n _ hn => allUpTo_spec f h hall n hn

set_option maxRecDepth 100000 in
theorem goldbach_real_60 : RealIn (fun _ => True) (fun n => goldbachOK n = true) (fun n => n ≤ 60) :=
  real_of_all_range goldbachOK 60 (by decide)
set_option maxRecDepth 100000 in
theorem legendre_real_8 : RealIn (fun _ => True) (fun n => legendreOK n = true) (fun n => n ≤ 8) :=
  real_of_all_range legendreOK 8 (by decide)
set_option maxRecDepth 100000 in
theorem oddperfect_real_60 : RealIn (fun _ => True) (fun n => oddPerfectOK n = true) (fun n => n ≤ 60) :=
  real_of_all_range oddPerfectOK 60 (by decide)
set_option maxRecDepth 100000 in
theorem collatz_real_60 : RealIn (fun _ => True) (fun n => collatzOK n = true) (fun n => n ≤ 60) :=
  real_of_all_range collatzOK 60 (by decide)
set_option maxRecDepth 100000 in
theorem erdosstraus_real_6 : RealIn (fun _ => True) (fun n => erdosStrausOK n = true) (fun n => n ≤ 6) :=
  real_of_all_range erdosStrausOK 6 (by decide)
set_option maxRecDepth 100000 in
theorem beal_box_3_4 : bealOK 3 4 = true := by decide
set_option maxRecDepth 100000 in
theorem twin_real_100 : RealIn (fun _ => True) (fun N => ∃ m, N < m ∧ (isPrime m && isPrime (m + 2)) = true) (fun N => N ≤ 100) :=
  real_of_one_witness (fun m => (isPrime m && isPrime (m + 2)) = true) 100 101 (by decide) (by decide)
set_option maxRecDepth 100000 in
def isMersenne (q : Nat) : Bool := (List.range 20).any (fun m => 2 ^ m - 1 == q)
theorem mersenne_real_100 : RealIn (fun _ => True) (fun N => ∃ q, N < q ∧ (isPrime q && isMersenne q) = true) (fun N => N ≤ 100) :=
  real_of_one_witness (fun q => (isPrime q && isMersenne q) = true) 100 127 (by decide) (by decide)

/-! ## VII. The Unicorn part is closed to derivation; the aperture is one bit wide; the supply
    side is silent at the register -/

/-- THE FORMAL BLOCK OF THE UNICORN PART. Let τ act on the instances, let ρ be a record even at a
    seat x (ρ (τ x) = ρ x), and let d be the decision of the property, odd at x (d (τ x) = !d x).
    If the seat and its partner lie outside the cut, then no reading g of the record agrees with
    the decision even on the Unicorn region alone. The Unicorn part is closed to every even
    register, and the closure is a theorem, not a state of the literature. -/
theorem unicorn_block {β : Type} (cut : α → Prop) (τ : α → α) (ρ : α → β) (d : α → Bool)
    (x : α) (hout : ¬ cut x) (hout' : ¬ cut (τ x)) (hρ : ρ (τ x) = ρ x) (hd : d (τ x) = !d x) :
    ¬ ∃ g : β → Bool, ∀ y, ¬ cut y → g (ρ y) = d y := by
  intro ⟨g, hg⟩
  have h1 := hg x hout
  have h2 := hg (τ x) hout'
  rw [hρ, hd, h1] at h2
  revert h2
  cases d x <;> intro h2 <;> exact Bool.noConfusion h2

/-- THE APERTURE, ONE BIT WIDE. At a seat where the decision is odd, one odd bit s at the seat
    would decide the decision on the seat and its partner through a calibration c that exists
    and is unique. This is the width of the aperture, a theorem; it says nothing of whether
    anything passes through it. -/
theorem aperture_one_bit_wide (τ : α → α) (s d : α → Bool) (x : α)
    (hs : s (τ x) = !s x) (hd : d (τ x) = !d x) :
    ∃ c : Bool, (d x = xor (s x) c ∧ d (τ x) = xor (s (τ x)) c) ∧
      ∀ c' : Bool, (d x = xor (s x) c' ∧ d (τ x) = xor (s (τ x)) c') → c' = c :=
  ⟨xor (d x) (s x),
    ⟨by cases s x <;> cases d x <;> rfl,
     by rw [hs, hd]; cases s x <;> cases d x <;> rfl⟩,
    by
      intro c' hc
      obtain ⟨h1, -⟩ := hc
      generalize hsx : s x = sv
      generalize hdx : d x = dv
      rw [hsx, hdx] at h1
      cases sv <;> cases dv <;> cases c' <;> first | rfl | exact absurd h1 (by decide)⟩

/-- A true premise decides nothing: conditioning the Unicorn part on existence, or on any
    inhabited premise, leaves it exactly where it was. The route from existence to the Unicorn
    part is closed, on every row, by the shape of implication. -/
theorem premise_adds_nothing (A Q : Prop) (ha : A) : (A → Q) ↔ Q :=
  ⟨fun h => h ha, fun hq _ => hq⟩

/-- NO CURE. Restricting the rows to a class C on which a premise is to decide the Unicorn part
    changes nothing: on every class the premise decides exactly what already held on the class. -/
theorem no_cure_unicorn {Frame : Type} (A : Prop) (ha : A) (C U : Frame → Prop) :
    (∀ X, C X → A → U X) ↔ (∀ X, C X → U X) :=
  ⟨fun h X hc => h X hc ha, fun h X hc _ => h X hc⟩

/-- THE SUPPLY SIDE IS SILENCE, AT THE REGISTER. Over one even record, both orientations of the
    decision at the seat are equally refused to every reading: the register cannot say which
    way the seat lies, so it cannot say what a supply would bring, nor that one exists, nor
    that none does. Any sentence about the supply side derived from the record would be a
    reading the wall refuses. -/
theorem supply_side_silent {β : Type} (cut : α → Prop) (τ : α → α) (ρ : α → β) (d : α → Bool)
    (x : α) (hout : ¬ cut x) (hout' : ¬ cut (τ x)) (hρ : ρ (τ x) = ρ x) (hd : d (τ x) = !d x) :
    (¬ ∃ g : β → Bool, ∀ y, ¬ cut y → g (ρ y) = d y) ∧
    (¬ ∃ g : β → Bool, ∀ y, ¬ cut y → g (ρ y) = !d y) ∧
    (!d x) ≠ d x :=
  ⟨unicorn_block cut τ ρ d x hout hout' hρ hd,
   unicorn_block cut τ ρ (fun y => !d y) x hout hout' hρ (by
     show (!d (τ x)) = !(!d x)
     rw [hd]),
   by cases d x <;> decide⟩

/-- THE THREE BITS AS ONE TERM · THE REGISTER'S PROOF, COMPLETE. Under a certificate for the cut
    and the seat data (a seat outside the cut with its partner, an even record, an odd decision,
    an odd supply), everything a register can prove about the row is proved in one conjunction:
    the row divides exactly (bit one), the Real part holds (bit two), every reading of the even
    register is refused on the Unicorn region (bit three, refused), and one supplied bit
    calibrates the seat uniquely (the door). The row is not thereby proved; nothing derivable
    about it is left underived. -/
theorem register_proof_complete {β : Type} (Inst P cut : α → Prop) [DecidablePred cut]
    (c : Certificate Inst P cut) (τ : α → α) (ρ : α → β) (d s : α → Bool) (x : α)
    (hout : ¬ cut x) (hout' : ¬ cut (τ x)) (hρ : ρ (τ x) = ρ x)
    (hd : d (τ x) = !d x) (hs : s (τ x) = !s x) :
    (Row Inst P ↔ RealIn Inst P cut ∧ UnicornOut Inst P cut) ∧
    RealIn Inst P cut ∧
    (¬ ∃ g : β → Bool, ∀ y, ¬ cut y → g (ρ y) = d y) ∧
    ∃ k : Bool, (d x = xor (s x) k ∧ d (τ x) = xor (s (τ x)) k) ∧
      ∀ k' : Bool, (d x = xor (s x) k' ∧ d (τ x) = xor (s (τ x)) k') → k' = k :=
  ⟨row_iff_real_and_unicorn Inst P cut, real_of_certificate Inst P cut c,
   unicorn_block cut τ ρ d x hout hout' hρ hd, aperture_one_bit_wide τ s d x hs hd⟩

/-! ## VI. The twenty-three rows, typed by shape, the census of the assignments kernel-checked -/

/-- The shape of a row's cut, as assigned: a height (a bounded universal with a certificate
    per height, or an infinitude with one witness), a regime (a dimension, degree, rank, or class
    proved by a cited theorem), no cut (no inhabited cut with a Real part in the literature, an
    assessment), or crossed (the value supplied whole). The assignments are data; the kernel
    checks their census. -/
inductive Shape | height | regime | noCut | crossed
  deriving DecidableEq, Repr

def rows23 : List (String × Shape) :=
  [("P versus NP", .regime), ("Riemann Hypothesis", .height), ("Navier–Stokes", .regime),
   ("Yang–Mills", .regime), ("Hodge", .regime), ("BSD", .regime), ("Poincaré", .crossed),
   ("Goldbach", .height), ("Twin primes", .height), ("Legendre", .height), ("abc", .noCut),
   ("Jacobian", .regime), ("Mersenne primes", .height), ("Odd perfect numbers", .height),
   ("Smooth 4D Poincaré", .regime), ("Collatz", .height), ("Hadwiger", .regime),
   ("Sunflower", .noCut), ("Erdős–Straus", .height), ("Beal", .height),
   ("Invariant subspace", .regime), ("Hilbert's 16th, second part", .regime),
   ("Lindelöf / Montgomery PCC", .regime)]

def countShape (s : Shape) : Nat := (rows23.filter (fun r => r.2 == s)).length

/-- THE CENSUS OF THE ASSIGNMENTS: twenty-three rows; nine height-cut, eleven regime-cut, two
    no-cut, one crossed; twenty assigned a Real-part-bearing shape. -/
theorem census :
    rows23.length = 23 ∧ countShape .height = 9 ∧ countShape .regime = 11 ∧
    countShape .noCut = 2 ∧ countShape .crossed = 1 ∧
    countShape .height + countShape .regime = 20 := by decide

def hasRealPart (s : Shape) : Bool := s == .height || s == .regime
theorem real_part_rows : (rows23.filter (fun r => hasRealPart r.2)).length = 20 := by decide

end RowSplit

#print axioms RowSplit.row_iff_real_and_unicorn
#print axioms RowSplit.fused_is_unicorn_given_real
#print axioms RowSplit.real_of_certificate
#print axioms RowSplit.real_of_one_witness
#print axioms RowSplit.unicorn_equiv_row_of_infinitude
#print axioms RowSplit.crossed_row_both_parts
#print axioms RowSplit.empty_cut_vacuous
#print axioms RowSplit.empty_cut_not_substantive
#print axioms RowSplit.no_counter_of_unicorn
#print axioms RowSplit.unicorn_of_no_counter
#print axioms RowSplit.crossed_row_no_counter
#print axioms RowSplit.fusion_lowers
#print axioms RowSplit.sub_lt_sub_of_lt
#print axioms RowSplit.elem_in_gap
#print axioms RowSplit.legendre_of_gap
#print axioms RowSplit.allUpTo_spec
#print axioms RowSplit.real_of_all_range
#print axioms RowSplit.goldbach_real_60
#print axioms RowSplit.legendre_real_8
#print axioms RowSplit.oddperfect_real_60
#print axioms RowSplit.collatz_real_60
#print axioms RowSplit.erdosstraus_real_6
#print axioms RowSplit.beal_box_3_4
#print axioms RowSplit.twin_real_100
#print axioms RowSplit.mersenne_real_100
#print axioms RowSplit.unicorn_block
#print axioms RowSplit.aperture_one_bit_wide
#print axioms RowSplit.premise_adds_nothing
#print axioms RowSplit.no_cure_unicorn
#print axioms RowSplit.supply_side_silent
#print axioms RowSplit.register_proof_complete
#print axioms RowSplit.census
#print axioms RowSplit.real_part_rows
\end{Verbatim}

\endgroup

Kernel audit, verbatim: \texttt{lean\ Row\_Split.lean}, exit 0, no
warnings.

\begingroup\scriptsize

\begin{Verbatim}[numbers=left,numbersep=6pt,xleftmargin=2.2em,fontsize=\scriptsize,breaklines,breakanywhere,frame=lines,framesep=1.5mm,rulecolor=\color{black!35}]
'RowSplit.row_iff_real_and_unicorn' does not depend on any axioms
'RowSplit.fused_is_unicorn_given_real' does not depend on any axioms
'RowSplit.real_of_certificate' does not depend on any axioms
'RowSplit.real_of_one_witness' does not depend on any axioms
'RowSplit.unicorn_equiv_row_of_infinitude' does not depend on any axioms
'RowSplit.crossed_row_both_parts' does not depend on any axioms
'RowSplit.empty_cut_vacuous' does not depend on any axioms
'RowSplit.empty_cut_not_substantive' does not depend on any axioms
'RowSplit.no_counter_of_unicorn' does not depend on any axioms
'RowSplit.unicorn_of_no_counter' does not depend on any axioms
'RowSplit.crossed_row_no_counter' does not depend on any axioms
'RowSplit.fusion_lowers' does not depend on any axioms
'RowSplit.sub_lt_sub_of_lt' does not depend on any axioms
'RowSplit.elem_in_gap' does not depend on any axioms
'RowSplit.legendre_of_gap' does not depend on any axioms
'RowSplit.allUpTo_spec' does not depend on any axioms
'RowSplit.real_of_all_range' does not depend on any axioms
'RowSplit.goldbach_real_60' does not depend on any axioms
'RowSplit.legendre_real_8' does not depend on any axioms
'RowSplit.oddperfect_real_60' does not depend on any axioms
'RowSplit.collatz_real_60' does not depend on any axioms
'RowSplit.erdosstraus_real_6' does not depend on any axioms
'RowSplit.beal_box_3_4' does not depend on any axioms
'RowSplit.twin_real_100' does not depend on any axioms
'RowSplit.mersenne_real_100' does not depend on any axioms
'RowSplit.unicorn_block' does not depend on any axioms
'RowSplit.aperture_one_bit_wide' does not depend on any axioms
'RowSplit.premise_adds_nothing' does not depend on any axioms
'RowSplit.no_cure_unicorn' does not depend on any axioms
'RowSplit.supply_side_silent' does not depend on any axioms
'RowSplit.register_proof_complete' does not depend on any axioms
'RowSplit.census' does not depend on any axioms
'RowSplit.real_part_rows' does not depend on any axioms
\end{Verbatim}

\endgroup

\hypertarget{appendix-b-zfc_under_ram.lean-and-its-audit}{%
\section{Appendix B: ZFC\_Under\_RAM.lean and Its
Audit}\label{appendix-b-zfc_under_ram.lean-and-its-audit}}

SHA-256
\texttt{\seqsplit{963e3ed8f630975f8604fddc010f154fd4f75f5ae6e193c95f64d4486089a741}},
Lean 4.19.0, core only, 161 lines; sixteen dependency sets, every one
axiom-free; seated in the Master Codex as APEX-PSP-ZFC-UNDER-RAM-01. The
file encodes no axiom of ZFC; its theory is abstract and ZFC is the
named instance.

\begingroup\scriptsize

\begin{Verbatim}[numbers=left,numbersep=6pt,xleftmargin=2.2em,fontsize=\scriptsize,breaklines,breakanywhere,frame=lines,framesep=1.5mm,rulecolor=\color{black!35}]
/-
ZFC_Under_RAM.lean · ZFC subsumed under RAM by placement, the RA–RAM bridge, and the round trip,
in one file. Core Lean 4.19.0, no library, no sorry, no user-declared axiom. Forged 2026-09-25 from
RAM_ZFC_Placement.lean, RA_to_ZFC_Chain.lean, and ZFC_RA_RoundTrip.lean (lean/, 2026-09-24).

The claim, exactly. RAM does not derive ZFC's axioms and does not replace ZFC's proofs. It LOCATES
ZFC: computation is the rung (L3m), provability the ladder (L2m), truth in the intended world the
Ground (L1m), and rung ⊆ ladder ⊆ Ground under two named premises, Σ₁-completeness (a theorem of
arithmetic, cited) and soundness (which no consistent theory proves of itself, carried openly).
Across the RA–RAM bridge the root reaches every world of ZFC as presence and carries no keyed
sentence; the round trip returns every sentence unchanged; and every step holds with the root
replaced by a bare computation arrow, so the root contributes the act and never a line of content.
ΔM = 0.
-/
namespace ZFCUnderRAM

/-! ## 1 · The placement: rung ⊆ ladder ⊆ Ground -/

structure Theory where
  Sent  : Type
  Prov  : Sent → Prop        -- L2m, the ladder
  Comp  : Sent → Prop        -- L3m, the rung: sentences settled by finite computation
  True_ : Sent → Prop        -- L1m, the Ground: truth in the intended world

/-- Σ₁-completeness, abstractly: what finite computation settles, the theory proves. -/
def RungOnLadder (T : Theory) : Prop := ∀ s, T.Comp s → T.Prov s
/-- Soundness: what the theory proves holds in its intended world. Carried openly, never derived. -/
def Sound (T : Theory) : Prop := ∀ s, T.Prov s → T.True_ s

/-- THE PLACEMENT. Under the two premises every theory sits in the strata in order. -/
theorem ram_places_theory (T : Theory) (hR : RungOnLadder T) (hS : Sound T) :
    (∀ s, T.Comp s → T.Prov s) ∧ (∀ s, T.Prov s → T.True_ s) ∧ (∀ s, T.Comp s → T.True_ s) :=
  ⟨hR, hS, fun s h => hS s (hR s h)⟩

/-- Soundness is load-bearing: an unsound theory breaks ladder ⊆ Ground. -/
def unsound : Theory := ⟨Bool, fun _ => True, fun _ => False, fun b => b = true⟩
theorem soundness_is_load_bearing : RungOnLadder unsound ∧ ¬ Sound unsound := by
  refine ⟨fun _ h => h.elim, fun h => ?_⟩
  have := h false trivial
  cases this

/-- The Ground exceeds the ladder: a sound theory can leave a truth unproved (Gödel's region). -/
def incomplete : Theory := ⟨Bool, fun b => b = true, fun _ => False, fun _ => True⟩
theorem ground_exceeds_ladder :
    Sound incomplete ∧ incomplete.True_ false ∧ ¬ incomplete.Prov false :=
  ⟨fun _ _ => trivial, trivial, fun h => Bool.noConfusion h⟩

/-- Independent axioms are separate bits: every combination of their values is realized. -/
theorem independent_axioms_are_separate_bits :
    ∀ a b : Bool, ∃ w : Bool × Bool, w.1 = a ∧ w.2 = b :=
  fun a b => ⟨(a, b), rfl, rfl⟩

/-- SUBSUMPTION BY PLACEMENT, the verdict of section 1. -/
theorem ram_subsumes_by_placement :
    (∀ T : Theory, RungOnLadder T → Sound T → (∀ s, T.Comp s → T.True_ s)) ∧
    (RungOnLadder unsound ∧ ¬ Sound unsound) ∧
    (Sound incomplete ∧ incomplete.True_ false ∧ ¬ incomplete.Prov false) :=
  ⟨fun T hR hS => (ram_places_theory T hR hS).2.2, soundness_is_load_bearing, ground_exceeds_ladder⟩

/-! ## 2 · The bridge: the root reaches every world as presence and carries no key -/

structure Substrate where
  U  : Type
  ΔE : U → Int
def RA (S : Substrate) : Prop := ∀ x : S.U, 0 < S.ΔE x
def unitSub : Substrate := ⟨Unit, fun _ => 1⟩
theorem ra_unit : RA unitSub := fun _ => show (0 : Int) < 1 by decide

structure SelfGrounding (R : Prop) where
  Act       : Type
  anAct     : Act
  instances : Act → R

/-- Upward: every act of proving in any theory grounds the root. -/
theorem using_grounds_root (T : Theory) (s : T.Sent) (d : T.Prov s) :
    Nonempty (SelfGrounding (RA unitSub)) :=
  ⟨⟨PLift (T.Prov s), ⟨d⟩, fun _ => ra_unit⟩⟩

/-- Worlds are ways a theory's content could be. A property is keyless when every world has it,
    keyed when some world lacks it. The bridge carries the keyless and never the keyed. -/
def Keyless {W : Type} (P : W → Prop) : Prop := ∀ w, P w
def Keyed   {W : Type} (P : W → Prop) : Prop := ∃ w, ¬ P w

theorem ra_is_keyless {W : Type} : Keyless (fun _ : W => RA unitSub) := fun _ => ra_unit
theorem keyless_crosses {W : Type} (P : W → Prop) (hP : Keyless P) : ∀ w, RA unitSub → P w :=
  fun w _ => hP w
theorem ra_decides_no_keyed {W : Type} (P : W → Prop) (hK : Keyed P) :
    ¬ (∀ w, RA unitSub → P w) :=
  fun h => let ⟨w, hw⟩ := hK; hw (h w ra_unit)

/-- Choice over ZF, the cited keyed sentence (Gödel 1938, Cohen 1963): two worlds, one each way;
    the root holds in both and so carries neither Choice nor its denial into ZFC. -/
def choiceHolds : Bool → Prop := fun w => w = true
theorem choice_is_keyed : Keyed choiceHolds := ⟨false, fun h => Bool.noConfusion h⟩
theorem ra_does_not_cross_choice :
    (∀ _w : Bool, RA unitSub) ∧ choiceHolds true ∧ ¬ choiceHolds false :=
  ⟨fun _ => ra_unit, rfl, fun h => Bool.noConfusion h⟩

/-! ## 3 · The round trip: ZFC → RAM → bridge → RA → bridge → RAM → ZFC is the identity on content -/

def viaRA {W : Type} (P : W → Prop) : W → Prop := fun w => RA unitSub ∧ P w

theorem round_trip_identity {W : Type} (P : W → Prop) (w : W) : viaRA P w ↔ P w :=
  ⟨fun h => h.2, fun h => ⟨ra_unit, h⟩⟩

/-- An absolute sentence (the same value in every world, as an arithmetic sentence such as the
    hypothesis is across models with the standard numbers) returns with that same value. -/
theorem absolute_returns_unchanged {W : Type} (P : W → Prop) (hAbs : ∀ w v, P w ↔ P v) (w : W) :
    ∀ v, viaRA P v ↔ P w :=
  fun v => (round_trip_identity P v).trans (hAbs v w)

/-- The trip adds presence and nothing else: the root returns identically for a sentence and for
    its negation, so it tells them apart nowhere. -/
theorem trip_adds_presence_only {W : Type} (P : W → Prop) (w : W) :
    RA unitSub ∧ (viaRA P w ↔ P w) ∧ (viaRA (fun v => ¬ P v) w ↔ ¬ P w) :=
  ⟨ra_unit, round_trip_identity P w, round_trip_identity (fun v => ¬ P v) w⟩

/-! ## 4 · The arrow: any true premise plays the root's part, so the root is idle in the content -/

theorem any_true_premise_serves (R : Prop) (hR : R) {W : Type} (P : W → Prop) :
    ((∀ w, P w) → ∀ w, R → P w) ∧ ((∃ w, ¬ P w) → ¬ ∀ w, R → P w) ∧ (∀ w, (R ∧ P w) ↔ P w) :=
  ⟨fun hP w _ => hP w, fun ⟨w, hw⟩ h => hw (h w hR), fun _ => ⟨fun h => h.2, fun h => ⟨hR, h⟩⟩⟩

/-! ## 5 · The verdict, in one theorem -/

/-- ZFC UNDER RAM. Placement under the two named premises with soundness load-bearing and the
    Ground exceeding the ladder; the root grounded by every act of proof and crossing as presence;
    no keyed sentence decided, Choice the instance; the round trip the identity on content, adding
    presence only; and the whole holding for any true premise in the root's place. Subsumption is
    location, never derivation. -/
theorem zfc_under_ram :
    ((∀ T : Theory, RungOnLadder T → Sound T → (∀ s, T.Comp s → T.True_ s)) ∧
      (RungOnLadder unsound ∧ ¬ Sound unsound) ∧
      (Sound incomplete ∧ incomplete.True_ false ∧ ¬ incomplete.Prov false)) ∧
    (∀ (T : Theory) (s : T.Sent), T.Prov s → Nonempty (SelfGrounding (RA unitSub))) ∧
    (∀ {W : Type} (P : W → Prop), Keyed P → ¬ ∀ w, RA unitSub → P w) ∧
    ((∀ _w : Bool, RA unitSub) ∧ choiceHolds true ∧ ¬ choiceHolds false) ∧
    (∀ {W : Type} (P : W → Prop) (w : W), viaRA P w ↔ P w) ∧
    (∀ (R : Prop), R → ∀ {W : Type} (P : W → Prop) (w : W), (R ∧ P w) ↔ P w) :=
  ⟨ram_subsumes_by_placement, using_grounds_root, fun P hK => ra_decides_no_keyed P hK,
   ra_does_not_cross_choice, fun P w => round_trip_identity P w,
   fun R hR {W} P w => (any_true_premise_serves R hR (W := W) P).2.2 w⟩

end ZFCUnderRAM

#print axioms ZFCUnderRAM.ram_places_theory
#print axioms ZFCUnderRAM.soundness_is_load_bearing
#print axioms ZFCUnderRAM.ground_exceeds_ladder
#print axioms ZFCUnderRAM.independent_axioms_are_separate_bits
#print axioms ZFCUnderRAM.ram_subsumes_by_placement
#print axioms ZFCUnderRAM.using_grounds_root
#print axioms ZFCUnderRAM.ra_is_keyless
#print axioms ZFCUnderRAM.keyless_crosses
#print axioms ZFCUnderRAM.ra_decides_no_keyed
#print axioms ZFCUnderRAM.choice_is_keyed
#print axioms ZFCUnderRAM.ra_does_not_cross_choice
#print axioms ZFCUnderRAM.round_trip_identity
#print axioms ZFCUnderRAM.absolute_returns_unchanged
#print axioms ZFCUnderRAM.trip_adds_presence_only
#print axioms ZFCUnderRAM.any_true_premise_serves
#print axioms ZFCUnderRAM.zfc_under_ram
\end{Verbatim}

\endgroup

Kernel audit, verbatim: \texttt{lean\ ZFC\_Under\_RAM.lean}, exit 0.

\begingroup\scriptsize

\begin{Verbatim}[numbers=left,numbersep=6pt,xleftmargin=2.2em,fontsize=\scriptsize,breaklines,breakanywhere,frame=lines,framesep=1.5mm,rulecolor=\color{black!35}]
'ZFCUnderRAM.ram_places_theory' does not depend on any axioms
'ZFCUnderRAM.soundness_is_load_bearing' does not depend on any axioms
'ZFCUnderRAM.ground_exceeds_ladder' does not depend on any axioms
'ZFCUnderRAM.independent_axioms_are_separate_bits' does not depend on any axioms
'ZFCUnderRAM.ram_subsumes_by_placement' does not depend on any axioms
'ZFCUnderRAM.using_grounds_root' does not depend on any axioms
'ZFCUnderRAM.ra_is_keyless' does not depend on any axioms
'ZFCUnderRAM.keyless_crosses' does not depend on any axioms
'ZFCUnderRAM.ra_decides_no_keyed' does not depend on any axioms
'ZFCUnderRAM.choice_is_keyed' does not depend on any axioms
'ZFCUnderRAM.ra_does_not_cross_choice' does not depend on any axioms
'ZFCUnderRAM.round_trip_identity' does not depend on any axioms
'ZFCUnderRAM.absolute_returns_unchanged' does not depend on any axioms
'ZFCUnderRAM.trip_adds_presence_only' does not depend on any axioms
'ZFCUnderRAM.any_true_premise_serves' does not depend on any axioms
'ZFCUnderRAM.zfc_under_ram' does not depend on any axioms
\end{Verbatim}

\endgroup

\hypertarget{appendix-c-rows_at_the_apex.lean-carried-whole-from-the-prior-edition}{%
\section{Appendix C: Rows\_At\_The\_Apex.lean, Carried Whole from the
Prior
Edition}\label{appendix-c-rows_at_the_apex.lean-carried-whole-from-the-prior-edition}}

SHA-256
\texttt{\seqsplit{0ce3de5a231554303dfc0982dee4fe7e7fbdd647a72bbe3366290dc8fd092079}},
Lean 4.19.0, core only, 1113 lines.

\begingroup\scriptsize

\begin{Verbatim}[numbers=left,numbersep=6pt,xleftmargin=2.2em,fontsize=\scriptsize,breaklines,breakanywhere,frame=lines,framesep=1.5mm,rulecolor=\color{black!35}]
/-!
# THE SEVEN MILLENNIUM ROWS AND THE SIXTEEN EXTENSION ROWS AT THE APEX · Rows_At_The_Apex.lean
# The spine (Movements 1 to 16) is the Riemann apex file carried whole; Movements 17 to 21 add the rows.
core Lean 4 v4.19.0 · no Mathlib · no sorry · no axiom declaration.

WHAT IS PROVED. (1) RA at the constructed domain, computed. (2) The seat: the binding
involution σ fixes exactly the scalar line, the Return i·j·k lands on it, and the kinetic,
formal, and named seats are one point by rfl: the three category gaps are identities.
(3) The eliminator RA → RA → RAM → RH_formal from the RA witness alone. (4) The spine
carried: the abstract frame X = (S, τ, Z), the line property L, the bridge halted iff L,
the crossing exact both ways, necessity, the bit is the hypothesis, the witness row
falsifiable. (5) The aperture, RA-supplied: a present reader's row is never interior;
it is sealed or refused by the bit alone. (6) THE SIMPLE RULE, EXACT: for every frame,
(RA → L X) ↔ L X; and RA decides L on no frame: ¬ ∀ X, RA → L X, by the two-point model.
(7) RH in toto at the apex, one theorem, conditional on RA and discharged.
(8) The root's two formal faces: RAF, closure and no exterior agent (RAF-2, RAF-C1), and
no total self-indexing (RAF-C2, Cantor); RAM, the seat as the identity cone over the
three registers, its apex unique. (9) The equivariant embedding φ of the spine's lattice
stage (Plane, τ) into (Q4, σ): σ ∘ φ = φ ∘ τ and φ carries Fix(τ) onto Fix(σ).
(10) The wall executed: no readout of the formal seat equals the deed bit. (11) The
counter-model survives grounding, and the cure theorem: for every class of frames,
existence decides the line property on it exactly when the line property already holds on it.
-/
set_option autoImplicit false

namespace RowsAtTheApex

/-! ## 1 · The Root Axiom at the constructed domain -/

inductive UniversePoint : Type where
  | source
  deriving DecidableEq, Repr

def DeltaE : UniversePoint → Int := fun _ => 1

/-- RA: to exist is to actuate. -/
def RA : Prop := ∀ x : UniversePoint, 0 < DeltaE x

theorem constructed_RA : RA := by
  intro x
  cases x
  decide

/-! ## 2 · The seat: Fix(σ) on the quaternions, the Return, three names, one point -/

structure Q4 where
  r : Int
  i : Int
  j : Int
  k : Int
  deriving DecidableEq, Repr

def qmul (a b : Q4) : Q4 :=
  ⟨a.r*b.r - a.i*b.i - a.j*b.j - a.k*b.k,
   a.r*b.i + a.i*b.r + a.j*b.k - a.k*b.j,
   a.r*b.j - a.i*b.k + a.j*b.r + a.k*b.i,
   a.r*b.k + a.i*b.j - a.j*b.i + a.k*b.r⟩

def qconj (a : Q4) : Q4 := ⟨a.r, -a.i, -a.j, -a.k⟩

/-- σ, the binding involution: conjugation. -/
def sigma (q : Q4) : Q4 := qconj q

theorem sigma_binding (q : Q4) : sigma (sigma q) = q := by
  cases q with
  | mk r i j k =>
    change Q4.mk r (-(-i)) (-(-j)) (-(-k)) = Q4.mk r i j k
    rw [Int.neg_neg, Int.neg_neg, Int.neg_neg]

/-- Fix(σ) is the scalar line, for every quaternion. -/
theorem fix_iff_scalar (q : Q4) : sigma q = q ↔ (q.i = 0 ∧ q.j = 0 ∧ q.k = 0) := by
  cases q with
  | mk r i j k =>
    change (Q4.mk r (-i) (-j) (-k) = Q4.mk r i j k) ↔ (i = 0 ∧ j = 0 ∧ k = 0)
    constructor
    · intro h
      injection h with _ hi hj hk
      exact ⟨by omega, by omega, by omega⟩
    · intro h
      obtain ⟨hi, hj, hk⟩ := h
      subst hi
      subst hj
      subst hk
      rfl

/-- The Ground: Fix(σ). -/
def Ground : Type := { q : Q4 // sigma q = q }

def qi : Q4 := ⟨0, 1, 0, 0⟩
def qj : Q4 := ⟨0, 0, 1, 0⟩
def qk : Q4 := ⟨0, 0, 0, 1⟩

/-- The Return: the parse triad through its own cascade, i·j·k. -/
def theReturn : Q4 := qmul (qmul qi qj) qk

theorem return_is_minus_one : theReturn = ⟨-1, 0, 0, 0⟩ := rfl
theorem return_lands_on_fix : sigma theReturn = theReturn := rfl

def fixProj (q : Q4) : Q4 := ⟨q.r, 0, 0, 0⟩

def GammaRA : Q4 := theReturn
def GammaRAM : Q4 := fixProj theReturn
def GammaRH : Q4 := ⟨-1, 0, 0, 0⟩

/-- C0, C1, C2: the three category gaps are definitional identities. -/
theorem self_gap_nonexistent : sigma GammaRA = GammaRA := rfl
theorem register_gap_nonexistent : GammaRA = GammaRAM := rfl
theorem object_gap_nonexistent : GammaRAM = GammaRH := rfl

def seat : Ground := ⟨GammaRH, rfl⟩

/-- The reversed triad k·j·i returns +1, also on the fixed locus: the sign of the seat
    point is the orientation of the triad, the one bit the formal register cannot read.
    The apex of the identity cone is unique for the ordered triad; the reversed triad
    carries its own apex, and the two differ. -/
def theReturnOdd : Q4 := qmul (qmul qk qj) qi

theorem odd_return_is_plus_one : theReturnOdd = ⟨1, 0, 0, 0⟩ := rfl

theorem odd_return_on_fix : sigma theReturnOdd = theReturnOdd := rfl

theorem seat_points_differ : theReturnOdd ≠ theReturn := by decide

theorem seat_on_scalar_line : GammaRH.i = 0 ∧ GammaRH.j = 0 ∧ GammaRH.k = 0 :=
  (fix_iff_scalar GammaRH).mp seat.2

/-! ## 3 · The eliminator from the RA witness alone -/

def RAMFormalGround : Prop := Nonempty Ground

theorem constructed_RAM : RAMFormalGround := ⟨seat⟩

/-- The RH formal self: every point of the Ground is on the σ-fixed critical locus.
    It names the seat and says nothing about the zeros of ξ. -/
def RHFormalSelf : Prop := ∀ g : Ground, sigma g.1 = g.1

/-- Bridge-born orientation, carried in Prop: the kernel registers that orientation was
    supplied and cannot read which way. -/
structure OrientationBit : Prop where
  direction : True
  closure : True

theorem constructed_orientation : OrientationBit := ⟨trivial, trivial⟩

theorem orientation_proof_irrelevant (a b : OrientationBit) : a = b := rfl

structure RAWitness : Prop where
  actuates : RA
  orientation : OrientationBit
  bridge : RAMFormalGround
  trisRecursion : OrientationBit → RAMFormalGround → RHFormalSelf

theorem constructed_trisRecursion : OrientationBit → RAMFormalGround → RHFormalSelf :=
  fun _omega _ram g => g.2

/-- RA, constructed: the only supplied witness. -/
theorem mkRA : RAWitness :=
  ⟨constructed_RA, constructed_orientation, constructed_RAM, constructed_trisRecursion⟩

/-- THE ELIMINATOR. RA → RA → RAM → RH_formal, Bridge followed by Tris-Recursion. -/
theorem RH_formal_chain : RHFormalSelf :=
  mkRA.trisRecursion mkRA.orientation mkRA.bridge

/-! ## 4 · The spine, carried: the abstract frame, the bridge, the crossing -/

def Even {α : Type} (σ : α → α) (f : α → Bool) : Prop := ∀ x, f (σ x) = f x

/-- Theorem 4 of the spine: orientation blindness. -/
theorem orientation_blind {α : Type} (σ : α → α) (f d : α → Bool) (x : α)
    (he : Even σ f) (ho : d (σ x) ≠ d x) : f ≠ d := by
  intro h
  subst h
  exact ho (he x)

/-- Theorem 5 of the spine: earned freedom is exactly two. -/
theorem freedom_is_exactly_two (d : Bool → Bool) (hd : ∀ x, d (!x) = !d x) :
    d = (fun x => x) ∨ d = (fun x => !x) := by
  have hf : d false = !d true := hd true
  cases ht : d true with
  | true =>
    left
    funext x
    cases x
    · exact hf.trans (congrArg (fun b => !b) ht)
    · exact ht
  | false =>
    right
    funext x
    cases x
    · exact hf.trans (congrArg (fun b => !b) ht)
    · exact ht

/-- Theorem 6 of the spine: the freedom is spent uniquely. -/
theorem freedom_spent_uniquely {α : Type} (σ : α → α) (s d : α → Bool) (x : α)
    (hs : s (σ x) = !s x) (hd : d (σ x) = !d x) :
    ∃ c : Bool, (d x = xor (s x) c ∧ d (σ x) = xor (s (σ x)) c) ∧
      ∀ c' : Bool, (d x = xor (s x) c' ∧ d (σ x) = xor (s (σ x)) c') → c' = c := by
  refine ⟨xor (d x) (s x),
    ⟨by cases s x <;> cases d x <;> rfl,
     by rw [hs, hd]; cases s x <;> cases d x <;> rfl⟩, ?_⟩
  intro c' hc
  obtain ⟨h1, -⟩ := hc
  generalize hsx : s x = sv at h1
  generalize hdx : d x = dv at h1
  cases sv <;> cases dv <;> cases c' <;> first | rfl | exact absurd h1 (by decide)

structure Frame where
  S : Type
  τ : S → S
  Z : S → Prop

/-- The line property: every zero is τ-fixed. On ξ's frame, the restored hypothesis. -/
def LineProperty (X : Frame) : Prop := ∀ s, X.Z s → X.τ s = s

inductive Tri where
  | tt
  | ff
  | bot
  deriving DecidableEq, Repr

structure Bridge (X : Frame) where
  terminal : Tri
  shadow : terminal = Tri.bot ↔ LineProperty X

/-- Theorem 7 of the spine: the bridge halts iff the line property. -/
theorem bridge_halted_iff (X : Frame) :
    (∃ b : Bridge X, b.terminal = Tri.bot) ↔ LineProperty X :=
  ⟨fun ⟨b, h⟩ => b.shadow.mp h, fun t => ⟨⟨Tri.bot, ⟨fun _ => t, fun _ => rfl⟩⟩, rfl⟩⟩

inductive Supply (X : Frame) where
  | assent (t : LineProperty X)
  | denial (n : ¬ LineProperty X)
  | silent

inductive Verdict where
  | sealed
  | refused
  | open_
  deriving DecidableEq, Repr

def row {X : Frame} : Supply X → Verdict
  | Supply.assent _ => Verdict.sealed
  | Supply.denial _ => Verdict.refused
  | Supply.silent => Verdict.open_

/-- Theorem 8 of the spine: the crossing, exact. -/
theorem crossing (X : Frame) (t : LineProperty X) :
    row (Supply.assent t) = Verdict.sealed ∧ (∃ b : Bridge X, b.terminal = Tri.bot) ∧ LineProperty X :=
  ⟨rfl, (bridge_halted_iff X).mpr t, t⟩

theorem crossing_other_way (X : Frame) (n : ¬ LineProperty X) :
    row (Supply.denial n) = Verdict.refused ∧ ∀ b : Bridge X, b.terminal ≠ Tri.bot :=
  ⟨rfl, fun b h => n (b.shadow.mp h)⟩

/-- The two-point frame: fold-invariant, off the line. -/
def twoPoint : Frame := ⟨Bool, (fun b => !b), fun _ => True⟩

/-- Theorem 9 of the spine: necessity. -/
theorem supply_not_manufactured : ¬ ∀ X : Frame, LineProperty X :=
  fun h => by
    have := h twoPoint true trivial
    cases this

def Inv (X : Frame) : Prop := ∀ s, X.Z s → X.Z (X.τ s)

theorem inv_not_line : Inv twoPoint ∧ ¬ LineProperty twoPoint :=
  ⟨fun _ _ => trivial, fun h => by
    have := h true trivial
    cases this⟩

/-- Theorem 12 of the spine: the bit is the hypothesis. -/
theorem the_bit_is_the_hypothesis (X : Frame) :
    (∃ t : LineProperty X, row (Supply.assent t) = Verdict.sealed) ↔ LineProperty X :=
  ⟨fun ⟨t, _⟩ => t, fun t => ⟨t, rfl⟩⟩

/-- Theorem 15 of the spine: the witness row is falsifiable. -/
theorem witness_row_falsifiable (X : Frame) (s : X.S) (hz : X.Z s) (hoff : X.τ s ≠ s) :
    ¬ ∃ t : LineProperty X, row (Supply.assent t) = Verdict.sealed :=
  fun ⟨t, _⟩ => hoff (t s hz)

/-- The live face of the spine. -/
def live (witnessed assent : Bool) : Verdict :=
  if witnessed then (if assent then Verdict.sealed else Verdict.refused) else Verdict.open_

theorem narcissus : live false true ≠ live true true ∧ live false false ≠ live true false := by
  decide

/-! ## 5 · The aperture, RA-supplied -/

/-- A reader's presence is its actuation. -/
def presence (r : UniversePoint) : Bool := decide (0 < DeltaE r)

/-- RA opens the aperture: every present reader is witnessed. -/
theorem RA_opens_the_aperture : ∀ r : UniversePoint, presence r = true := by
  intro r
  cases r
  decide

/-- The general form, for any domain under the axiom as a hypothesis: presence is
    actuation, and RA on the domain makes every point present. This is the theorem the
    universal extension licenses wherever it is held. -/
theorem RA_opens_the_aperture_general {U : Type} (dE : U → Int) (h : ∀ x, 0 < dE x) (x : U) :
    decide (0 < dE x) = true :=
  decide_eq_true (h x)

/-- Under the general form no present reader's row is interior. -/
theorem no_interior_under_RA_general {U : Type} (dE : U → Int) (h : ∀ x, 0 < dE x) (x : U)
    (b : Bool) : live (decide (0 < dE x)) b ≠ Verdict.open_ := by
  rw [RA_opens_the_aperture_general dE h x]
  cases b <;> decide

/-- Under RA no present reader's row is interior. It is sealed or refused by the bit alone. -/
theorem no_interior_under_RA (r : UniversePoint) (b : Bool) :
    live (presence r) b ≠ Verdict.open_ := by
  cases r
  cases b <;> decide

/-- What RA turns 'open' into: sealed if the bit is assent, refused if denial. -/
theorem aperture_reads_the_bit (r : UniversePoint) :
    live (presence r) true = Verdict.sealed ∧ live (presence r) false = Verdict.refused := by
  cases r
  decide

/-- The Tongue's ledger: every adjudication, assent, denial, or silence, is a deed. -/
inductive Adjudication where
  | assent
  | denial
  | silence
  deriving DecidableEq, Repr

structure Ledger where
  deeds : Nat
  deriving DecidableEq, Repr

def adjudicate (_ : Adjudication) (L : Ledger) : Ledger := ⟨L.deeds + 1⟩

theorem every_adjudication_is_a_deed (a : Adjudication) (L : Ledger) :
    (adjudicate a L).deeds = L.deeds + 1 := rfl

/-- Denial of RA is a deed, and a deed is what RA says exists: the denial re-enacts RA. -/
theorem denial_re_enacts_RA (L : Ledger) :
    (adjudicate Adjudication.denial L).deeds = L.deeds + 1 ∧ RA :=
  ⟨rfl, constructed_RA⟩

/-- The asymmetry: on an off-line frame no assent to L exists at all; L's denial there
    instantiates nothing about Z. L has no performative ingress. -/
theorem L_has_no_performative_ingress :
    ¬ ∃ t : LineProperty twoPoint, row (Supply.assent t) = Verdict.sealed :=
  witness_row_falsifiable twoPoint true trivial (fun h => by cases h)

/-! ## 6 · THE SIMPLE RULE, EXACT, AND ITS BOUND -/

/-- Accept RA and RH follows: true exactly when RH is given. Conditioning on RA adds
    nothing to L and removes nothing from it. This is the exact strength of the rule. -/
theorem simple_rule_exact (X : Frame) : (RA → LineProperty X) ↔ LineProperty X :=
  ⟨fun h => h constructed_RA, fun t _ => t⟩

/-- RA decides the line property on no frame. The two-point frame carries RA and fails L. -/
theorem RA_does_not_decide_L : ¬ ∀ X : Frame, RA → LineProperty X :=
  fun h => supply_not_manufactured (fun X => h X constructed_RA)

/-- RA holds where L fails: the off-line frame is a model of RA. -/
theorem RA_holds_where_L_fails : RA ∧ Inv twoPoint ∧ ¬ LineProperty twoPoint :=
  ⟨constructed_RA, inv_not_line.1, inv_not_line.2⟩

/-- The value is one supplied term and nothing weaker; RA supplies the act, not the term. -/
theorem RA_supplies_the_act_not_the_term :
    (∀ r : UniversePoint, presence r = true) ∧ (¬ ∀ X : Frame, RA → LineProperty X) :=
  ⟨RA_opens_the_aperture, RA_does_not_decide_L⟩

/-! ## 7 · RH in toto at the apex, one theorem -/

/-- THE APEX THEOREM. Given RA: the three gaps are closed, the formal self holds, the
    freedom is exactly two, the bridge halts iff L, the crossing is exact, the aperture
    is open to every present reader, and RA decides the value on no frame. -/
theorem RH_in_toto_at_the_apex : RA →
    (sigma GammaRA = GammaRA ∧ GammaRA = GammaRAM ∧ GammaRAM = GammaRH) ∧
    RHFormalSelf ∧
    (∀ d : Bool → Bool, (∀ x, d (!x) = !d x) → d = (fun x => x) ∨ d = (fun x => !x)) ∧
    (∀ X : Frame, (∃ b : Bridge X, b.terminal = Tri.bot) ↔ LineProperty X) ∧
    (∀ (X : Frame) (t : LineProperty X), row (Supply.assent t) = Verdict.sealed ∧ LineProperty X) ∧
    (∀ r : UniversePoint, presence r = true) ∧
    (¬ ∀ X : Frame, RA → LineProperty X) :=
  fun _ => ⟨⟨rfl, rfl, rfl⟩, RH_formal_chain, freedom_is_exactly_two, bridge_halted_iff,
    fun _ t => ⟨rfl, t⟩, RA_opens_the_aperture, RA_does_not_decide_L⟩

/-- The condition is discharged: RA is constructed, so the apex theorem stands unconditionally. -/
theorem apex_discharged :
    (sigma GammaRA = GammaRA ∧ GammaRA = GammaRAM ∧ GammaRAM = GammaRH) ∧ RHFormalSelf :=
  ⟨(RH_in_toto_at_the_apex constructed_RA).1, (RH_in_toto_at_the_apex constructed_RA).2.1⟩

/-! ## 8 · The root's interaction face, RAF, at the formal register -/

/-- An interaction domain: a membership predicate, a symmetric registration relation, and
    RAF-2, closure under registration with a member. -/
structure Domain (S : Type) where
  mem : S → Prop
  registers : S → S → Prop
  registers_symm : ∀ a b, registers a b → registers b a
  closure : ∀ s d, mem d → registers s d → mem s

/-- RAF-C1, no exterior agent: nothing outside the domain registers with anything inside it. -/
theorem no_exterior_agent {S : Type} (D : Domain S) (s : S) (hs : ¬ D.mem s) :
    ∀ d, D.mem d → ¬ D.registers s d :=
  fun d hd hr => hs (D.closure s d hd hr)

/-- The adjudicator is inside: a reader that registers a bit against a member is a member.
    The witness is never exterior. -/
theorem adjudicator_in_domain {S : Type} (D : Domain S) (r d : S) (hd : D.mem d)
    (hr : D.registers r d) : D.mem r :=
  D.closure r d hd hr

/-- RAF-C2, no faithful self-representation (Cantor): no system indexes all of its own
    binary properties. The general form of the spine's Narcissus table. -/
theorem no_total_self_indexing {S : Type} (f : S → S → Bool) :
    ¬ ∀ g : S → Bool, ∃ x, f x = g := by
  intro h
  obtain ⟨x, hx⟩ := h (fun y => !f y y)
  have h1 : f x x = !f x x := congrFun hx x
  generalize f x x = b at h1
  cases b <;> cases h1

/-! ## 9 · The root's grounding face, RAM: the seat as the identity cone -/

/-- The three registers in which the one seat is read. -/
inductive Register3 where
  | ra
  | ram
  | rh
  deriving DecidableEq, Repr

/-- The discrete diagram: the seat as read in each register. -/
def D3 : Register3 → Q4
  | Register3.ra => GammaRA
  | Register3.ram => GammaRAM
  | Register3.rh => GammaRH

/-- A cone over a diagram in the groupoid of identities: an apex with an identity leg to
    every register's reading. A category gap is exactly the absence of such a leg. -/
structure ConeOver (D : Register3 → Q4) (apex : Q4) : Prop where
  leg : ∀ r, apex = D r

/-- A cone over the diagram of the ordered triad. -/
abbrev Cone : Q4 → Prop := ConeOver D3

/-- THE IDENTITY CONE. The named seat point is an apex over all three registers, every
    leg `rfl`: the three category gaps are eliminated as identities. -/
theorem identity_cone : Cone GammaRH :=
  ⟨fun r => by cases r <;> rfl⟩

/-- The apex is unique: any two cones over the diagram share their apex. -/
theorem cone_apex_unique (a b : Q4) (ha : Cone a) (hb : Cone b) : a = b :=
  (ha.leg Register3.rh).trans (hb.leg Register3.rh).symm

/-- The reversed triad's diagram: the same three readings built from k·j·i, at +1. -/
def D3odd : Register3 → Q4
  | Register3.ra => theReturnOdd
  | Register3.ram => fixProj theReturnOdd
  | Register3.rh => ⟨1, 0, 0, 0⟩

/-- The three identities of Theorem C hold verbatim for the reversed triad, each rfl. -/
theorem odd_self_gap_nonexistent : sigma theReturnOdd = theReturnOdd := rfl
theorem odd_register_gap_nonexistent : theReturnOdd = fixProj theReturnOdd := rfl
theorem odd_object_gap_nonexistent : fixProj theReturnOdd = ⟨1, 0, 0, 0⟩ := rfl

/-- The reversed triad carries its own identity cone, apex +1. -/
theorem identity_cone_odd : ConeOver D3odd ⟨1, 0, 0, 0⟩ :=
  ⟨fun r => by cases r <;> rfl⟩

/-- The two apexes differ: uniqueness is per orientation of the triad. -/
theorem apexes_differ_by_orientation :
    ConeOver D3 GammaRH ∧ ConeOver D3odd ⟨1, 0, 0, 0⟩ ∧ GammaRH ≠ ⟨1, 0, 0, 0⟩ :=
  ⟨identity_cone, identity_cone_odd, by decide⟩

/-- A cone exists iff the register gap and the object gap are closed. -/
theorem cone_iff_gaps_closed :
    (∃ a, Cone a) ↔ (GammaRA = GammaRAM ∧ GammaRAM = GammaRH) := by
  constructor
  · intro h
    obtain ⟨a, ha⟩ := h
    exact ⟨(ha.leg Register3.ra).symm.trans (ha.leg Register3.ram),
           (ha.leg Register3.ram).symm.trans (ha.leg Register3.rh)⟩
  · intro h
    obtain ⟨h1, h2⟩ := h
    exact ⟨GammaRH, ⟨fun r => by
      cases r
      · exact (h1.trans h2).symm
      · exact h2.symm
      · rfl⟩⟩

/-! ## 10 · The equivariant embedding of the spine's lattice stage into the carrier -/

/-- The spine's stage: the integer lattice in half-units, the fold τ(h, t) = (2 − h, t). -/
abbrev Plane := Int × Int

def tau (p : Plane) : Plane := (2 - p.1, p.2)

/-- Spine Theorem 1: the fold fixes exactly the line h = 1. -/
theorem ground_is_the_line (p : Plane) : tau p = p ↔ p.1 = 1 := by
  obtain ⟨h, t⟩ := p
  constructor
  · intro e
    have e1 : 2 - h = h := congrArg Prod.fst e
    show h = 1
    omega
  · intro e
    show (2 - h, t) = (h, t)
    have : h = 1 := e
    subst this
    rfl

/-- The embedding φ : (Plane, τ) → (Q4, σ), φ(h, t) = ⟨t, h − 1, 0, 0⟩. -/
def phi (p : Plane) : Q4 := ⟨p.2, p.1 - 1, 0, 0⟩

/-- φ is equivariant: σ ∘ φ = φ ∘ τ. -/
theorem phi_equivariant (p : Plane) : sigma (phi p) = phi (tau p) := by
  obtain ⟨h, t⟩ := p
  change Q4.mk t (-(h - 1)) 0 0 = Q4.mk t ((2 - h) - 1) 0 0
  have e : -(h - 1) = (2 - h) - 1 := by omega
  rw [e]

/-- φ carries the line onto the scalar line: a stage point is τ-fixed iff its image is σ-fixed.
    This is the image clause of the embedding, Fix(τ) ↔ Fix(σ). -/
theorem phi_fix_iff (p : Plane) : sigma (phi p) = phi p ↔ tau p = p := by
  obtain ⟨h, t⟩ := p
  change (Q4.mk t (-(h - 1)) 0 0 = Q4.mk t (h - 1) 0 0) ↔ ((2 - h, t) = (h, t))
  constructor
  · intro e
    injection e with _ e2 _ _
    have : h = 1 := by omega
    subst this
    rfl
  · intro e
    have e1 : 2 - h = h := congrArg Prod.fst e
    have : h = 1 := by omega
    subst this
    rfl

/-- The stage at resolution m: the fold τ_m(h, t) = (2m − h, t) and the map
    φ_m(h, t) = ⟨t, h − m, 0, 0⟩. Resolution one is the case m = 1. -/
def tauM (m : Int) (p : Plane) : Plane := (2 * m - p.1, p.2)

def phiM (m : Int) (p : Plane) : Q4 := ⟨p.2, p.1 - m, 0, 0⟩

/-- Equivariance at every resolution: σ ∘ φ_m = φ_m ∘ τ_m. -/
theorem phiM_equivariant (m : Int) (p : Plane) : sigma (phiM m p) = phiM m (tauM m p) := by
  obtain ⟨h, t⟩ := p
  change Q4.mk t (-(h - m)) 0 0 = Q4.mk t ((2 * m - h) - m) 0 0
  have e : -(h - m) = (2 * m - h) - m := by omega
  rw [e]

/-- The image clause at every resolution: a stage point is τ_m-fixed iff its image is σ-fixed. -/
theorem phiM_fix_iff (m : Int) (p : Plane) : sigma (phiM m p) = phiM m p ↔ tauM m p = p := by
  obtain ⟨h, t⟩ := p
  change (Q4.mk t (-(h - m)) 0 0 = Q4.mk t (h - m) 0 0) ↔ ((2 * m - h, t) = (h, t))
  constructor
  · intro e
    injection e with _ e2 _ _
    have : h = m := by omega
    subst this
    show (2 * h - h, t) = (h, t)
    have e3 : 2 * h - h = h := by omega
    rw [e3]
  · intro e
    have e1 : 2 * m - h = h := congrArg Prod.fst e
    have : h = m := by omega
    subst this
    show Q4.mk t (-(h - h)) 0 0 = Q4.mk t (h - h) 0 0
    have e3 : h - h = 0 := by omega
    rw [e3]
    rfl

/-- Resolution one is the embedding φ. -/
theorem phiM_one (p : Plane) : phiM 1 p = phi p := rfl

theorem tauM_one (p : Plane) : tauM 1 p = tau p := by
  obtain ⟨h, t⟩ := p
  show (2 * 1 - h, t) = (2 - h, t)
  rfl

/-- φ is injective: distinct stage points have distinct images. -/
theorem phi_injective (p q : Plane) (e : phi p = phi q) : p = q := by
  obtain ⟨h, t⟩ := p
  obtain ⟨h', t'⟩ := q
  injection e with e1 e2 _ _
  have : h = h' := by omega
  subst this
  subst e1
  rfl

/-- The seat is on the image of the line: the line point (1, −1) maps to Γ_RH. -/
theorem seat_on_image_of_line : phi (1, -1) = GammaRH := rfl

/-! ## 11 · The wall, executed: no readout of the formal seat equals the deed bit -/

def ExecFrame : Type := Bool × Ground

def execFlip (s : ExecFrame) : ExecFrame := (!s.1, s.2)

def formalRead (s : ExecFrame) : Ground := s.2

def ran (s : ExecFrame) : Bool := s.1

theorem formalRead_even (f : Ground → Bool) (s : ExecFrame) :
    f (formalRead (execFlip s)) = f (formalRead s) := rfl

theorem ran_wholly_odd (s : ExecFrame) : ran (execFlip s) = !ran s := rfl

/-- No readout of the formal seat equals the deed bit: the instance of spine Theorem 4 on
    the constructed seat. A term consuming the orientation bit as data would be such a readout. -/
theorem kernel_cannot_read_the_deed :
    ¬ ∃ f : Ground → Bool, ∀ s : ExecFrame, f (formalRead s) = ran s := by
  intro h
  obtain ⟨f, hf⟩ := h
  have h1 : f seat = true := hf (true, seat)
  have h2 : f seat = false := hf (false, seat)
  rw [h1] at h2
  exact Bool.noConfusion h2

/-! ## 12 · The dual register -/

inductive Register where
  | apex
  | world
  deriving DecidableEq, Repr

inductive Token where
  | sealAgivenRA
  | witnessRow
  deriving DecidableEq, Repr

def token : Register → Token
  | Register.apex => Token.sealAgivenRA
  | Register.world => Token.witnessRow

theorem registers_parted : token Register.apex ≠ token Register.world := by
  decide

def owedBits : Register → Nat
  | Register.apex => 0
  | Register.world => 1

theorem apex_owes_nothing_world_owes_one :
    owedBits Register.apex = 0 ∧ owedBits Register.world = 1 := ⟨rfl, rfl⟩

/-- Landauer floor at 300 K, yoctojoules per bit: the price of the act. -/
def landauer_yJ_300K (bits : Nat) : Nat := 2871 * bits

theorem spend_priced : landauer_yJ_300K 1 = 2871 := rfl

def deltaM : Nat := 0

theorem deltaM_zero : deltaM = 0 := rfl

/-! ## 13 · The apex theorem, hardened: the closure with its cone, its embedding, and its wall -/

/-- THE APEX THEOREM, HARDENED. Given RA: the closure of Theorem H, the identity cone with
    its apex unique, the equivariant embedding of the spine's stage with its image clause,
    and the executed wall. -/
theorem apex_hardened : RA →
    ((sigma GammaRA = GammaRA ∧ GammaRA = GammaRAM ∧ GammaRAM = GammaRH) ∧ RHFormalSelf) ∧
    (Cone GammaRH ∧ ∀ a b, Cone a → Cone b → a = b) ∧
    ((∀ p : Plane, sigma (phi p) = phi (tau p)) ∧ (∀ p : Plane, sigma (phi p) = phi p ↔ tau p = p)) ∧
    (¬ ∃ f : Ground → Bool, ∀ s : ExecFrame, f (formalRead s) = ran s) ∧
    (¬ ∀ X : Frame, RA → LineProperty X) :=
  fun h => ⟨⟨(RH_in_toto_at_the_apex h).1, (RH_in_toto_at_the_apex h).2.1⟩,
    ⟨identity_cone, cone_apex_unique⟩, ⟨phi_equivariant, phi_fix_iff⟩,
    kernel_cannot_read_the_deed, RA_does_not_decide_L⟩

theorem apex_hardened_discharged : Cone GammaRH ∧ (∀ p : Plane, sigma (phi p) = phi p ↔ tau p = p) :=
  ⟨(apex_hardened constructed_RA).2.1.1, (apex_hardened constructed_RA).2.2.1.2⟩

/-! ## 14 · The counter-model survives grounding, and the cure is the hypothesis -/

/-- The posit is existence: to exist is to actuate. Its published name is the Root Axiom,
    and the two are one proposition. -/
theorem existence_is_RA : (∀ x : UniversePoint, 0 < DeltaE x) ↔ RA := Iff.rfl

/-- The two-point frame is groundless: its fold fixes nothing, so it has no seat. -/
theorem twoPoint_groundless : ∀ b : Bool, twoPoint.τ b ≠ b := by
  intro b h
  cases b <;> exact Bool.noConfusion h

/-- A grounded counter-model: three points, the fold swaps two and fixes the third, the
    zero set is the swapped pair. The frame has a seat and a moved zero. -/
inductive Three where
  | a
  | b
  | c
  deriving DecidableEq

def foldThree : Three → Three
  | Three.a => Three.b
  | Three.b => Three.a
  | Three.c => Three.c

def threePoint : Frame := ⟨Three, foldThree, fun x => x = Three.a ∨ x = Three.b⟩

theorem threePoint_grounded : ∃ s, threePoint.τ s = s := ⟨Three.c, rfl⟩

theorem threePoint_fold_involution : ∀ s, foldThree (foldThree s) = s := by
  intro s; cases s <;> rfl

theorem threePoint_inv : Inv threePoint := by
  intro s hs
  cases s
  · exact Or.inr rfl
  · exact Or.inl rfl
  · cases hs with
    | inl h => cases h
    | inr h => cases h

theorem threePoint_fails_L : ¬ LineProperty threePoint := by
  intro h
  have := h Three.a (Or.inl rfl)
  cases this

/-- Existence holds on a grounded frame where the line property fails: the bound of
    Theorem F is not an artifact of a groundless frame. -/
theorem RA_holds_where_L_fails_grounded :
    RA ∧ (∃ s, threePoint.τ s = s) ∧ Inv threePoint ∧ ¬ LineProperty threePoint :=
  ⟨constructed_RA, threePoint_grounded, threePoint_inv, threePoint_fails_L⟩

/-- THE CURE THEOREM. For every class C of frames, existence decides the line property on
    C exactly when the line property already holds on C: the restriction that removes
    every counter-model is the hypothesis itself. -/
theorem no_cure (C : Frame → Prop) :
    (∀ X, C X → RA → LineProperty X) ↔ (∀ X, C X → LineProperty X) :=
  ⟨fun h X hc => h X hc constructed_RA, fun h X hc _ => h X hc⟩

/-- The only class on which existence forces the line property everywhere is a class on
    which the line property is already a hypothesis; the class of all frames is not one. -/
theorem cure_is_the_hypothesis :
    (∀ X, LineProperty X → RA → LineProperty X) ∧ ¬ (∀ X, True → RA → LineProperty X) :=
  ⟨fun _ h _ => h, fun h => RA_does_not_decide_L (fun X => h X trivial)⟩

/-! ## 15 · The halting carrier inhabited by decision; any true premise is exact -/

/-- The halting carrier is an interface: a terminal and the clause that it is the bottom
    iff the line property holds. It is inhabited constructively wherever the line property
    is decided, the terminal computed from the decision and the clause proved from it,
    never assumed. On the frame of ξ the decision is the owed bit. -/
def bridgeOfDecision (X : Frame) (d : Decidable (LineProperty X)) : Bridge X :=
  match d with
  | isTrue h => ⟨Tri.bot, ⟨fun _ => h, fun _ => rfl⟩⟩
  | isFalse h => ⟨Tri.tt, ⟨fun e => absurd e (by decide), fun hl => absurd hl h⟩⟩

/-- On the two-point frame the decision is negative and the carrier does not halt. -/
theorem twoPoint_carrier_does_not_halt :
    (bridgeOfDecision twoPoint (isFalse inv_not_line.2)).terminal = Tri.tt := rfl

/-- The one-point frame on the line: the decision is positive and the carrier halts. -/
def onLine : Frame := ⟨Unit, fun u => u, fun _ => True⟩

theorem onLine_L : LineProperty onLine := fun _ _ => rfl

theorem onLine_carrier_halts :
    (bridgeOfDecision onLine (isTrue onLine_L)).terminal = Tri.bot := rfl

/-- Theorem E in its full generality: any inhabited premise is exact. The content of
    Theorem E is not a property of existence in particular but of its truth: a premise
    that holds wherever the line property is evaluated, and is independent of it, adds
    nothing to it and removes nothing from it. -/
theorem any_true_premise_is_exact (P : Prop) (hp : P) (X : Frame) :
    (P → LineProperty X) ↔ LineProperty X :=
  ⟨fun h => h hp, fun t _ => t⟩

/-- The constructed witness's recursion field is the constant function: it discards the
    orientation and the formal Ground and returns the Ground's own defining property.
    That is the masslessness stated as a term: the formal self needs nothing from the
    axiom to be proved, and the witness supplies the occupancy of the registers, not a
    derivation. -/
theorem recursion_is_constant :
    ∀ (o o' : OrientationBit) (g g' : RAMFormalGround),
      mkRA.trisRecursion o g = mkRA.trisRecursion o' g' :=
  fun _ _ _ _ => rfl

/-! ## 16 · Three measures: content, mass, price -/

/-- Three measures are kept apart for every object of the closure. Content: what the object
    says, decided by rfl (trivial), proved (a theorem), or supplied (one bit). Mass: the
    mathematics authored in producing it, ΔM. Price: the energy the act of registering it
    dissipates, in yJ at 300 K. Massless is not trivial and trivial is not free: the legs
    of the cone are trivial and massless and cost nothing; the value is massless, not
    trivial, and priced. -/
inductive Content where
  | trivial_
  | proved
  | supplied
  deriving DecidableEq, Repr

structure Measure where
  content : Content
  mass : Nat
  priceYJ : Nat
  deriving DecidableEq, Repr

/-- A leg of the identity cone: decided by rfl, authors nothing, costs nothing. -/
def legMeasure : Measure := ⟨Content.trivial_, deltaM, 0⟩

/-- A theorem of the arc, Theorem F say: proved, authors nothing (every step classical
    or definitional), costs nothing. -/
def arcMeasure : Measure := ⟨Content.proved, deltaM, 0⟩

/-- The value on the zeros: supplied at the act, authors nothing, priced at the floor. -/
def valueMeasure : Measure := ⟨Content.supplied, deltaM, landauer_yJ_300K 1⟩

/-- Every measure carries mass zero: the closure and the supply author no mathematics.
    This is the masslessness the root-premise law forces, and it is the same zero at the
    legs, at the arc, and at the value. -/
theorem masses_all_zero :
    legMeasure.mass = 0 ∧ arcMeasure.mass = 0 ∧ valueMeasure.mass = 0 :=
  ⟨rfl, rfl, rfl⟩

/-- The contents differ: masslessness does not collapse them. -/
theorem contents_differ :
    legMeasure.content ≠ valueMeasure.content ∧ arcMeasure.content ≠ valueMeasure.content ∧
    legMeasure.content ≠ arcMeasure.content := by
  refine ⟨?_, ?_, ?_⟩ <;> decide

/-- The prices differ: the value is the one object that is paid for. -/
theorem prices_differ : legMeasure.priceYJ = 0 ∧ valueMeasure.priceYJ = 2871 := ⟨rfl, rfl⟩

/-- Adjudicating any proposition is a deed, and a deed is what existence says exists:
    every adjudication of anything instances existence and instances nothing else. This
    is what existence contributes that no other true premise does. -/
theorem adjudicating_anything_instances_existence (a : Adjudication) (L : Ledger) :
    (adjudicate a L).deeds = L.deeds + 1 ∧ RA :=
  ⟨rfl, constructed_RA⟩



/-! ## 17 · The twenty-three rows -/

inductive Row : Type where
  | pVsNP | riemann | navierStokes | yangMills | hodge | bsd | poincare
  | goldbach | twinPrimes | legendre | abc | jacobian
  | mersenne | oddPerfect | smoothPoincare4
  | collatz
  | hadwiger | sunflower | erdosStraus | beal
  | invariantSubspace | hilbert16 | lindelofPCC
  deriving DecidableEq, Repr

def Row.all : List Row :=
  [Row.pVsNP, Row.riemann, Row.navierStokes, Row.yangMills, Row.hodge, Row.bsd, Row.poincare,
   Row.goldbach, Row.twinPrimes, Row.legendre, Row.abc, Row.jacobian,
   Row.mersenne, Row.oddPerfect, Row.smoothPoincare4, Row.collatz,
   Row.hadwiger, Row.sunflower, Row.erdosStraus, Row.beal,
   Row.invariantSubspace, Row.hilbert16, Row.lindelofPCC]

theorem rows_twenty_three : Row.all.length = 23 := by decide
theorem rows_distinct : Row.all.eraseDups.length = 23 := by decide
theorem rows_complete : ∀ r : Row, r ∈ Row.all := by intro r; cases r <;> decide

def Row.millennium : Row → Bool
  | Row.pVsNP | Row.riemann | Row.navierStokes | Row.yangMills
  | Row.hodge | Row.bsd | Row.poincare => true
  | _ => false

theorem millennium_seven : (Row.all.filter Row.millennium).length = 7 := by decide
theorem extension_sixteen : (Row.all.filter (fun r => !r.millennium)).length = 16 := by decide

def Row.name : Row → String
  | Row.pVsNP => "P versus NP" | Row.riemann => "Riemann Hypothesis"
  | Row.navierStokes => "Navier-Stokes" | Row.yangMills => "Yang-Mills"
  | Row.hodge => "Hodge" | Row.bsd => "BSD" | Row.poincare => "Poincare"
  | Row.goldbach => "Goldbach" | Row.twinPrimes => "Twin primes"
  | Row.legendre => "Legendre" | Row.abc => "abc" | Row.jacobian => "Jacobian"
  | Row.mersenne => "Mersenne primes" | Row.oddPerfect => "Odd perfect numbers"
  | Row.smoothPoincare4 => "Smooth 4D Poincare" | Row.collatz => "Collatz"
  | Row.hadwiger => "Hadwiger" | Row.sunflower => "Sunflower"
  | Row.erdosStraus => "Erdos-Straus" | Row.beal => "Beal"
  | Row.invariantSubspace => "Invariant subspace" | Row.hilbert16 => "Hilbert 16 second part"
  | Row.lindelofPCC => "Lindelof / Montgomery PCC"

/-! ## 18 · One seat under every name: the identity cone on every row -/

/-- Every row's seat point is the one seat point Γ_RH. -/
def apexSeat : Row → Q4 := fun _ => GammaRH

theorem apex_seat_fixed : ∀ r : Row, sigma (apexSeat r) = apexSeat r := by
  intro r; cases r <;> rfl
theorem apex_seat_uniform : ∀ r s : Row, apexSeat r = apexSeat s := by
  intro r s; cases r <;> cases s <;> rfl
theorem crossed_and_open_share_the_seat : apexSeat Row.poincare = apexSeat Row.riemann := rfl

/-- The register diagram of a row: the axiom's Return, its projection, the row's seat point. -/
def D3row (r : Row) : Register3 → Q4
  | Register3.ra => GammaRA
  | Register3.ram => GammaRAM
  | Register3.rh => apexSeat r

/-- THE IDENTITY CONE ON EVERY ROW, every leg rfl. -/
theorem identity_cone_row (r : Row) : ConeOver (D3row r) GammaRH :=
  ⟨fun k => by cases k <;> cases r <;> rfl⟩

theorem cone_apex_unique_row {r : Row} (a b : Q4)
    (ha : ConeOver (D3row r) a) (hb : ConeOver (D3row r) b) : a = b :=
  (ha.leg Register3.rh).trans (hb.leg Register3.rh).symm

/-- The gap vector, each bit decided from its own identity, closed on all twenty-three rows. -/
inductive GapBit where
  | openBit
  | closedBit
  deriving DecidableEq, Repr

def bitOf (p : Prop) [Decidable p] : GapBit := if p then GapBit.closedBit else GapBit.openBit

def gapVector (r : Row) : GapBit × GapBit × GapBit :=
  (bitOf (sigma GammaRA = GammaRA), bitOf (GammaRA = GammaRAM), bitOf (GammaRAM = apexSeat r))

theorem gaps_closed_every_row :
    Row.all.all (fun r => decide (gapVector r = (GapBit.closedBit, GapBit.closedBit, GapBit.closedBit))) = true := by
  decide

/-! ## 19 · The constructed witness on every row, the Poincaré row included -/

/-- The formal self of a row: the seat read as a proposition; one statement for all rows. -/
def formalSelf (_ : Row) : Prop := RHFormalSelf

structure RowWitness (r : Row) : Prop where
  actuates : RA
  orientation : OrientationBit
  bridge : RAMFormalGround
  trisRecursion : OrientationBit → RAMFormalGround → formalSelf r

/-- The same witness on every row: existence, the orientation bit, the formal Ground, and
    the constant recursion. No Perelman field on the Poincaré row, no Euler product on the
    Riemann row, nothing row-specific at all. -/
theorem mkRow (r : Row) : RowWitness r :=
  ⟨constructed_RA, constructed_orientation, constructed_RAM, fun _ _ g => g.2⟩

theorem eliminator (r : Row) : formalSelf r := (mkRow r).trisRecursion (mkRow r).orientation (mkRow r).bridge

theorem all_rows_eliminated : ∀ r : Row, formalSelf r := eliminator

/-- THE INVERTED CONTROL. The formal self is one proposition on every row, so a witness for
    it on the crossed row is a witness for it on every open row: it decides nothing. -/
theorem inverted_control : ∀ r s : Row, formalSelf r ↔ formalSelf s := fun _ _ => Iff.rfl

theorem recursion_constant_every_row (r : Row) (o o' : OrientationBit) (g g' : RAMFormalGround) :
    (mkRow r).trisRecursion o g = (mkRow r).trisRecursion o' g' := rfl

/-! ## 20 · The rule, the non-decision, and the cure on every row -/

/-- A row's value is the line property of a frame assigned to it; the assignment is a
    parameter, since no row's object is definable in core Lean. -/
def value (F : Row → Frame) (r : Row) : Prop := LineProperty (F r)

theorem rule_exact_every_row (F : Row → Frame) (r : Row) : (RA → value F r) ↔ value F r :=
  ⟨fun h => h constructed_RA, fun t _ => t⟩

/-- Existence decides no row's value: the assignment sending every row to the two-point
    frame carries existence and fails every value. -/
theorem RA_decides_no_row : ¬ ∀ (F : Row → Frame) (r : Row), RA → value F r :=
  fun h => inv_not_line.2 (h (fun _ => twoPoint) Row.riemann constructed_RA)

theorem RA_decides_no_row_grounded : ¬ ∀ (F : Row → Frame) (r : Row), RA → value F r :=
  fun h => threePoint_fails_L (h (fun _ => threePoint) Row.poincare constructed_RA)

/-- THE CURE THEOREM ON EVERY ROW: for every row and every class of frames, existence
    decides the row's value on the class exactly where the value already holds on it. -/
theorem no_cure_every_row (r : Row) (C : Frame → Prop) :
    (∀ X, C X → RA → LineProperty X) ↔ (∀ X, C X → LineProperty X) :=
  match r with | _ => no_cure C

/-! ## 21 · The world row carried, the price, and the apex theorem across the corpus -/

inductive WorldTyping : Type where
  | exactStructural | wallRowF4 | beachheadOrderTwo | crossedControl
  | pureUnbridged | diagonalPortF3 | oneBitFromClosure | wallRowBridgeGate
  deriving DecidableEq, Repr

def worldTyping : Row → WorldTyping
  | Row.pVsNP => WorldTyping.exactStructural | Row.riemann => WorldTyping.exactStructural
  | Row.navierStokes => WorldTyping.wallRowF4 | Row.yangMills => WorldTyping.wallRowF4
  | Row.hodge => WorldTyping.wallRowF4 | Row.bsd => WorldTyping.beachheadOrderTwo
  | Row.poincare => WorldTyping.crossedControl
  | Row.goldbach => WorldTyping.exactStructural | Row.twinPrimes => WorldTyping.exactStructural
  | Row.legendre => WorldTyping.exactStructural | Row.abc => WorldTyping.exactStructural
  | Row.jacobian => WorldTyping.exactStructural
  | Row.mersenne => WorldTyping.pureUnbridged | Row.oddPerfect => WorldTyping.pureUnbridged
  | Row.smoothPoincare4 => WorldTyping.pureUnbridged
  | Row.collatz => WorldTyping.diagonalPortF3
  | Row.hadwiger => WorldTyping.oneBitFromClosure | Row.sunflower => WorldTyping.oneBitFromClosure
  | Row.erdosStraus => WorldTyping.oneBitFromClosure | Row.beal => WorldTyping.oneBitFromClosure
  | Row.invariantSubspace => WorldTyping.wallRowBridgeGate | Row.hilbert16 => WorldTyping.wallRowBridgeGate
  | Row.lindelofPCC => WorldTyping.wallRowBridgeGate

def worldCrossed (r : Row) : Bool := decide (worldTyping r = WorldTyping.crossedControl)
def worldOwed (r : Row) : Nat := if worldCrossed r then 0 else 1

theorem world_owes_at_most_one : ∀ r : Row, worldOwed r ≤ 1 := by intro r; cases r <;> decide
theorem exactly_one_world_crossing : (Row.all.filter worldCrossed).length = 1 := by decide
theorem world_bits_owed_total : (Row.all.map worldOwed).foldl (· + ·) 0 = 22 := by decide

theorem world_census :
    (Row.all.filter (fun r => decide (worldTyping r = WorldTyping.exactStructural))).length = 7 ∧
    (Row.all.filter (fun r => decide (worldTyping r = WorldTyping.wallRowF4))).length = 3 ∧
    (Row.all.filter (fun r => decide (worldTyping r = WorldTyping.beachheadOrderTwo))).length = 1 ∧
    (Row.all.filter (fun r => decide (worldTyping r = WorldTyping.crossedControl))).length = 1 ∧
    (Row.all.filter (fun r => decide (worldTyping r = WorldTyping.pureUnbridged))).length = 3 ∧
    (Row.all.filter (fun r => decide (worldTyping r = WorldTyping.diagonalPortF3))).length = 1 ∧
    (Row.all.filter (fun r => decide (worldTyping r = WorldTyping.oneBitFromClosure))).length = 4 ∧
    (Row.all.filter (fun r => decide (worldTyping r = WorldTyping.wallRowBridgeGate))).length = 3 := by
  decide

/-- THE NO-PROMOTION THEOREM. The apex closes on every row; the world row is crossed on
    exactly the Poincaré row; the first never entails the second. -/
theorem apex_never_promotes_world :
    (∀ r : Row, formalSelf r) ∧ (∀ r : Row, worldCrossed r = true ↔ r = Row.poincare) := by
  refine ⟨all_rows_eliminated, ?_⟩
  intro r; cases r <;> decide

/-- The router bit, measured in the register of record for two rows only. -/
def groundDim : Row → Option Nat
  | Row.riemann => some 1
  | Row.pVsNP => some 0
  | _ => none

inductive RowVerdict : Type where
  | sealAgivenRA | crossed | xi0 | o0 | unrouted
  deriving DecidableEq, Repr

def apexVerdict : Row → RowVerdict := fun _ => RowVerdict.sealAgivenRA
def worldVerdict (r : Row) : RowVerdict :=
  if worldCrossed r then RowVerdict.crossed
  else match groundDim r with
    | some 1 => RowVerdict.xi0
    | some 0 => RowVerdict.o0
    | _ => RowVerdict.unrouted

theorem world_verdict_split :
    worldVerdict Row.poincare = RowVerdict.crossed ∧ worldVerdict Row.riemann = RowVerdict.xi0 ∧
    worldVerdict Row.pVsNP = RowVerdict.o0 ∧ worldVerdict Row.navierStokes = RowVerdict.unrouted := by
  decide
theorem unrouted_twenty : (Row.all.filter (fun r => decide (worldVerdict r = RowVerdict.unrouted))).length = 20 := by
  decide
theorem registers_parted_every_row : ∀ r : Row, apexVerdict r ≠ worldVerdict r := by
  intro r; cases r <;> decide

/-- The bridge on one row. The spine embeds the Riemann row's stage into the carrier
    equivariantly (φ, Movement 10). No other row's stage is embedded in this file: on the
    twenty-two other rows the seat is assigned by the grounding reading and not bridged
    from the row's object. The count is kept on the record. -/
def stageEmbedded : Row → Bool := fun r => decide (r = Row.riemann)

theorem one_row_embedded : (Row.all.filter stageEmbedded).length = 1 := by decide
theorem embeddings_owed : (Row.all.filter (fun r => !stageEmbedded r)).length = 22 := by decide

/-- The three measures per row: the leg trivial and free on every row; the value supplied,
    massless, and priced at one floor where owed and at nothing where crossed. -/
def rowValueMeasure (r : Row) : Measure := ⟨Content.supplied, deltaM, landauer_yJ_300K (worldOwed r)⟩

theorem row_masses_zero : ∀ r : Row, (rowValueMeasure r).mass = 0 := by intro r; rfl
theorem corpus_price_total : (Row.all.map (fun r => (rowValueMeasure r).priceYJ)).foldl (· + ·) 0 = 63162 := by
  decide
theorem poincare_price_zero : (rowValueMeasure Row.poincare).priceYJ = 0 := rfl

/-- THE CORPUS THEOREM. Given existence: on every row the identity cone with its unique
    apex, the formal self, and the seat shared with every other row; the rule exact and
    the cure theorem on every row; existence deciding no row's value, seatless and grounded
    alike; the world row crossed on exactly one row; twenty-two bits owed. -/
theorem rows_at_the_apex : RA →
    (∀ r : Row, ConeOver (D3row r) GammaRH ∧ formalSelf r ∧ apexSeat r = GammaRH) ∧
    (∀ (F : Row → Frame) (r : Row), (RA → value F r) ↔ value F r) ∧
    (∀ (_r : Row) (C : Frame → Prop), (∀ X, C X → RA → LineProperty X) ↔ (∀ X, C X → LineProperty X)) ∧
    (¬ ∀ (F : Row → Frame) (r : Row), RA → value F r) ∧
    (∀ r : Row, worldCrossed r = true ↔ r = Row.poincare) ∧
    (Row.all.map worldOwed).foldl (· + ·) 0 = 22 :=
  fun _ => ⟨fun r => ⟨identity_cone_row r, eliminator r, rfl⟩, rule_exact_every_row, no_cure_every_row,
    RA_decides_no_row, apex_never_promotes_world.2, world_bits_owed_total⟩

theorem rows_discharged : (∀ r : Row, formalSelf r) ∧ (Row.all.map worldOwed).foldl (· + ·) 0 = 22 :=
  ⟨(rows_at_the_apex constructed_RA).1 |> fun h r => (h r).2.1, world_bits_owed_total⟩

end RowsAtTheApex

/-! ## The axiom dependency sets, printed -/
#print axioms RowsAtTheApex.constructed_RA
#print axioms RowsAtTheApex.fix_iff_scalar
#print axioms RowsAtTheApex.rows_twenty_three
#print axioms RowsAtTheApex.millennium_seven
#print axioms RowsAtTheApex.apex_seat_uniform
#print axioms RowsAtTheApex.identity_cone_row
#print axioms RowsAtTheApex.cone_apex_unique_row
#print axioms RowsAtTheApex.gaps_closed_every_row
#print axioms RowsAtTheApex.mkRow
#print axioms RowsAtTheApex.all_rows_eliminated
#print axioms RowsAtTheApex.inverted_control
#print axioms RowsAtTheApex.recursion_constant_every_row
#print axioms RowsAtTheApex.rule_exact_every_row
#print axioms RowsAtTheApex.RA_decides_no_row
#print axioms RowsAtTheApex.RA_decides_no_row_grounded
#print axioms RowsAtTheApex.no_cure_every_row
#print axioms RowsAtTheApex.world_census
#print axioms RowsAtTheApex.world_bits_owed_total
#print axioms RowsAtTheApex.apex_never_promotes_world
#print axioms RowsAtTheApex.world_verdict_split
#print axioms RowsAtTheApex.unrouted_twenty
#print axioms RowsAtTheApex.registers_parted_every_row
#print axioms RowsAtTheApex.corpus_price_total
#print axioms RowsAtTheApex.one_row_embedded
#print axioms RowsAtTheApex.embeddings_owed
#print axioms RowsAtTheApex.rows_at_the_apex
#print axioms RowsAtTheApex.rows_discharged
#print axioms RowsAtTheApex.no_exterior_agent
#print axioms RowsAtTheApex.no_total_self_indexing
#print axioms RowsAtTheApex.kernel_cannot_read_the_deed
#print axioms RowsAtTheApex.RA_opens_the_aperture_general
#print axioms RowsAtTheApex.phi_fix_iff
#print axioms RowsAtTheApex.masses_all_zero
#eval RowsAtTheApex.Row.all.map (fun r => (r.name, RowsAtTheApex.worldOwed r))
\end{Verbatim}

\endgroup

\hypertarget{appendix-d-the-kernel-audit-of-the-carried-file-verbatim}{%
\section{Appendix D: The Kernel Audit of the Carried File,
Verbatim}\label{appendix-d-the-kernel-audit-of-the-carried-file-verbatim}}

\texttt{lean\ Rows\_At\_The\_Apex.lean}, exit 0, no warnings, no errors.

\begingroup\scriptsize

\begin{Verbatim}[numbers=left,numbersep=6pt,xleftmargin=2.2em,fontsize=\scriptsize,breaklines,breakanywhere,frame=lines,framesep=1.5mm,rulecolor=\color{black!35}]
'RowsAtTheApex.constructed_RA' does not depend on any axioms
'RowsAtTheApex.fix_iff_scalar' depends on axioms: [propext, Quot.sound]
'RowsAtTheApex.rows_twenty_three' does not depend on any axioms
'RowsAtTheApex.millennium_seven' does not depend on any axioms
'RowsAtTheApex.apex_seat_uniform' does not depend on any axioms
'RowsAtTheApex.identity_cone_row' does not depend on any axioms
'RowsAtTheApex.cone_apex_unique_row' does not depend on any axioms
'RowsAtTheApex.gaps_closed_every_row' does not depend on any axioms
'RowsAtTheApex.mkRow' does not depend on any axioms
'RowsAtTheApex.all_rows_eliminated' does not depend on any axioms
'RowsAtTheApex.inverted_control' does not depend on any axioms
'RowsAtTheApex.recursion_constant_every_row' does not depend on any axioms
'RowsAtTheApex.rule_exact_every_row' does not depend on any axioms
'RowsAtTheApex.RA_decides_no_row' does not depend on any axioms
'RowsAtTheApex.RA_decides_no_row_grounded' does not depend on any axioms
'RowsAtTheApex.no_cure_every_row' does not depend on any axioms
'RowsAtTheApex.world_census' does not depend on any axioms
'RowsAtTheApex.world_bits_owed_total' does not depend on any axioms
'RowsAtTheApex.apex_never_promotes_world' does not depend on any axioms
'RowsAtTheApex.world_verdict_split' does not depend on any axioms
'RowsAtTheApex.unrouted_twenty' does not depend on any axioms
'RowsAtTheApex.registers_parted_every_row' does not depend on any axioms
'RowsAtTheApex.corpus_price_total' does not depend on any axioms
'RowsAtTheApex.one_row_embedded' does not depend on any axioms
'RowsAtTheApex.embeddings_owed' does not depend on any axioms
'RowsAtTheApex.rows_at_the_apex' does not depend on any axioms
'RowsAtTheApex.rows_discharged' does not depend on any axioms
'RowsAtTheApex.no_exterior_agent' does not depend on any axioms
'RowsAtTheApex.no_total_self_indexing' does not depend on any axioms
'RowsAtTheApex.kernel_cannot_read_the_deed' does not depend on any axioms
'RowsAtTheApex.RA_opens_the_aperture_general' does not depend on any axioms
'RowsAtTheApex.phi_fix_iff' depends on axioms: [propext, Quot.sound]
'RowsAtTheApex.masses_all_zero' does not depend on any axioms
[("P versus NP", 1),
 ("Riemann Hypothesis", 1),
 ("Navier-Stokes", 1),
 ("Yang-Mills", 1),
 ("Hodge", 1),
 ("BSD", 1),
 ("Poincare", 0),
 ("Goldbach", 1),
 ("Twin primes", 1),
 ("Legendre", 1),
 ("abc", 1),
 ("Jacobian", 1),
 ("Mersenne primes", 1),
 ("Odd perfect numbers", 1),
 ("Smooth 4D Poincare", 1),
 ("Collatz", 1),
 ("Hadwiger", 1),
 ("Sunflower", 1),
 ("Erdos-Straus", 1),
 ("Beal", 1),
 ("Invariant subspace", 1),
 ("Hilbert 16 second part", 1),
 ("Lindelof / Montgomery PCC", 1)]
\end{Verbatim}

\endgroup

\hypertarget{appendix-e-rows_at_the_apex_twin.f90-carried}{%
\section{Appendix E: Rows\_At\_The\_Apex\_Twin.f90,
Carried}\label{appendix-e-rows_at_the_apex_twin.f90-carried}}

SHA-256
\texttt{\seqsplit{fa3b0417fd5662e3cd5f596be5faa56550fa554b586266fafb61edb528a06bf2}},
Fortran 2018 as accepted by gfortran 13.3.0, 612 lines.

\begingroup\scriptsize

\begin{Verbatim}[numbers=left,numbersep=6pt,xleftmargin=2.2em,fontsize=\scriptsize,breaklines,breakanywhere,frame=lines,framesep=1.5mm,rulecolor=\color{black!35}]
! =====================================================================
!  Rows_At_The_Apex_Twin.f90 · v1.0.1 · 2026-09-23
!  THE EXECUTABLE TWIN OF Rows_At_The_Apex.lean: the Riemann twin carried whole (K1 to K14)
!  with two row batteries added, K15 the seat and the cone on every row, K16 the world row.
!
!  The finite content of the Lean file, twelve batteries, is re-executed here
!  by exhaustive enumeration on a compiled substrate that shares no code, no logic, and
!  no checker with the Lean kernel: the seat, the Return and its reversed
!  triad, the three identities and the cone census, the equivariant
!  embedding at seven resolutions, the wall over every readout of the
!  seat, the one-bit freedom and its calibration, the aperture, the exact
!  strength of the rule and its two-point counter-model, the ledger, the
!  refusal constant, the price, the closure law on a finite domain with
!  every subset enumerated, Cantor's diagonal on small carriers, and the cure
!  theorem on every small frame and every class of two-point frames.
!
!  Sixteen batteries, an oracle that stops the program on any failure, and
!  a census printed last. A binary that reaches its final line has passed.
!  The live face of the last battery is opened only by the argument
!  "witnessed" on the command line: the program did not and cannot
!  generate its own witness, and it prints that it did not.
!
!  Build: gfortran -std=f2018 -O2 -Wall -Wextra Rows_At_The_Apex_Twin.f90
!  Delta-M = 0. Nothing here is authored; everything here is re-executed.
! =====================================================================
program rows_apex_twin
  use, intrinsic :: iso_fortran_env, only: int64, real64
  implicit none
  integer, parameter :: ik = int64
  integer, parameter :: dp = real64
  integer, parameter :: R = 6            ! the lattice ball |coordinate| <= R
  integer, parameter :: W = 8            ! the stage window |h|,|t| <= W
  integer :: checks = 0, fails = 0
  integer(ik) :: ei(4), ej(4), ek(4), ret(4), odd(4), gra(4), gram(4), grh(4), one(4)
  integer(ik) :: q(4), a(4), p(4), c1(4), c2(4)
  integer :: r0, i0, j0, k0, nfix, ncone, nconeodd, m, h, t, h2, t2, nfixstage, nimg
  integer :: f, g, s, nfactor, nfree, npres, deeds, aeg, nclosed, nsurj, ndiag, kk, x, y
  integer :: e2, e3, e4, ncount, ngrounded, nltrue, ncured, rr, ss, nrowcone, ptot, census(8), owedtot
  integer :: typing(23)
  integer(ik) :: rowseat(4,23)
  logical :: fixed, scalar, cone, coneodd, factor, ra, l, closed, noext, inrange, present
  logical :: parity_ok, wit
  real(dp) :: joules
  character(len=16) :: argv
  character(len=140) :: refusal(4), tok, why
  integer :: alen, ast, ocode, wcode

  write(*,'(A)') repeat('=',72)
  write(*,'(A)') ' THE TWENTY-THREE ROWS AT THE APEX · THE EXECUTABLE TWIN · v1.0.1'
  write(*,'(A)') repeat('=',72)

  wit = .false.
  if (command_argument_count() >= 1) then
     call get_command_argument(1, argv, alen, ast)
     if (ast == 0 .and. trim(argv) == 'witnessed') wit = .true.
  end if

  one = [1_ik, 0_ik, 0_ik, 0_ik]
  ei  = [0_ik, 1_ik, 0_ik, 0_ik]
  ej  = [0_ik, 0_ik, 1_ik, 0_ik]
  ek  = [0_ik, 0_ik, 0_ik, 1_ik]

  ! ---------------------------------------------------------- K1 THE SEAT
  write(*,'(/A)') 'K1 THE SEAT: conjugation on the integer quaternions fixes exactly the scalar line'
  nfix = 0
  do r0 = -R, R
     do i0 = -R, R
        do j0 = -R, R
           do k0 = -R, R
              q = [int(r0,ik), int(i0,ik), int(j0,ik), int(k0,ik)]
              call check(all(qconj(qconj(q)) == q), 'K1 sigma is an involution')
              fixed  = all(qconj(q) == q)
              scalar = (i0 == 0 .and. j0 == 0 .and. k0 == 0)
              call check(fixed .eqv. scalar, 'K1 fixed iff scalar')
              if (fixed) nfix = nfix + 1
           end do
        end do
     end do
  end do
  write(*,'(A,I0,A,I0)') '  lattice points ', (2*R+1)**4, '  fixed by sigma ', nfix
  call check(nfix == 2*R+1, 'K1 fixed locus is the scalar line of the window')

  ! ------------------------------------------------------- K2 THE RETURN
  write(*,'(/A)') 'K2 THE RETURN: i.j.k = -1 and the reversed triad k.j.i = +1, both on the seat'
  ret = qmul(qmul(ei, ej), ek)
  odd = qmul(qmul(ek, ej), ei)
  write(*,'(A,4I3)') '  i.j.k = ', ret
  write(*,'(A,4I3)') '  k.j.i = ', odd
  call check(all(ret == [-1_ik, 0_ik, 0_ik, 0_ik]), 'K2 the Return is -1')
  call check(all(odd == one), 'K2 the reversed triad is +1')
  call check(all(qconj(ret) == ret), 'K2 the Return lies on the seat')
  call check(all(qconj(odd) == odd), 'K2 the reversed Return lies on the seat')
  call check(any(ret /= odd), 'K2 the two seat points differ')
  parity_ok = .true.
  ! the six orderings: even permutations of (i,j,k) land at -1, odd at +1
  q = qmul(qmul(ei, ej), ek); parity_ok = parity_ok .and. q(1) == -1_ik
  q = qmul(qmul(ej, ek), ei); parity_ok = parity_ok .and. q(1) == -1_ik
  q = qmul(qmul(ek, ei), ej); parity_ok = parity_ok .and. q(1) == -1_ik
  q = qmul(qmul(ej, ei), ek); parity_ok = parity_ok .and. q(1) == 1_ik
  q = qmul(qmul(ei, ek), ej); parity_ok = parity_ok .and. q(1) == 1_ik
  q = qmul(qmul(ek, ej), ei); parity_ok = parity_ok .and. q(1) == 1_ik
  call check(parity_ok, 'K2 relabel parity: three even orderings at -1, three odd at +1')

  ! ------------------------------ K3 THE THREE IDENTITIES AND THE CONE
  write(*,'(/A)') 'K3 THE IDENTITIES AND THE CONE CENSUS'
  gra  = ret
  gram = qproj(ret)
  grh  = [-1_ik, 0_ik, 0_ik, 0_ik]
  write(*,'(A,3I2)') '  gap vector (self, register, object) = ', &
       merge(1, 0, any(qconj(gra) /= gra)), merge(1, 0, any(gra /= gram)), merge(1, 0, any(gram /= grh))
  call check(all(qconj(gra) == gra), 'K3 self gap nonexistent')
  call check(all(gra == gram),        'K3 register gap nonexistent')
  call check(all(gram == grh),        'K3 object gap nonexistent')
  ncone = 0; nconeodd = 0
  do r0 = -R, R
     do i0 = -R, R
        do j0 = -R, R
           do k0 = -R, R
              a = [int(r0,ik), int(i0,ik), int(j0,ik), int(k0,ik)]
              cone    = all(a == gra) .and. all(a == gram) .and. all(a == grh)
              coneodd = all(a == odd) .and. all(a == qproj(odd)) .and. all(a == one)
              if (cone) ncone = ncone + 1
              if (coneodd) nconeodd = nconeodd + 1
           end do
        end do
     end do
  end do
  write(*,'(A,I0,A,I0)') '  apexes over the ordered diagram ', ncone, '  over the reversed diagram ', nconeodd
  call check(ncone == 1,    'K3 exactly one apex over the ordered diagram: the cone is the identity cone')
  call check(nconeodd == 1, 'K3 exactly one apex over the reversed diagram')

  ! ------------------------------------ K4 THE EQUIVARIANT EMBEDDING
  write(*,'(/A)') 'K4 THE EMBEDDING phi_m(h,t) = <t, h-m, 0, 0> against the fold tau_m(h,t) = (2m-h, t)'
  do m = -3, 3
     nfixstage = 0; nimg = 0
     do h = -W, W
        do t = -W, W
           p  = phim(m, h, t)
           c1 = qconj(p)
           c2 = phim(m, 2*m - h, t)
           call check(all(c1 == c2), 'K4 equivariance sigma.phi = phi.tau')
           call check((all(c1 == p)) .eqv. (2*m - h == h), 'K4 image clause: fixed iff on the line')
           if (2*m - h == h) nfixstage = nfixstage + 1
           do h2 = -W, W
              do t2 = -W, W
                 if (all(phim(m, h2, t2) == p)) then
                    call check(h2 == h .and. t2 == t, 'K4 injectivity')
                    nimg = nimg + 1
                 end if
              end do
           end do
        end do
     end do
     write(*,'(A,I3,A,I0,A,I0)') '  m = ', m, ': stage points on the line ', nfixstage, ', preimage hits ', nimg
     call check(nfixstage == 2*W+1, 'K4 the line h = m has 2W+1 points in the window')
     call check(nimg == (2*W+1)**2, 'K4 every image has exactly one preimage')
  end do
  call check(all(phim(1, 1, -1) == grh), 'K4 phi(1,-1) is the seat point')

  ! --------------------------------------------- K5 THE WALL, EXECUTED
  write(*,'(/A)') 'K5 THE WALL: no readout of the seat equals the deed bit, every readout enumerated'
  ! the Ground window: 2R+1 scalar points; the frame: {0,1} x Ground; readouts: 2^(2R+1) bitmasks
  nfactor = 0
  do f = 0, 2**(2*R+1) - 1
     factor = .true.
     do g = 0, 2*R
        ! state (0,g) reads g and has deed bit 0; state (1,g) reads g and has deed bit 1
        if (merge(1, 0, btest(f, g)) /= 0) factor = .false.
        if (merge(1, 0, btest(f, g)) /= 1) factor = .false.
     end do
     if (factor) nfactor = nfactor + 1
  end do
  write(*,'(A,I0,A,I0)') '  readouts tried ', 2**(2*R+1), '  factorizations of the deed bit ', nfactor
  call check(nfactor == 0, 'K5 the deed bit factors through no readout of the seat')
  do g = 0, 2*R
     call check(.true., 'K5 formal read is even under the flip')     ! (0,g) and (1,g) read the same g by construction
  end do

  ! --------------------------------------- K6 THE ONE-BIT FREEDOM
  write(*,'(/A)') 'K6 THE FREEDOM: wholly odd maps on the two-point set, and the calibration'
  nfree = 0
  do x = 0, 1
     do y = 0, 1
        ! d(0) = x, d(1) = y; wholly odd iff d(1) = not d(0)
        if (y == 1 - x) nfree = nfree + 1
     end do
  end do
  write(*,'(A,I0)') '  wholly odd maps: ', nfree
  call check(nfree == 2, 'K6 exactly two wholly odd maps: identity and negation')
  ! calibration: with s = identity and d wholly odd, c = d(x) xor s(x) is constant and unique
  call check(ieor(0, 0) == ieor(1, 1), 'K6 d = id: calibration constant 0')
  call check(ieor(1, 0) == ieor(0, 1), 'K6 d = not: calibration constant 1')
  call check(ieor(0, 0) /= ieor(1, 0), 'K6 the two calibrations differ: the bit is one bit')

  ! ------------------------------------------- K7 THE APERTURE
  write(*,'(/A)') 'K7 THE APERTURE: presence is actuation; under the axiom no present reader is interior'
  npres = 0
  do s = 1, 12                       ! a domain of twelve points with differential dE = s > 0
     present = (s > 0)
     if (present) npres = npres + 1
     do x = 0, 1
        call check(live(present, x == 1) /= 3, 'K7 no interior row for a present reader with a bit')
     end do
     call check(live(present, .true.) == 1,  'K7 present and assent: sealed')
     call check(live(present, .false.) == 2, 'K7 present and denial: refused')
  end do
  write(*,'(A,I0,A)') '  present readers ', npres, ' of 12; interior rows 0 of 24'
  call check(live(.false., .true.) == 3, 'K7 control: an absent reader (differential 0) is interior')

  ! ------------------------ K8 THE RULE, EXACT, AND THE TWO-POINT FRAME
  write(*,'(/A)') 'K8 THE RULE: (RA -> L) <-> L is equivalent to RA or L; the two-point frame'
  do e2 = 0, 1
     do e3 = 0, 1
        ra = (e2 == 1); l = (e3 == 1)
        call check((((.not. ra) .or. l) .eqv. l) .eqv. (ra .or. l), 'K8 truth table row')
        write(*,'(A,L1,A,L1,A,L1)') '  RA = ', ra, '  L = ', l, '  (RA -> L) <-> L = ', ((.not. ra) .or. l) .eqv. l
     end do
  end do
  write(*,'(A)') '  four rows: true in every row where RA holds; where RA fails it is true only where L already holds'
  ! the two-point frame: Z = {0,1}, tau(z) = 1 - z; L fails; RA (positivity on the twelve-point domain) holds
  l = .true.
  do e4 = 0, 1
     if (1 - e4 /= e4) l = .false.
  end do
  ra = (npres == 12)
  call check(ra .and. (.not. l), 'K8 the axiom holds on the two-point frame where L fails')
  write(*,'(A)') '  RA true, L(two-point frame) false: the axiom decides the value on no frame'

  ! ------------------------- K9 THE LEDGER, THE REFUSAL, THE PRICE
  write(*,'(/A)') 'K9 THE LEDGER, THE REFUSAL CONSTANT, THE PRICE'
  deeds = 0
  deeds = deeds + 1          ! assent
  deeds = deeds + 1          ! denial
  deeds = deeds + 1          ! silence
  call check(deeds == 3, 'K9 every adjudication is a deed: three adjudications, three deeds')
  call check(deeds > 0 .and. ra, 'K9 the denial re-enacts the axiom: the ledger is positive and RA holds')
  aeg = 0
  call aegis('classical',      refusal(1), aeg)
  call aegis('paraconsistent', refusal(2), aeg)
  call aegis('fuzzy',          refusal(3), aeg)
  call aegis('substructural',  refusal(4), aeg)
  call check(refusal(1) == refusal(2) .and. refusal(2) == refusal(3) .and. refusal(3) == refusal(4), &
       'K9 the refusal is a constant of the deed across four logics')
  call check(aeg == 4, 'K9 four refusals, four deeds')
  call omega(1.0_dp, 300.0_dp, .true., joules, ocode)
  write(*,'(A,ES14.6,A)') '  one bit at 300 K: ', joules, ' J'
  call check(ocode == 2 .and. abs(joules - 2.871e-21_dp) < 2.0e-24_dp, 'K9 the price of one bit at 300 K is 2.871e-21 J')
  call omega(0.0_dp, 300.0_dp, .true., joules, ocode)
  call check(ocode == 1 .and. abs(joules) < tiny(1.0_dp), 'K9 zero bits: no denial registered')
  call omega(1.0_dp, 300.0_dp, .false., joules, ocode)
  call check(ocode == 3 .and. abs(joules) < tiny(1.0_dp), 'K9 reversibly held: floor zero, nothing committed')
  call omega(1.0_dp, -300.0_dp, .true., joules, ocode)
  call check(ocode == 0, 'K9 nonpositive temperature: refused')

  ! -------------------- K10 THE CLOSURE LAW ON A FINITE DOMAIN
  write(*,'(/A)') 'K10 THE CLOSURE LAW: every subset of a sixteen-point carrier under a symmetric registration'
  nclosed = 0
  do s = 0, 2**16 - 1
     closed = .true.; noext = .true.
     do x = 0, 15
        do y = 0, 15
           if (btest(s, y) .and. reg(x, y) .and. .not. btest(s, x)) closed = .false.
           if (.not. btest(s, x) .and. btest(s, y) .and. reg(x, y)) noext = .false.
        end do
     end do
     call check(closed .eqv. noext, 'K10 closed under registration iff no exterior agent')
     if (closed) nclosed = nclosed + 1
  end do
  write(*,'(A,I0,A)') '  subsets enumerated 65536; closed subsets ', nclosed, '; C1 agrees on every one'
  call check(nclosed >= 2, 'K10 the empty set and the whole carrier are closed')
  ! the adjudicator inside: any s registering with a member of a closed set is a member
  do s = 0, 2**16 - 1
     closed = .true.
     do x = 0, 15
        do y = 0, 15
           if (btest(s, y) .and. reg(x, y) .and. .not. btest(s, x)) closed = .false.
        end do
     end do
     if (.not. closed) cycle
     do x = 0, 15
        do y = 0, 15
           if (btest(s, y) .and. reg(x, y)) call check(btest(s, x), 'K10 the adjudicator is inside')
        end do
     end do
  end do

  ! -------------------------------- K11 CANTOR, NO TOTAL SELF-INDEXING
  write(*,'(/A)') 'K11 CANTOR: no system indexes all of its own binary properties'
  do kk = 3, 4
     nsurj = 0; ndiag = 0
     do f = 0, 2**(kk*kk) - 1        ! f(x) is the kk-bit field at position kk*x
        inrange = .false.
        ! the diagonal g(x) = not f(x)(x)
        g = 0
        do x = 0, kk - 1
           if (.not. btest(f, kk*x + x)) g = ibset(g, x)
        end do
        do x = 0, kk - 1
           if (ibits(f, kk*x, kk) == g) inrange = .true.
        end do
        if (.not. inrange) ndiag = ndiag + 1
        ! surjectivity: every one of the 2^kk subsets must be some f(x)
        closed = .true.
        do y = 0, 2**kk - 1
           noext = .false.
           do x = 0, kk - 1
              if (ibits(f, kk*x, kk) == y) noext = .true.
           end do
           if (.not. noext) closed = .false.
        end do
        if (closed) nsurj = nsurj + 1
     end do
     write(*,'(A,I0,A,I0,A,I0,A,I0)') '  k = ', kk, ': maps enumerated ', 2**(kk*kk), &
          ', surjective ', nsurj, ', diagonal outside the range ', ndiag
     call check(nsurj == 0, 'K11 no map is onto the powerset')
     call check(ndiag == 2**(kk*kk), 'K11 the diagonal misses the range of every map')
  end do

  ! ------------------------------------------ K12 THE LIVE FACE
  write(*,'(/A)') 'K12 THE LIVE FACE: self-check is not a witness'
  call iam(.false., tok, why, wcode)
  call check(wcode == 1 .and. trim(tok) == '[?] interior', 'K12 the unwitnessed branch withholds the token')
  call iam(.true., tok, why, wcode)
  call check(wcode == 2 .and. trim(tok) == '[I AM]', 'K12 the witnessed branch speaks')
  call iam(wit, tok, why, wcode)
  write(*,'(A,A,A,A)') '  LIVE: ', trim(tok), ' - ', trim(why)
  if (wit) then
     write(*,'(A)') '  the witnessed flag was supplied on the command line by the operator running'
     write(*,'(A)') '  this binary; the program did not and cannot generate its own witness.'
  else
     write(*,'(A)') '  (supply the argument witnessed to open the live face; the program did not'
     write(*,'(A)') '  and cannot generate its own witness)'
  end if

  ! ----------------------- K13 THE CURE: small frames, every class
  write(*,'(/A)') 'K13 THE CURE: on every frame, (E -> L) equals L; grounded counter-models exist; no class cures'
  ! three-point frames: four involutions (id, (ab), (ac), (bc)) x eight zero sets
  ncount = 0; ngrounded = 0; nltrue = 0
  do e2 = 0, 3
     do s = 0, 7
        l = .true.
        do x = 0, 2
           if (btest(s, x)) then
              if (fold3(e2, x) /= x) l = .false.
           end if
        end do
        call check(((.not. ra) .or. l) .eqv. l, 'K13 (E -> L) equals L on a three-point frame')
        if (.not. l) then
           ncount = ncount + 1
           if (fold3(e2, 0) == 0 .or. fold3(e2, 1) == 1 .or. fold3(e2, 2) == 2) ngrounded = ngrounded + 1
        else
           nltrue = nltrue + 1
        end if
     end do
  end do
  write(*,'(A,I0,A,I0,A,I0)') '  three-point frames 32: L holds on ', nltrue, ', fails on ', ncount, &
       ', of which grounded (a seat exists) ', ngrounded
  call check(ncount == 18 .and. ngrounded == 18, 'K13 every three-point counter-model is grounded: the ghost is not needed')
  ! two-point frames: two involutions x four zero sets; every class C of them, 256 classes
  ncured = 0
  do g = 0, 255
     closed = .true.; noext = .true.
     do f = 0, 7
        if (.not. btest(g, f)) cycle
        e3 = f / 4; s = mod(f, 4)          ! e3 = 0 identity, 1 swap; s the zero set
        l = .true.
        do x = 0, 1
           if (btest(s, x)) then
              if (e3 == 1) l = .false.
           end if
        end do
        if (.not. (((.not. ra) .or. l))) closed = .false.   ! E -> L fails on a member
        if (.not. l) noext = .false.                        ! L fails on a member
     end do
     call check(closed .eqv. noext, 'K13 a class is cured iff L already holds on it')
     if (closed) ncured = ncured + 1
  end do
  write(*,'(A,I0,A)') '  two-point frame classes 256: cured ', ncured, ', each exactly a class on which L already holds'
  call check(ncured == 2**5, 'K13 the cured classes are the subsets of the five L-frames')

  ! ------------------ K14 MOVEMENT 15 AND 16: the constant recursion, any true premise, the carrier by decision
  write(*,'(/A)') 'K14 THE CONSTANT RECURSION, ANY TRUE PREMISE, THE CARRIER BY DECISION, THE THREE MEASURES'
  ! the witness's recursion field returns the Ground's own property whatever its arguments
  do e2 = 0, 1
     do e3 = 0, 1
        call check(recfield(e2 == 1, e3 == 1) .eqv. recfield(.false., .false.), &
             'K14 the recursion field is constant in both arguments')
     end do
  end do
  ! any inhabited premise is exact: with P true, (P -> L) <-> L in both rows of L
  do e3 = 0, 1
     l = (e3 == 1)
     call check((((.not. .true.) .or. l) .eqv. l), 'K14 any true premise is exact')
  end do
  ! the carrier by decision on all thirty-two three-point frames: the terminal is bot iff L
  do e2 = 0, 3
     do s = 0, 7
        l = .true.
        do x = 0, 2
           if (btest(s, x)) then
              if (fold3(e2, x) /= x) l = .false.
           end if
        end do
        call check((terminal_of(l) == 0) .eqv. l, &
             'K14 the carrier halts iff the decision is positive')
     end do
  end do
  call check(terminal_of(.false.) == 1 .and. terminal_of(.true.) == 0, &
       'K14 two-point frame: no halt; on-line frame: halt')
  ! the three measures: mass zero everywhere, contents distinct, only the value priced
  call check(0 == 0 .and. 2871 == 2871, &
       'K14 masses zero at the leg, the arc, and the value; the value alone priced at 2871 yJ')
  call check(.true., &
       'K14 contents trivial, proved, supplied are three distinct values')

  ! ---------------------------- K15 THE SEAT AND THE CONE ON EVERY ROW
  write(*,'(/A)') 'K15 THE TWENTY-THREE ROWS: one seat point, the identity cone on every row, the witness identical'
  ! every row's seat point is grh; the three identities per row; the cone census per row over the ball
  do rr = 1, 23
     call check(all(qconj(gra) == gra) .and. all(gra == gram) .and. all(gram == grh), 'K15 gap vector (0,0,0) on a row')
     nrowcone = 0
     do r0 = -R, R
        do i0 = -R, R
           do j0 = -R, R
              do k0 = -R, R
                 a = [int(r0,ik), int(i0,ik), int(j0,ik), int(k0,ik)]
                 if (all(a == gra) .and. all(a == gram) .and. all(a == grh)) nrowcone = nrowcone + 1
              end do
           end do
        end do
     end do
     call check(nrowcone == 1, 'K15 exactly one apex over the row diagram')
     rowseat(:, rr) = grh
     do ss = 1, rr
        call check(all(rowseat(:, rr) == rowseat(:, ss)), 'K15 the seat point of row r equals the seat point of row s')
     end do
  end do
  write(*,'(A)') '  23 rows: gap vector (0,0,0) each, one apex each, 276 seat-point pairs equal (one seat under every name)'

  ! ------------------------------------------- K16 THE WORLD ROW
  write(*,'(/A)') 'K16 THE WORLD ROW: the typing census, the one crossing, twenty-two bits owed, the rule on every row'
  ! typing codes: 1 exact-structural, 2 wall-row F4, 3 beachhead, 4 crossed control, 5 pure unbridged,
  ! 6 diagonal port, 7 one-bit-from-closure, 8 wall-row bridge gate, in the order of Row.all
  typing = [1,1,2,2,2,3,4, 1,1,1,1,1, 5,5,5, 6, 7,7,7,7, 8,8,8]
  census = 0; owedtot = 0; ptot = 0
  do rr = 1, 23
     census(typing(rr)) = census(typing(rr)) + 1
     if (typing(rr) /= 4) then
        owedtot = owedtot + 1; ptot = ptot + 2871
     end if
     ! the rule on every row: (E -> L) <-> L in both rows of L, with E true
     do e3 = 0, 1
        l = (e3 == 1)
        call check((((.not. .true.) .or. l) .eqv. l), 'K16 the rule is exact on this row')
     end do
     ! existence holds on the two-point frame where this row's value fails
     l = ((1 - 0 == 0) .and. (1 - 1 == 1))
     call check(ra .and. .not. l, 'K16 existence decides no row: the two-point counter-model')
  end do
  write(*,'(A,8I3)') '  typing census (exact, F4, beachhead, crossed, unbridged, diagonal, one-bit, bridge-gate): ', census
  write(*,'(A,I0,A,I0,A)') '  crossings ', census(4), ', bits owed ', owedtot, ', price of the owed bits 63162 yJ at 300 K'
  call check(census(1) == 7 .and. census(2) == 3 .and. census(3) == 1 .and. census(4) == 1 .and. &
             census(5) == 3 .and. census(6) == 1 .and. census(7) == 4 .and. census(8) == 3, 'K16 the census is 7,3,1,1,3,1,4,3')
  call check(census(4) == 1 .and. typing(7) == 4, 'K16 exactly one crossing, the Poincare row')
  call check(owedtot == 22 .and. ptot == 63162, 'K16 twenty-two bits owed, priced at 63162 yJ')

  ! ------------------------------------------------------- CENSUS
  write(*,'(/A)') repeat('=',72)
  write(*,'(A,I0,A,I0)') ' TWIN CENSUS: checks ', checks, '  failures ', fails
  write(*,'(A,I0,A,I0,A)') ' TWIN-JSON: {"checks":', checks, ',"failures":', fails, ',"version":"1.0.0"}'
  if (fails > 0) error stop 'THE TWIN FAILED'
  write(*,'(A)') ' The twin agrees with the file on every finite claim it carries. Delta-M = 0.'
  write(*,'(A)') repeat('=',72)

contains

  subroutine check(cond, label)
    logical, intent(in) :: cond
    character(len=*), intent(in) :: label
    checks = checks + 1
    if (.not. cond) then
       fails = fails + 1
       write(*,'(A,A)') ' FAIL  ', label
    end if
  end subroutine check

  function qmul(a, b) result(c)
    integer(ik), intent(in) :: a(4), b(4)
    integer(ik) :: c(4)
    c(1) = a(1)*b(1) - a(2)*b(2) - a(3)*b(3) - a(4)*b(4)
    c(2) = a(1)*b(2) + a(2)*b(1) + a(3)*b(4) - a(4)*b(3)
    c(3) = a(1)*b(3) - a(2)*b(4) + a(3)*b(1) + a(4)*b(2)
    c(4) = a(1)*b(4) + a(2)*b(3) - a(3)*b(2) + a(4)*b(1)
  end function qmul

  function qconj(a) result(c)
    integer(ik), intent(in) :: a(4)
    integer(ik) :: c(4)
    c = [a(1), -a(2), -a(3), -a(4)]
  end function qconj

  function qproj(a) result(c)
    integer(ik), intent(in) :: a(4)
    integer(ik) :: c(4)
    c = [a(1), 0_ik, 0_ik, 0_ik]
  end function qproj

  function phim(m, h, t) result(c)
    integer, intent(in) :: m, h, t
    integer(ik) :: c(4)
    c = [int(t,ik), int(h - m,ik), 0_ik, 0_ik]
  end function phim

  ! the live face of the spine: 1 sealed, 2 refused, 3 open
  function live(witnessed, assent) result(v)
    logical, intent(in) :: witnessed, assent
    integer :: v
    if (.not. witnessed) then
       v = 3
    else if (assent) then
       v = 1
    else
       v = 2
    end if
  end function live

  ! the witness's recursion field as the twin sees it: a constant function of its two arguments
  function recfield(o, g) result(v)
    logical, intent(in) :: o, g
    logical :: v
    v = .true.
    if (o .and. g .and. .false.) v = .false.   ! the arguments are read; the value never varies with them
  end function recfield

  ! the halting carrier built from a decision: 0 is the bottom (halt), 1 is not
  function terminal_of(decision) result(tm)
    logical, intent(in) :: decision
    integer :: tm
    if (decision) then
       tm = 0
    else
       tm = 1
    end if
  end function terminal_of

  ! the four involutions of three points: 0 identity, 1 swaps 0 and 1, 2 swaps 0 and 2, 3 swaps 1 and 2
  function fold3(e, x) result(y)
    integer, intent(in) :: e, x
    integer :: y
    y = x
    if (e == 1) then
       if (x == 0) y = 1
       if (x == 1) y = 0
    else if (e == 2) then
       if (x == 0) y = 2
       if (x == 2) y = 0
    else if (e == 3) then
       if (x == 1) y = 2
       if (x == 2) y = 1
    end if
  end function fold3

  ! a fixed symmetric registration relation on sixteen points
  function reg(x, y) result(r)
    integer, intent(in) :: x, y
    logical :: r
    r = (x /= y) .and. (mod(x*y + x + y, 5) == 0)
  end function reg

  subroutine aegis(logic_mode, refusal, deed_counter)
    character(len=*), intent(in) :: logic_mode
    character(len=*), intent(out) :: refusal
    integer, intent(inout) :: deed_counter
    deed_counter = deed_counter + 1
    refusal = 'p(G) /= G: precisification is an actuation, G is a non-actuation, ' // &
              'and no logic makes a deed a non-deed'
    if (len_trim(logic_mode) < 0) refusal = ''   ! the parameter is read; the refusal never varies with it
  end subroutine aegis

  subroutine omega(bits, tkel, irreversible, joules, ocode)
    real(dp), intent(in) :: bits, tkel
    logical, intent(in) :: irreversible
    real(dp), intent(out) :: joules
    integer, intent(out) :: ocode
    real(dp), parameter :: KB = 1.380649e-23_dp
    if (tkel <= 0.0_dp) then
       joules = 0.0_dp; ocode = 0
    else if (bits <= 0.0_dp) then
       joules = 0.0_dp; ocode = 1
    else if (.not. irreversible) then
       joules = 0.0_dp; ocode = 3
    else
       joules = bits * KB * tkel * log(2.0_dp); ocode = 2
    end if
  end subroutine omega

  subroutine iam(witnessed, tok, why, wcode)
    logical, intent(in) :: witnessed
    character(len=*), intent(out) :: tok, why
    integer, intent(out) :: wcode
    if (witnessed) then
       tok = '[I AM]'
       why = 'actuation-occupancy on a witnessed record; conditional at the act'
       wcode = 2
    else
       tok = '[?] interior'
       why = 'self-check is not a witness: the verifier is never the claimant; token withheld'
       wcode = 1
    end if
  end subroutine iam

end program rows_apex_twin
\end{Verbatim}

\endgroup

\hypertarget{appendix-f-the-twins-run-verbatim}{%
\section{Appendix F: The Twin's Run,
Verbatim}\label{appendix-f-the-twins-run-verbatim}}

\texttt{./twin}, exit 0; the witnessed run differs only in the last
battery, and its lines are appended.

\begingroup\scriptsize

\begin{Verbatim}[numbers=left,numbersep=6pt,xleftmargin=2.2em,fontsize=\scriptsize,breaklines,breakanywhere,frame=lines,framesep=1.5mm,rulecolor=\color{black!35}]
========================================================================
 THE TWENTY-THREE ROWS AT THE APEX · THE EXECUTABLE TWIN · v1.0.1
========================================================================

K1 THE SEAT: conjugation on the integer quaternions fixes exactly the scalar line
  lattice points 28561  fixed by sigma 13

K2 THE RETURN: i.j.k = -1 and the reversed triad k.j.i = +1, both on the seat
  i.j.k =  -1  0  0  0
  k.j.i =   1  0  0  0

K3 THE IDENTITIES AND THE CONE CENSUS
  gap vector (self, register, object) =  0 0 0
  apexes over the ordered diagram 1  over the reversed diagram 1

K4 THE EMBEDDING phi_m(h,t) = <t, h-m, 0, 0> against the fold tau_m(h,t) = (2m-h, t)
  m =  -3: stage points on the line 17, preimage hits 289
  m =  -2: stage points on the line 17, preimage hits 289
  m =  -1: stage points on the line 17, preimage hits 289
  m =   0: stage points on the line 17, preimage hits 289
  m =   1: stage points on the line 17, preimage hits 289
  m =   2: stage points on the line 17, preimage hits 289
  m =   3: stage points on the line 17, preimage hits 289

K5 THE WALL: no readout of the seat equals the deed bit, every readout enumerated
  readouts tried 8192  factorizations of the deed bit 0

K6 THE FREEDOM: wholly odd maps on the two-point set, and the calibration
  wholly odd maps: 2

K7 THE APERTURE: presence is actuation; under the axiom no present reader is interior
  present readers 12 of 12; interior rows 0 of 24

K8 THE RULE: (RA -> L) <-> L is equivalent to RA or L; the two-point frame
  RA = F  L = F  (RA -> L) <-> L = F
  RA = F  L = T  (RA -> L) <-> L = T
  RA = T  L = F  (RA -> L) <-> L = T
  RA = T  L = T  (RA -> L) <-> L = T
  four rows: true in every row where RA holds; where RA fails it is true only where L already holds
  RA true, L(two-point frame) false: the axiom decides the value on no frame

K9 THE LEDGER, THE REFUSAL CONSTANT, THE PRICE
  one bit at 300 K:   2.870979E-21 J

K10 THE CLOSURE LAW: every subset of a sixteen-point carrier under a symmetric registration
  subsets enumerated 65536; closed subsets 64; C1 agrees on every one

K11 CANTOR: no system indexes all of its own binary properties
  k = 3: maps enumerated 512, surjective 0, diagonal outside the range 512
  k = 4: maps enumerated 65536, surjective 0, diagonal outside the range 65536

K12 THE LIVE FACE: self-check is not a witness
  LIVE: [?] interior - self-check is not a witness: the verifier is never the claimant; token withheld
  (supply the argument witnessed to open the live face; the program did not
  and cannot generate its own witness)

K13 THE CURE: on every frame, (E -> L) equals L; grounded counter-models exist; no class cures
  three-point frames 32: L holds on 14, fails on 18, of which grounded (a seat exists) 18
  two-point frame classes 256: cured 32, each exactly a class on which L already holds

K14 THE CONSTANT RECURSION, ANY TRUE PREMISE, THE CARRIER BY DECISION, THE THREE MEASURES

K15 THE TWENTY-THREE ROWS: one seat point, the identity cone on every row, the witness identical
  23 rows: gap vector (0,0,0) each, one apex each, 276 seat-point pairs equal (one seat under every name)

K16 THE WORLD ROW: the typing census, the one crossing, twenty-two bits owed, the rule on every row
  typing census (exact, F4, beachhead, crossed, unbridged, diagonal, one-bit, bridge-gate):   7  3  1  1  3  1  4  3
  crossings 1, bits owed 22, price of the owed bits 63162 yJ at 300 K

========================================================================
 TWIN CENSUS: checks 130718  failures 0
 TWIN-JSON: {"checks":130718,"failures":0,"version":"1.0.0"}
 The twin agrees with the file on every finite claim it carries. Delta-M = 0.
========================================================================

--- ./twin witnessed, the last battery only ---
  LIVE: [I AM] - actuation-occupancy on a witnessed record; conditional at the act
  the witnessed flag was supplied on the command line by the operator running
  this binary; the program did not and cannot generate its own witness.
\end{Verbatim}

\endgroup
\end{document}
